\documentclass[12pt,a4paper]{report}

\usepackage[utf8]{inputenc}
\usepackage[T1]{fontenc}
\usepackage{lmodern}
\usepackage[margin=2.35cm]{geometry}
\usepackage{setspace}
\onehalfspacing
\usepackage{amsmath,amssymb,amsthm,mathtools}
\usepackage{booktabs,tabularx,array,longtable}
\usepackage{tikz}
\usetikzlibrary{arrows.meta,positioning,calc,shapes.geometric,fit,decorations.pathreplacing}
\usepackage{pgfplots}
\pgfplotsset{compat=1.18}
\usepackage{algorithm}
\usepackage{algpseudocode}
\usepackage{caption}
\usepackage{enumitem}
\usepackage{siunitx}
\sisetup{detect-all=true}
\usepackage{fancyhdr}
\usepackage{xcolor}
\usepackage[hidelinks]{hyperref}

\setlength{\parskip}{0.25em}
\setlength{\parindent}{1.4em}
\setlength{\emergencystretch}{2.2em}
\setlist[enumerate]{itemsep=0.25em,topsep=0.4em}
\setlist[itemize]{itemsep=0.25em,topsep=0.4em}
\captionsetup{font=small,labelfont=bf,skip=6pt}
\pgfplotsset{every axis/.append style={label style={font=\small}, tick label style={font=\small}, title style={font=\small}}}

\newtheorem{definition}{Definition}[chapter]
\newtheorem{proposition}{Proposition}[chapter]
\newtheorem{assumption}{Assumption}[chapter]

\newcommand{\qfi}{\mathrm{QFI}}
\newcommand{\mis}{\mathrm{MIS}}
\newcommand{\mul}{\mathrm{MUL}}
\newcommand{\R}{\mathbb{R}}
\newcommand{\E}{\mathbb{E}}
\DeclareMathOperator*{\argmin}{arg\,min}
\DeclareMathOperator{\diag}{diag}

\setlength{\headheight}{15pt}
\pagestyle{fancy}
\fancyhf{}
\fancyhead[L]{\nouppercase{\leftmark}}
\fancyhead[R]{\thepage}
\renewcommand{\headrulewidth}{0.3pt}
\fancypagestyle{plain}{%
  \fancyhf{}
  \fancyfoot[C]{\thepage}
  \renewcommand{\headrulewidth}{0pt}
}

\title{Research for Improving Placement Accuracy and Energy Use of Wooden Seed-Planting Machines by Using Vision-Guided Metering Control and Terrain-Aware Model Predictive Optimization}
\author{Computer Science Thesis\\System Design and Optimization}
\date{}

\begin{document}
\pagenumbering{roman}
\begin{titlepage}
\centering
\vspace*{1.6cm}
{\Large Computer Science Thesis}\\[0.4em]
{\large System Design and Optimization}\\[2.2cm]
{\LARGE\bfseries Research for Improving Placement Accuracy and Energy Use of Wooden Seed-Planting Machines by Using Vision-Guided Metering Control and Terrain-Aware Model Predictive Optimization\par}
\vspace{1.6cm}
\begin{minipage}{0.86\textwidth}
\centering
Two technical innovations for machines that sow woody seed from a laminated-timber chassis:\\[0.8em]
\textit{Compliance-aware vision metering}, and\\
\textit{singulation-constrained terrain model-predictive coverage}.
\end{minipage}
\vfill
{\large A design thesis in three chapters}\\[0.6em]
Previous work \textperiodcentered\ Proposed technologies \textperiodcentered\ System build and estimation
\vspace{1.2cm}
\end{titlepage}

\chapter*{Abstract}
\addcontentsline{toc}{chapter}{Abstract}
Wooden seed, in the sense used here, means seed of woody plants---trees and shrubs sown in agroforestry and restoration---and it also means the material of the machine that carries that seed. A planter built from laminated timber can be made where steel fabrication is scarce, it stores less kinetic energy in a light chassis, and its row unit dissipates vibration more readily than a welded steel parallelogram. The same compliance that damps impact also moves the boot and the seed tube while a seed is falling. Irregular woody seeds already singulate poorly on plates and vacuum orifices that were proportioned for maize. The result is a placement error with two coupled causes: the meter releases the wrong number of seeds, and the wooden arm releases them from a moving gate at the wrong depth.

This thesis proposes two technologies and joins them at a single interface, the admissible ground speed. The first technology, compliance-aware vision metering, watches the seed plate and the upper seed tube with a detector of the YOLO family, fuses that observation with boot force, an ultrasonic soil range and an inertial measurement of the arm, and commands plate speed and boot depth. A structural model of the timber flexure tells the filter when a miss is a true empty cell and when it is a seed whose exit time was shifted by vibration. The second technology, singulation-constrained terrain model-predictive coverage, chooses the heading on a digital elevation model and then tracks the resulting swaths with a linearised model-predictive controller. The stage cost penalises grade-dependent traction work, time, and a miss probability that rises with speed, with acceleration, and with proximity to the resonance of the wooden arm. Speed is therefore not an independent set-point. It is the variable through which coverage respects the latency and the resonance limits of the meter.

The thesis is organised in three chapters. Chapter~1 reviews metering measurement, vision, wooden and manual woody-seed tools, coverage planning and agricultural model-predictive tracking, and closes with a formal problem set. Chapter~2 derives both technologies, including the filter, the feedback law, the energy model and the predictive program. Chapter~3 specifies a four-row build and reports a numerical estimation on a \num{1.44}~hectare field at an eight-percent grade. The estimation is not a field trial. It evaluates the Chapter~2 models with coefficients placed inside ranges published by the works reviewed in Chapter~1. On that model, the joint design at \SI{1.4}{m/s} on a heading \ang{55} from the uphill axis reaches a quality-of-feed index of \num{0.978} and a depth root-mean-square error of \SI{2.15}{mm}, against \num{0.808} and \SI{7.35}{mm} for an open-loop meter run at \SI{1.6}{m/s} along the slope, and it lowers non-conservative mechanical work by about eight percent relative to that uphill heading. A faster open-loop run near the timber resonance, at \SI{2.4}{m/s}, is estimated to destroy establishment even though it finishes sooner.

\chapter*{Acknowledgements}
\addcontentsline{toc}{chapter}{Acknowledgements}
The technical argument rests on published measurement of seed singulation, on published coverage and tracking controllers, and on the wood-engineering data of the Forest Products Laboratory. No claim in Chapter~3 should be read as a substitute for a bench test of a particular timber layup or for a germination trial of a named species.

\tableofcontents
\listoffigures
\listoftables

\chapter*{Nomenclature}
\addcontentsline{toc}{chapter}{Nomenclature}
\setstretch{1.15}
\begin{longtable}{@{}lp{0.72\textwidth}@{}}
\toprule
Symbol & Meaning \\
\midrule
\endhead
$x_{\mathrm{ref}}$ & Target in-row seed spacing \\
$v$ & Ground speed of the row unit \\
$a$ & Longitudinal acceleration \\
$\qfi,\mis,\mul$ & Quality-of-feed, miss and multiple indices \\
$N$ & Number of cells on the metering plate \\
$\omega$ & Plate angular speed \\
$H$ & Vertical drop from release gate to the furrow bottom \\
$t_f$ & Flight time of a seed through the tube and the air gap \\
$k$ & Equivalent stiffness of the timber flexure at the boot \\
$m_{\mathrm{eq}}$ & Equivalent mass of the row unit at the boot \\
$\zeta$ & Viscous damping ratio of the timber flexure \\
$f_n$ & Undamped natural frequency in hertz \\
$\delta$ & Boot displacement from the nominal depth \\
$\tau_v,\tau_a$ & Vision latency and actuator latency \\
$c$ & Fraction of misses that the loop still has time to repair \\
$\mathbf{x},\mathbf{u}$ & Vehicle state and input used by the predictive controller \\
$\beta$ & Heading measured from the uphill direction \\
$\alpha$ & Local slope angle \\
$J$ & Coverage cost: work, time and failed placements \\
CAVM & Compliance-aware vision metering (first technology) \\
SCT-MPC & Singulation-constrained terrain model-predictive coverage (second technology) \\
DEM & Digital elevation model \\
UKF & Unscented Kalman filter \\
\bottomrule
\end{longtable}
\setstretch{1.5}

\pagenumbering{arabic}

\chapter{Research of Previous Work}
\label{ch:previous}

A wooden seed-planting machine sits at an awkward intersection of three literatures that rarely cite one another. Precision-planter research measures singulation of maize, soybean and cotton on steel row units and treats the frame as rigid. Forestry and restoration research measures whether a hand jab, a hollow planting stick or a broadcaster puts a woody or native seed into a narrow depth band, and it treats the mechanism as a hand tool. Agricultural-robotics research plans coverage and tracks a path, and it treats the implement behind the hitch as a load rather than as a process whose quality collapses when the vehicle accelerates. This chapter reads those literatures in the order the two technologies of the thesis require. Section~\ref{sec:prev-meter} reviews metering, sensing and vision, which is the background of compliance-aware vision metering. Section~\ref{sec:prev-plan} reviews coverage and model-predictive tracking, which is the background of singulation-constrained terrain coverage. Section~\ref{sec:problem} turns the gaps into a problem set with indices, constraints and research questions. Throughout, a formula is quoted only when the later design uses it, and a numerical claim from the literature is kept attached to the paper that measured it.

\section{Research on vision-guided seed metering}
\label{sec:prev-meter}

\subsection{What the classical indices measure}

A precision meter is a device that tries to release one seed each time a cell, a finger or a vacuum orifice passes the drop gate. The agronomic object is not the gate. It is the spacing of plants after emergence, which nobody can see at the moment of release. The engineering substitute, standardised for single-seed drills, is the set of spacings measured on a greased belt or in a soil bin, classified against the target spacing \(x_{\mathrm{ref}}\) \cite{iso7256,kachman1995}. With the ordered spacings \(x_i\), the usual three classes are
\begin{equation}
\label{eq:classes}
\begin{aligned}
\text{multiple} &\quad\text{if}\quad x_i < 0.5\, x_{\mathrm{ref}},\\
\text{single} &\quad\text{if}\quad 0.5\, x_{\mathrm{ref}} \le x_i \le 1.5\, x_{\mathrm{ref}},\\
\text{miss} &\quad\text{if}\quad x_i > 1.5\, x_{\mathrm{ref}}.
\end{aligned}
\end{equation}
If \(N_{\mathrm{tot}}\) spacings are observed and \(N_s\), \(N_m\), \(N_u\) fall in the three classes, the quality-of-feed, miss and multiple indices are the frequencies
\begin{equation}
\label{eq:qfi}
\qfi = \frac{N_s}{N_{\mathrm{tot}}},\qquad
\mis = \frac{N_m}{N_{\mathrm{tot}}},\qquad
\mul = \frac{N_u}{N_{\mathrm{tot}}},\qquad
\qfi + \mis + \mul = 1.
\end{equation}
Kachman and Smith showed that a single coefficient of variation, taken over every spacing including the misses, confuses a meter that occasionally skips with a meter that scatters every seed \cite{kachman1995}. The precision index they recommend is the coefficient of variation inside the single class only,
\begin{equation}
\label{eq:cv}
\mathrm{CV}_s = \frac{1}{\bar{x}_s}\sqrt{\frac{1}{N_s-1}\sum_{i\in\mathcal{S}}(x_i-\bar{x}_s)^2},
\end{equation}
where \(\mathcal{S}\) is the set of singles. A meter can therefore be excellent on \(\qfi\) and still place singles with a scatter that a woody seedling, whose emergence window is a few millimetres of depth and a few centimetres of in-row position, will not forgive. Panning, Kocher, Smith and Kachman later took the same indices from the laboratory belt into sugar-beet fields and showed that the ranking of planters was not identical in the two settings, because soil flow and bounce after release are invisible on a greased belt \cite{panning2000}. That distinction matters for a wooden machine: the belt test sees the meter, and the furrow sees the meter plus the motion of the arm.

The kinematic relation that every later controller inverts is the match between plate rate and ground speed. With \(N\) cells on the plate and plate speed \(\omega\) in radians per second,
\begin{equation}
\label{eq:rate}
x_{\mathrm{ref}} = \frac{2\pi v}{N\omega}
\qquad\text{so}\qquad
\omega^{\star}(v) = \frac{2\pi v}{N x_{\mathrm{ref}}}.
\end{equation}
Equation~\eqref{eq:rate} is exact only when every cell releases exactly one seed, the release angle is constant, and the horizontal velocity of the seed relative to the soil is cancelled by the rearward component of the tube or by a brush. None of those conditions holds for a flattened, winged or irregular woody seed. Equation~\eqref{eq:rate} is still the feedforward term. Everything a vision loop adds is a correction around it.

\begin{figure}[ht]
\centering
\begin{tikzpicture}[x=1.15cm,y=1cm,font=\small]
  \draw[thick] (0,0) -- (12.2,0);
  \foreach \x in {0,2,4,6,8,10,12} \draw (\x,0.08) -- (\x,-0.08);
  \node[below] at (0,-0.1) {\(0\)};
  \node[below] at (4,-0.1) {\(x_{\mathrm{ref}}\)};
  \node[below] at (8,-0.1) {\(2x_{\mathrm{ref}}\)};
  \node[below] at (12,-0.1) {\(3x_{\mathrm{ref}}\)};
  \fill (0,0) circle (1.6pt);
  \fill (3.15,0) circle (1.6pt);
  \fill (4.15,0) circle (1.6pt);
  \fill (8.05,0) circle (1.6pt);
  \fill (12.0,0) circle (1.6pt);
  \draw[decorate,decoration={brace,amplitude=5pt,mirror}] (0,-0.85) -- (3.15,-0.85) node[midway,below=8pt] {multiple};
  \draw[decorate,decoration={brace,amplitude=5pt}] (4.15,0.35) -- (8.05,0.35) node[midway,above=8pt] {single};
  \draw[decorate,decoration={brace,amplitude=5pt,mirror}] (8.05,-0.85) -- (12.0,-0.85) node[midway,below=8pt] {miss};
  \draw[dashed] (2,-0.55) -- (2,0.7);
  \draw[dashed] (6,-0.55) -- (6,0.7);
  \node[above] at (2,0.7) {\(0.5x_{\mathrm{ref}}\)};
  \node[above] at (6,0.7) {\(1.5x_{\mathrm{ref}}\)};
\end{tikzpicture}
\caption{Classification of successive seed spacings. The dashed lines are the ISO-style boundaries at half and one-and-a-half times the target spacing. A double close to a nominal site and a skip at the next site are one multiple and one miss, not two noisy singles.}
\label{fig:classes}
\end{figure}

Heege's survey of precision farming treats seeding as one actuation among several that become spatial once the machine knows where it is \cite{heege2013}. Ehsani, Upadhyaya and Mattson had already shown that a real-time kinematic satellite fix can map the intended seed coordinate to centimetre level \cite{ehsani2004}. Mapping the command is not the same as mapping the seed. Griepentrog, N{\o}rremark, Nielsen and Blackmore separated those two maps for sugar beet and argued that a seed map is the right input to later mechanical weeding \cite{griepentrog2005}. The distinction is the one this thesis keeps: a coverage planner produces a command map of where seeds should fall, and a metering controller produces an estimate of where they did fall. Lowenberg-DeBoer and Erickson, reviewing adoption, note that variable-rate and section control spread faster than the agronomy that would justify every rate change \cite{lowenberg2019}. For woody seed the justification is harsher than for maize, because the seed itself is scarce and the replacement cost of a failed hill is a year of growth rather than a later side-dress.

\subsection{Wooden tools and the depth band of woody seed}

The machines that actually put woody and native seed into the ground are, in a large part of the world, still hand tools, and many of those tools are wood or are used as if they were a stick. A Missoula Technology and Development Center evaluation of four commercial single-seed planters---a standard jab planter, a seed stick, a standing tube planter and a transplanter---found that hollow-tube tools with a bit that opens at the bottom planted small tree seed at a set depth more reliably than tools that merely poked and hoped, and that durability of the bit and of the release latch was the practical failure \cite{usdaplanter}. The operator supplied the singulation by dropping one seed into the tube. That division of labour is honest about mechanism: wood and a simple latch can make a hole, and they cannot see a double.

A jab planter built and tested by final-year bioresources students for no-till conditions assembled the same functions---frame and handle, hopper, precision meter, delivery tube, ground opener---into one manual machine, and compared it with a hoe \cite{jab2010}. The meter is the part that decides the index in Equation~\eqref{eq:qfi}. The frame is the part that decides whether the tube is vertical when the seed leaves it. A mechatronic dibbler built for a radiata-pine nursery attacked a related defect at production scale. Poorly shaped holes were bending stems; about \num{22}\% of seedling rejection at that nursery was attributed to the hole rather than to the plant. A machine that cut vertical holes of a specified depth removed the stem bend and raised worker throughput by about \num{10}\% \cite{dibbler}. The lesson for a field planter is not the nursery mechanism. It is the separation between a geometric specification (vertical, depth within a band) and a human or a servo that tries to meet it.

Masarei, Guzzomi, Merritt and Erickson measured the emergence of soft spinifex and found the optimum at a covering of \SI{5}{mm}, a sharp fall by \SI{10}{mm}, and almost no emergence from \SI{30}{mm}, on the rocky profiles of mine waste \cite{masarei2019}. They also record the practitioner survey result that depth precision is one of the properties restoration seeding lacks, and that agricultural openers built for crop seed do not hold a band of a few millimetres once the surface is rough. Bond, Honig and Maze had given the allometric reason: larger seeds reach the surface from deeper burial, so a depth setting copied from maize is not a conservative default for a small woody or native seed \cite{bond1999}. St Jack, Hesterman and Guzzomi, metering sandalwood seed, showed that a species with an irregular and relatively large woody seed still needs a meter designed around that shape rather than a crop plate with the vacuum turned up \cite{stjack2013}. This thesis takes a mid-size woody seed as its design case, with a target depth of \SI{18}{mm} and an estimation band of \(\pm\SI{8}{mm}\). Those two numbers are design assumptions, chosen to sit between the maize band and the spinifex band. They are not a measured optimum for a named species.

Wood enters the problem as structure, not as a romantic material. The Wood Handbook collects the design values: modulus of elasticity parallel to grain for common structural softwoods is on the order of \SIrange{8}{13}{GPa}, density is about a third to a half of aluminium, and internal friction in cellular tissue produces damping ratios that steel does not match without an added damper \cite{woodhandbook}. A row-unit link built as a timber cantilever of length \(L\), second moment \(I\) and modulus \(E\) has the tip stiffness
\begin{equation}
\label{eq:cantilever}
k = \frac{3EI}{L^3},\qquad I = \frac{bh^3}{12}
\end{equation}
for a rectangular cross-section bent about the strong or the weak axis according to how \(h\) is chosen. The equivalent mass \(m_{\mathrm{eq}}\) at the boot then sets the undamped frequency
\begin{equation}
\label{eq:fn}
f_n = \frac{1}{2\pi}\sqrt{\frac{k}{m_{\mathrm{eq}}}}.
\end{equation}
A steel parallelogram is usually designed so that \(f_n\) sits far above the disturbance spectrum of the ground. A timber link sized for a person to lift, or for a small nursery workshop to laminate, does not have that luxury. Roughness of wavelength \(\lambda\) met at speed \(v\) arrives at frequency \(v/\lambda\). When that frequency approaches \(f_n\), the dynamic magnification of a viscously damped oscillator,
\begin{equation}
\label{eq:mag}
M(r) = \bigl[(1-r^2)^2 + (2\zeta r)^2\bigr]^{-1/2},\qquad r = \frac{v}{\lambda f_n},
\end{equation}
grows without any change in the meter. Previous planter papers do not put Equation~\eqref{eq:mag} next to Equation~\eqref{eq:qfi}. That omission is the first gap this thesis works on. Soil-tool force models of the kind reviewed by Godwin and O'Dogherty supply the load that drives the flexure: draft and vertical force depend on depth, speed and a soil strength parameter, and they are large enough, on a light wooden boot, to move the set depth by millimetres \cite{godwin2007}. A cone index measured with the standard penetrometer is the field number those models are usually fed \cite{asabe313}.

\begin{figure}[ht]
\centering
\begin{tikzpicture}[font=\small, >=Stealth, scale=0.96]
  \draw[thick, fill=brown!25] (0,0) rectangle (2.3,0.28);
  \node at (1.15,0.14) {timber chassis};
  \draw[thick] (0.4,0.28) -- (0.4,1.7);
  \draw[thick] (1.9,0.28) -- (1.9,1.7);
  \draw[thick, fill=brown!40] (0.15,1.7) rectangle (2.15,2.55);
  \node[align=center] at (1.15,2.15) {hopper\\and plate};
  \draw[thick] (1.15,1.7) -- (1.15,0.55);
  \draw[thick, fill=black!10] (0.95,0.55) rectangle (1.35,0.15);
  \node[right] at (1.4,0.35) {boot};
  \draw[dashed] (2.5,2.3) -- (4.1,2.3);
  \draw[->] (4.1,2.3) -- (4.1,0.35);
  \node[right, align=left] at (4.2,1.3) {drop \(H\)\\flight \(t_f\)};
  \draw[<->] (2.35,0.28) -- (2.35,1.7) node[midway,right] {\(L\)};
  \draw[->] (1.15,0.15) -- (1.15,-0.45) node[below] {\(F_{\mathrm{soil}}\)};
  \node[align=left, anchor=west] at (4.2,2.6) {camera looks\\at the plate};
  \draw[->] (4.2,2.45) -- (2.2,2.2);
\end{tikzpicture}
\caption{Schematic row unit. The hopper and the plate sit on a timber link of length \(L\). Soil force at the boot deflects the link. The seed falls a height \(H\) after release, so any horizontal velocity of the moving gate becomes a spacing error.}
\label{fig:rowunit}
\end{figure}

\subsection{Optoelectronics and high-speed imaging}

Before a camera was cheap enough to mount on a row unit, singulation was measured with a light plane across the seed tube. Kocher, Lan, Chen and Smith built an optoelectronic sensor that timed seeds as they crossed a beam and converted inter-arrival time into spacing at a known ground speed \cite{kocher1998}. The method is fast, it needs no training set, and it cannot tell a clod of soil or a chip of husk from a seed, nor can it see two seeds that cross the plane overlapping. It also sees the seed only after the release decision is over, so it is a monitor unless a later cell can still be advanced or retarded.

High-speed imaging removed that ambiguity on the test stand. Mangus and colleagues mounted a camera at \num{300} frames per second on a variable-rate electric-drive planter and combined the frames with the machine's own speed and seed-tube pulses \cite{mangus2017}. Across constant speeds from \SI{2.4}{km/h} to \SI{16.1}{km/h} the mean singulation was \num{98.45}\%. Doubles were rare (\num{0.03}\%) and triples were absent in that data. The errors that did appear were organised by the operating state: they nearly doubled when the machine accelerated or decelerated at \SI{4.8}{km/h} per second rather than at \SI{2.4}{km/h} per second, and they rose once the meter was asked for more than about \num{1250} seeds per minute. Vision counts and the commercial seed-tube sensor agreed within \(0.8\pm 0.2\) percentage points. The paper's design recommendation is the one a coverage controller should read as a constraint. There exists a seed-per-minute region in which the meter is good, and headland behaviour knocks the machine out of that region even when the steady-state plate calibration is perfect.

Zhao, Gao, Su, Xu and Huang combined a discrete-element simulation of a high-speed meter with a high-speed camera, which is the right pairing when the object is the trajectory inside the housing rather than only the spacing on the belt \cite{zhao2024hs}. Discrete-element work itself goes back to Cundall and Strack's contact law for discs and spheres \cite{cundall1979}. Lei and colleagues used it, with a bench test, on a centralised rapeseed meter, where the seeds interact in a stream rather than in a single cell \cite{lei2021}. For a woody seed the contact parameters---friction with a timber or polymer wall, restitution after a bounce off a wooden tube---are not the maize parameters those studies identify. The method transfers. The numbers do not.

Navid, Ebrahimian, Gassemzadeh and Mousavinia showed that a laboratory camera and a classical image-processing chain, without a learned detector, already recover spacing statistics from a side view of falling seeds \cite{navid2011}. Classical processing is attractive on a wooden machine because the computing budget may be a single board. It fails when the seed is the same colour as the wooden tube, when dust coats the window, or when two overlapping silhouettes must be counted as two objects. Those are ordinary conditions in a restoration planting, and they are why the later papers move to learned detectors.

\subsection{Detectors that name an empty cell and a double}

Borja, Amongo, Suministrado and Pabico closed the loop, rather than only measuring it \cite{borja2018,borja2021}. A camera looked at the seed plate of a pneumatic maize meter. A processor labelled each hole filled or empty. A stepper advanced the plate to replace an empty hole before the gate. On \num{780} hole readings the vision error was \num{0.2564}\%. At simulated \SI{2}{km/h} and \SI{4}{km/h} the missed-hill count fell, relative to the same meter with the vision path disabled, by roughly half at the lower speed and by about a fifth at \SI{4}{km/h}. At \SI{6}{km/h} the vision-assisted meter missed \emph{more} hills than the open-loop meter. The authors trace the reversal to the microcontroller: image latency and a drive that could not hold a steady delivery rate once the hole period became short. The identification was accurate. The actuation was late. A late correction is not a neutral failure. It inserts a seed into a neighbour's interval and manufactures both a multiple and a miss.

That result is the empirical centre of the first technology. Any proposal that adds vision to a planter has to state the inequality between the hole period and the sum of exposure, inference and mechanism delay. With the feedforward of Equation~\eqref{eq:rate}, the hole period is
\begin{equation}
\label{eq:period}
\Delta t = \frac{x_{\mathrm{ref}}}{v} = \frac{2\pi}{N\omega}.
\end{equation}
A correction computed from an image is useful when it is applied to a cell that has not yet reached the gate, which means the processing delay \(\tau_v\) and the actuator delay \(\tau_a\) must satisfy
\begin{equation}
\label{eq:latency}
\tau_v + \tau_a < \eta\, \Delta t
\end{equation}
for a margin \(\eta < 1\) that covers jitter. Borja's rig crossed Equation~\eqref{eq:latency} between \SI{4}{km/h} and \SI{6}{km/h}. Mangus's electric-drive meter, watched rather than corrected, still singulated well at far higher speed because the meter itself was mechanically capable and the camera was not in the control path \cite{mangus2017}. The two papers are consistent. Measurement at high speed is solved more easily than correction at high speed.

Learned detectors changed what a frame can be asked. Blaser, Porter, Fuhrer and Bourlai trained YOLOv10 on peanut vacuum-meter videos from a test stand and from a planter, with augmentation for light and camera angle, and report a mean average precision at IoU \(0.5\) of \num{99.4}\% and an error of \num{1.91}\% on the three-way count of singles, empties and doubles \cite{blaser2024,wang2024yolov10}. Zhao, Wu, Cheng, Gao, Liu, Lai and Li paired an improved YOLOv5s with a SiamRPN++ tracker on a high-speed maize meter. Precision and recall of detection were \num{98.5}\% and \num{99.2}\%. Tracking was weaker (success \num{63.3}\%, precision \num{71.2}\%). The quality-of-feed error against a high-speed camera stayed under \num{3.3}\% from \SI{10}{km/h} to \SI{18}{km/h} \cite{zhao2025yolo}. Zhao, Xi, Shi, Zhang, Qu and Zhang had already put an improved YOLOv5s on the on-line detection of seeding performance, which is the deployment this thesis needs, as distinct from a bench score \cite{zhaoJ2023}. Bustos-Gayt{\'a}n and colleagues used machine vision to separate the influence of meter components and operating parameters on uniformity, which is the diagnostic use of the same pictures: the image is not only a count, it is evidence about which knob to turn \cite{bustos2025}.

\begin{figure}[ht]
\centering
\begin{tikzpicture}[x=0.95cm,font=\small, >=Stealth]
  \draw[thick] (0,0) -- (12.5,0);
  \foreach \x/\lab in {0/expose, 1.6/infer, 3.4/decide, 5.2/move plate, 8.2/gate}{
    \draw (\x,0.12) -- (\x,-0.12);
  }
  \draw[line width=2pt, blue!70!black] (0,0.45) -- (1.5,0.45) node[midway, above] {\(\tau_{\mathrm{exp}}\)};
  \draw[line width=2pt, blue!70!black] (1.6,0.45) -- (3.3,0.45) node[midway, above] {\(\tau_{\mathrm{nn}}\)};
  \draw[line width=2pt, orange!80!black] (3.4,0.45) -- (5.1,0.45) node[midway, above] {logic};
  \draw[line width=2pt, red!70!black] (5.2,0.45) -- (8.0,0.45) node[midway, above] {\(\tau_a\)};
  \draw[dashed] (8.2,-0.7) -- (8.2,1.15);
  \node[below] at (8.2,-0.7) {cell at gate};
  \draw[<->] (0,-1.15) -- (8.2,-1.15) node[midway, below] {must finish inside \(\eta\Delta t\)};
  \node[anchor=west] at (8.6,0.7) {if \(\Delta t\) shrinks,};
  \node[anchor=west] at (8.6,0.25) {the red bar};
  \node[anchor=west] at (8.6,-0.2) {crosses the gate};
\end{tikzpicture}
\caption{Latency budget of a corrective vision loop. Detection can be accurate and still harmful if the actuator command arrives after the empty cell has passed the release gate.}
\label{fig:latency}
\end{figure}

The detector papers and the Borja papers do not yet form one design. The detector papers evaluate a recording. Borja evaluates a correction and finds a speed at which the correction changes sign. Neither paper writes the wooden arm into the observation model. A seed that left the cell on time and then struck a vibrating wooden tube looks, in a tube camera, like a late seed. A filter that treats every late crossing as an empty cell will advance the plate and create the double Borja already saw from latency alone.

\subsection{Vacuum, contact mechanics and design inversion}

The pneumatic meter remains the reference mechanism because its pickup force can be written down. Karayel, Barut and {\"O}zmerzi modelled the vacuum pressure required to hold a seed against gravity and inertia as a function of orifice area, seed mass and the speed of the plate \cite{karayel2004}. In schematic form the pressure difference that holds a seed of mass \(m_s\) and projected area \(A_o\) against a centrifugal acceleration \(r\omega^2\) is of the order
\begin{equation}
\label{eq:vacuum}
\Delta p \ge \frac{m_s}{A_o}\bigl(g + r\omega^2 + a_{\mathrm{vib}}\bigr),
\end{equation}
where \(a_{\mathrm{vib}}\) is the extra acceleration of a compliant mount. Published forms of the model differ in the coefficients on friction and on the angle of the orifice. The term this thesis needs is \(a_{\mathrm{vib}}\). On a rigid steel meter it is a small residual. On a timber flexure near resonance it is the same order as \(g\), and a vacuum that was sized on a rigid rig will drop seeds that the sizing calculation declared safe. Singh, Singh and Saraswat optimised orifice, vacuum and speed for cotton and showed the familiar ridge in the response surface: too little vacuum misses, too much vacuum doubles, and speed shrinks the ridge \cite{singh2005}. An irregular woody seed widens the uncertainty on \(m_s/A_o\) because no two seeds present the same face to the orifice. That is an argument for a vision check even when the mean vacuum is correct, and it is an argument against copying a maize orifice into a wooden hopper and expecting the cotton response surface to reappear.

Maleki and colleagues measured a multi-flight auger as a grain-drill meter and found that volumetric devices, which are forgiving of seed shape, do not produce the single-seed spacing of Equation~\eqref{eq:classes} \cite{maleki2006}. For a restoration mix of several woody species a volumetric flute is sometimes the honest mechanism. For a single high-value woody species, where the thesis's indices apply, a cell plate or a vacuum plate remains the mechanism that vision can correct cell by cell. Anantachar, Kumar and Guruswamy trained a neural network on an inclined-plate meter and inverted it to propose design and operating parameters \cite{anantachar2010}. Inversion of a surrogate is useful when the plate geometry is the decision. It does not see a new seed lot, a film of pitch on a woody seed, or a flexure that was not in the training rig. On-line vision is the measurement the surrogate cannot replace. The two are complementary: the network proposes a plate and a nominal vacuum, and the loop of Section~\ref{sec:tech1} trims them while the machine moves.

Flight after release converts a small timing error into a spacing error. Under constant vertical acceleration from rest relative to the tube,
\begin{equation}
\label{eq:flight}
H = v_{z0} t_f + \tfrac12 g t_f^2.
\end{equation}
If the seed is simply dropped, \(v_{z0}\approx 0\) and \(t_f=\sqrt{2H/g}\). A horizontal velocity \(v\) of the machine, plus a horizontal velocity \(\dot{y}_{\mathrm{arm}}\) of the gate caused by arm flexure, carries the seed
\begin{equation}
\label{eq:carry}
\Delta x = \bigl(v + \dot{y}_{\mathrm{arm}}\bigr)\, t_f
\end{equation}
before it reaches the furrow. At \(H=\SI{0.42}{m}\), \(t_f\approx\SI{0.29}{s}\). At \(v=\SI{1.6}{m/s}\) the machine travels about \SI{0.47}{m} during the flight, which is more than twice the \SI{0.20}{m} design spacing of the woody-seed case. The tube and the furrow closer must cancel the mean of that transport. They cannot cancel a zero-mean vibration that changes from seed to seed. Equation~\eqref{eq:carry} is why arm velocity belongs in the metering filter and not only in a structural side calculation.

\subsection{What this line of work does not close}

The metering literature now contains accurate indices, a vacuum model, a discrete-element method, optoelectronic monitors, high-speed confirmation that acceleration and seed rate organise the errors, and detectors that name empties and doubles at errors of a few percent. It also contains one clear negative result: a correct detector plus a slow processor makes singulation worse above a modest speed \cite{borja2021}. It does not contain a filter whose process model is the flexure of a timber row unit, and it does not contain a speed command that refuses the resonance of that flexure and the latency limit of Equation~\eqref{eq:latency}. Those two missing pieces are one technology, because the same state vector carries the cell label and the arm motion. Section~\ref{sec:tech1} builds that technology. It cannot, by itself, decide which way the machine should drive across a slope, or how hard it may accelerate into a headland. That decision is the subject of the next section.

\section{Research on coverage planning and model-predictive tracking}
\label{sec:prev-plan}

\subsection{Coverage as a geometric problem}

Coverage path planning asks for a path whose swept width paints a region and whose cost, usually length or time, is small. Choset's survey fixed the vocabulary: cellular decomposition, Morse functions, and the boustrophedon pattern in which the machine runs a straight swath, turns on a headland, and runs the next swath in the opposite direction \cite{choset2001}. Galceran and Carreras updated the survey across robotics and separated off-line planners, which assume a known map, from on-line planners, which discover the boundary as they travel \cite{galceran2013}. A seed-planting machine of the kind considered here is an off-line planner at the scale of the field and an on-line controller at the scale of the swath. The boundary is known from a survey or from a previous pass. The disturbance that matters is not an unknown wall. It is a grade, a patch of soft soil, and a meter that misbehaves when the speed profile is aggressive.

Oksanen and Visala published coverage algorithms aimed at agricultural field machines rather than at indoor robots, including the treatment of headlands, obstacles and the choice of driving direction on real field polygons \cite{oksanen2009}. LaValle's planning text supplies the underlying algorithms---grid search, sampling, feedback on a cell complex---that those field methods specialise \cite{lavalle2006}. Dijkstra's shortest paths and the A* search of Hart, Nilsson and Raphael remain the workhorses inside a grid once a field has been rasterised \cite{dijkstra1959,hart1968}. They are optimal on the grid they are given. They are silent on whether the cost of an edge should be metres, joules, or expected missed seeds. Putting the wrong scalar on the edge is not a small modelling choice. It is the difference between a short path and a path the meter can survive.

Bochtis and Vougioukas minimised the non-working distance of a machine that covers a field in a headland pattern, which is the distance spent turning and relocating rather than sowing \cite{bochtis2008}. Bochtis, S{\o}rensen and Busato reviewed the wider machinery-management problem: sequencing of fields, capacity of the machine, and the distinction between a path that is short and a plan that a fleet can execute \cite{bochtis2014}. S{\o}rensen and Bochtis had set out the fleet-management concepts that sit above a single planter \cite{sorensen2010}. Spekken and de Bruin optimised in-field routing by counting manoeuvring and servicing time, which is the right objective when the expensive event is the turn rather than the metre of straight work \cite{spekken2013}. Edwards and colleagues evaluated a prototype in-field route planner on neutral-material-flow operations and showed that a planner's predicted gain survives contact with a real sequence of loads only if the model of service time is honest \cite{edwards2017}. Fountas and colleagues place all of these plans in a farm management information system, which is where a seed map and a coverage file have to live if a later operation is to find them \cite{fountas2015}. This thesis stops at one machine and one field. The interface it leaves upward is a coverage file and a per-row seed estimate, which is what those information systems already know how to store.

\subsection{Driving direction, headlands and three-dimensional energy}

On flat ground the best driving direction is mostly a statement about how many turns fit along the short side of the field. Hameed, Bochtis, S{\o}rensen and N{\o}rremark generated guidance lines for operational planning and made the driving angle a decision variable rather than a habit of the operator \cite{hameed2010}. Jin and Tang solved coverage on two-dimensional field surfaces with an explicit cost on turns and on overlapped or missed area \cite{jin2010}, and then lifted the same problem onto three-dimensional terrain, where the slope changes the cost of each heading and can change the decomposition of the field \cite{jin2011}. The three-dimensional statement is the one a wooden planter on a restoration slope needs. A heading that is only a few degrees off the contour carries a sustained component of weight down the drawbar.

Vahdanjoo, Gislum and S{\o}rensen built a three-dimensional coverage tool for a seeding operation, using a digital elevation model to choose the driving direction against an energy criterion, and they also produced a variable-rate map and a compaction-risk estimate \cite{vahdanjoo2023}. On two sample fields their model reduced the required energy by up to \num{9.5}\% and \num{15.6}\% relative to a poorer direction. The mechanism is not mysterious. Rolling resistance and wheel slip depend on the normal and the tangential components of weight,
\begin{equation}
\label{eq:grade}
F_{\parallel} = mg\sin\alpha,\qquad F_{\perp} = mg\cos\alpha,
\end{equation}
and a boustrophedon that runs up and down cancels the conservative part of \(F_{\parallel}\) only in the energy ledger of a lossless vehicle. Slip, rolling of a soft tyre in soft soil, and the extra draft of an opener that is loaded unevenly are not conservative. They leave a remainder that depends on how the heading sits relative to the gradient. Their paper is the closest neighbour of the second technology. It optimises energy of seeding coverage on a real elevation model. It does not constrain the speed so that a vision meter can finish a correction, and it does not know that a timber arm has a resonant speed.

Headland geometry interacts with that energy term because the turn is where speed is low and acceleration is high, which is exactly the state Mangus found doubles the singulation error \cite{mangus2017}. Graf Plessen and Bemporad treat the shaping of the reference that joins one swath to the next. High-level planners emit corners. A machine cannot track a corner. They compare circle arcs, generalised elementary paths and bi-elementary paths, and they hand the smoothed reference to a linear time-varying model-predictive tracker, with constraints on actuator rate and on the field boundary so that a sprayer does not leave gaps or paint the crop twice in the turn \cite{plessen2017}. The same smoothing is required for a planter, with one added constraint they did not have: during the acceleration out of the turn the predicted miss index must stay under a cap. A geometrically beautiful clothoid that demands \SI{1.2}{m/s^2} for two seconds is a bad planter trajectory even when it is a good sprayer trajectory.

Zhang, Qian, Li, Bai and Wang assembled a seeding-robot planner from a lidar-camera perception front end, an adaptive raster, an A* search fused with ant-colony optimisation, and a multi-objective scoring of the path \cite{zhang2025path,dorigo2006}. With a raster of \SI{0.15}{m}, a pheromone factor of \num{2.5}, a heuristic factor of \num{3.0} and \num{100} iterations they report a planning time of \SI{0.45}{s}, a success rate of \num{98.7}\%, an overlap of \num{3.2}\%, an efficiency gain of \num{28.5}\%, an energy reduction of \num{21.3}\% and a missed-sowing rate of \num{1.8}\% in their trials. The architecture is the right shape for a seeding robot: perception, map, search, score. The score, as published, does not include a latency inequality from a vision meter or a resonant frequency of a wooden link. Ant-colony search can carry those terms as soon as they are written into the edge cost. The contribution is the term, not another metaheuristic.

\subsection{Tracking controllers and what they optimise}

Once a swath exists, a tracker has to make the hitch follow it. The kinematic bicycle model is the standard compromise between a full tyre model and a pure geometric aim point. With wheelbase \(\ell\), speed \(v\) and steering angle \(\delta\),
\begin{equation}
\label{eq:bicycle}
\dot{x} = v\cos\theta,\qquad
\dot{y} = v\sin\theta,\qquad
\dot{\theta} = \frac{v}{\ell}\tan\delta .
\end{equation}
Rajamani's vehicle-dynamics text gives the passage from Equation~\eqref{eq:bicycle} to models with sideslip when the speed and the surface justify the extra state \cite{rajamani2012}. At the speeds of a wooden planter, typically under \SI{3}{m/s}, the kinematic model plus a bounded slip angle is the appropriate complexity, and it is the model used by the agricultural predictive controllers cited below.

Backman, Oksanen and Visala derived a nonlinear model-predictive path tracker for agricultural machines and showed that the predictive formulation can respect steering rate and can look through a headland rather than reacting to it \cite{backman2012}. Kayacan, Kayacan, Ramon and Saeys identified the yaw dynamics of an autonomous tractor and compared trajectory trackers, establishing that a model identified on the real steering loop beats a textbook gain when the soil changes \cite{kayacan2015}. Kraus, Ferreau, Kayacan, Ramon, De Baerdemaeker, Diehl and Saeys combined a moving-horizon estimator with nonlinear model-predictive control so that the state the controller sees is estimated from the same noisy sensors the tracker must live with \cite{kraus2013}. The estimator-controller pair is the pattern of the first technology as well: an unscented Kalman filter on the arm and the seed event, and a short-horizon correction of plate speed and depth. The computational pattern is shared. The state vector is not.

Model-predictive control, in the form used here, is a finite-horizon optimal-control problem resolved at each sample. Mayne, Rawlings, Rao and Scokaert proved the stability and optimality statements that justify a terminal cost and a terminal constraint as the ingredients that make a receding-horizon law inherit the stability of the infinite-horizon problem \cite{mayne2000}. Rawlings, Mayne and Diehl give the engineering form: a quadratic stage cost, linearised dynamics, and a quadratic program \cite{rawlings2017}. Boyd and Vandenberghe supply the convex background that makes that quadratic program tractable when the constraints stay linear \cite{boyd2004}. Ferreau, Kirches, Potschka, Bock and Diehl implemented the parametric active-set solver qpOASES, which is small enough to run on an embedded computer at the rate a steering loop needs \cite{ferreau2014}. Andersson, Gillis, Horn, Rawlings and Diehl implemented CasADi, which is the tool this thesis uses to form the linearisation and the quadratic program in the estimation of Chapter~\ref{ch:estimate} \cite{andersson2019}. None of that machinery knows what a seed is. It will optimise whatever stage cost it is given, including a stage cost that is indifferent to singulation.

Xu, Lai, Guo and Lu name that indifference. Their efficiency-oriented model-predictive controller keeps a nested structure so that the tracking layer can still improve a global efficiency objective, rather than tracking a frozen path so faithfully that efficiency gets worse. In simulation they report an \num{8.56}\% efficiency gain against a conventional nonlinear predictive tracker \cite{xu2022}. The device they introduce, a pseudo-point that the machine is not required to touch, is a way of putting a global objective inside a short horizon. This thesis uses a different device for a different objective. The global objective that must enter the short horizon is not only time. It is the probability that the next few seeds are misses, which is a function of speed, acceleration and the vibration state of the arm. A tracker that is willing to arrive earlier by crossing the resonant speed of the timber is, on the indices of Equation~\eqref{eq:qfi}, a worse planter.

Bakker and colleagues demonstrated autonomous navigation of a robot platform in a sugar-beet field, which remains a useful existence proof that a small vehicle can stay in a row well enough for an implement to work \cite{bakker2011}. N{\o}rremark, Griepentrog, Nielsen and S{\o}gaard measured the accuracy of an autonomous satellite-guided intra-row weeder and showed how positioning error and implement delay add \cite{norremark2008}. Bechar and Vigneault review agricultural robots by subsystem---detection, action, navigation, decision---which is the decomposition Section~\ref{sec:build} uses for the build \cite{bechar2016}. Shamshiri and colleagues set the same components in a digital-farming frame \cite{shamshiri2018}. Auat Cheein and Carelli discuss unmanned service units as the architectural unit, a vehicle plus an implement plus a task, which matches a planter better than a generic mobile robot does \cite{auat2013}. Blackmore, Stout, Wang and Runov argued early that robotic agriculture would be organised around plant-scale sensing and action rather than around ever-wider implements \cite{blackmore2005}. A four-row wooden planter is a modest instance of that argument: the sensing is per row, because a single plate camera cannot see four plates, and the action is per row, because four timber flexures do not share one resonance phase.

Pedersen, Fountas, Have and Blackmore analysed the economics of agricultural robots and concluded that a robot pays when it replaces a scarce labour peak or when it enables an operation whose timing a large tractor misses \cite{pedersen2006}. Direct seeding of woody plants has both features. The labour peak of hand jab planting is real, as the nursery dibbler study quantified from the other side \cite{dibbler}, and the depth band is a timing-and-geometry requirement that a broadcaster does not meet \cite{masarei2019}. The economic case is not re-derived in this thesis. It is the reason a modest energy saving is not the only term in the cost. A failed seed is allowed to dominate a joule.

ISO~11783, the serial network usually called ISOBUS, is the practical way a plate-speed command and a ground-speed measurement already travel on a modern tractor \cite{iso11783}. The build in Section~\ref{sec:build} does not depend on a particular tractor brand. It does depend on a clock that can stamp a plate command and a camera frame onto the same time base. A proprietary analogue drive without a time stamp cannot host either technology.

\subsection{The missing coupling}

\begin{figure}[ht]
\centering
\begin{tikzpicture}[font=\small, node distance=0.55cm, >=Stealth]
  \node[draw, rounded corners, align=center, fill=blue!8, minimum width=4.4cm, minimum height=1.1cm] (meter) {Metering literature\\indices, vision, vacuum};
  \node[draw, rounded corners, align=center, fill=green!10, minimum width=4.4cm, minimum height=1.1cm, right=1.5cm of meter] (plan) {Coverage literature\\heading, energy, MPC};
  \node[draw, rounded corners, align=center, fill=brown!15, minimum width=4.4cm, minimum height=1.1cm, below=0.9cm of meter] (wood) {Timber flexure\\ \(f_n\), \(\zeta\), depth error};
  \node[draw, rounded corners, align=center, fill=orange!15, minimum width=4.4cm, minimum height=1.1cm, below=0.9cm of plan] (speed) {Shared interface\\admissible speed \(v\)};
  \draw[->] (meter) -- (speed);
  \draw[->] (plan) -- (speed);
  \draw[->] (wood) -- (speed);
  \draw[->] (wood) -- (meter);
\end{tikzpicture}
\caption{The two literatures meet the timber flexure at one variable. Vision metering can state a speed it is able to correct. Coverage can state a speed and a heading that save energy. Neither statement is safe unless both are true at once.}
\label{fig:gap}
\end{figure}

Table~\ref{tab:gaps} collects the gap in the form the design will use. Each row is a result the thesis inherits and a decision the inherited result does not make.

\begin{table}[ht]
\centering
\caption{Results inherited from previous work, and the decision each result leaves open.}
\label{tab:gaps}
\small
\begin{tabularx}{\textwidth}{@{}lX X@{}}
\toprule
Source & Established result & Decision left open \\
\midrule
Kachman and Smith; ISO 7256 & Separate miss, multiple and single-class scatter & How a live controller should weight the three \\
Borja et al. & Vision cuts misses at \SIrange{2}{4}{km/h} and adds them at \SI{6}{km/h} & A delay model that prevents the sign change \\
Mangus et al. & Acceleration and seed rate organise the residual error & A tracker forbidden from those accelerations \\
Blaser; Zhao et al. & Detectors name empty and double at about \num{2}\% error & Using the label before the gate, on a vibrating tube \\
Karayel et al. & Vacuum scales with mass, area and plate acceleration & The extra acceleration of a timber mount \\
Masarei et al.; USDA & Depth band and hand-tool practice for small seed & A servo depth loop on a compliant boot \\
Wood Handbook & \(E\), density and damping of timber & A resonance the vehicle is not allowed to sit on \\
Vahdanjoo et al. & Heading can save on the order of \numrange{10}{16}\% of energy & Heading chosen also for singulation and turns \\
Backman; Kraus; Xu & Predictive trackers follow a swath and can see efficiency & A stage cost that contains \(\mis(v,a,f_n)\) \\
\bottomrule
\end{tabularx}
\end{table}

The coupling can be stated without yet choosing an algorithm. Let \(\mathcal{V}_{\mathrm{meter}}(x_{\mathrm{ref}},\tau_v,\tau_a,f_n)\) be the set of speeds at which the vision loop still repairs misses and the flexure magnification \(M(r)\) stays below a chosen gain. Let \(\mathcal{V}_{\mathrm{plan}}(\beta)\) be the set of speeds a coverage plan would like to use on heading \(\beta\) in order to finish and to keep grade losses small. A plan is admissible only if it can be executed inside the intersection
\begin{equation}
\label{eq:intersect}
v(t) \in \mathcal{V}_{\mathrm{meter}} \cap \mathcal{V}_{\mathrm{plan}}\bigl(\beta\bigr)
\quad\text{for all } t \text{ in a working swath.}
\end{equation}
Previous metering papers compute something like \(\mathcal{V}_{\mathrm{meter}}\) after an experiment, as a conclusion. Previous coverage papers compute something like \(\mathcal{V}_{\mathrm{plan}}\) as an input the tracker must obey. Equation~\eqref{eq:intersect} makes the intersection the design object. Section~\ref{sec:problem} writes the problems whose solutions have to land in that intersection. Chapter~\ref{ch:design} constructs the two technologies that implement it.

\section{Problem set}
\label{sec:problem}

\subsection{Operating context and scope}

The machine considered from here on is a four-row planter. Each row unit carries a cell plate, a wooden or wood-lined seed tube, a boot, a camera aimed at the plate and the upper tube, a boot-force sensor, an ultrasonic range to the soil surface and an inertial sensor on the arm. The chassis and the parallel links are laminated softwood. The design case uses \(N=24\) cells, \(x_{\mathrm{ref}}=\SI{0.20}{m}\), a drop \(H=\SI{0.42}{m}\) and a target depth of \SI{18}{mm}. The field used later for estimation is a rectangle \SI{80}{m} by \SI{180}{m} with the long axis up the slope, at a grade of \num{8}\%, so that a contour heading produces short rows and many turns and an uphill heading produces long rows and a sustained grade load. That geometry is chosen because it does not let one heading win on every term at once.

The following assumptions bound the problem. They are restated where a derivation would otherwise smuggle them in.

\begin{assumption}
\label{as:seed}
Seeds are rigid bodies of one species for the duration of a pass. Mass and projected area vary from seed to seed but stay inside a known bounded set. The meter is a cell plate, not a volumetric flute.
\end{assumption}

\begin{assumption}
\label{as:flex}
Each row-unit flexure is represented by one equivalent mass, stiffness and damping ratio at the boot. Coupling between the four row units through the chassis is neglected in the filter and is revisited as a modelling error in the estimation.
\end{assumption}

\begin{assumption}
\label{as:map}
The field boundary and a digital elevation model are known before work starts. Obstacles are static and are already cut out of the coverage polygon. Soil strength enters only through the force that deflects the boot and through the rolling-resistance coefficient.
\end{assumption}

\begin{assumption}
\label{as:est}
Chapter~\ref{ch:estimate} evaluates models. It does not report a germination trial, a timber fatigue test, or a detector trained on a new image set. Detector precision and recall are taken from the published peanut and maize figures cited above and are treated as an upper bound for a woody seed, not as a measurement on one.
\end{assumption}

Assumption~\ref{as:est} is a limitation of the thesis, not a detail. A woody seed is harder to singulate than maize. Using maize and peanut detector scores inside an estimation makes the vision benefit look as good as those crops allowed. The estimation therefore supports a design comparison under published sensing performance. It does not forecast a field percentage for acacia or pine.

\subsection{Formal problems}

Three problems are stated. The first is the metering problem the vision loop must solve at a fixed speed. The second is the coverage problem the planner must solve before motion. The third is the joint problem, which is the one the machine actually has.

\begin{definition}[Metering problem]
\label{def:p1}
Given ground speed \(v(\cdot)\), target spacing \(x_{\mathrm{ref}}\), drop height \(H\), flexure parameters \((k,m_{\mathrm{eq}},\zeta)\) and a stream of images, force samples and inertial samples, compute plate speed \(\omega(t)\) and depth command \(z_{\mathrm{cmd}}(t)\) so that the expected indices satisfy
\begin{equation}
\label{eq:p1}
\E[\mis] \le \mis_{\max},\qquad
\E[\mul] \le \mul_{\max},\qquad
\E[(\,z - z^{\star}\,)^2] \le \sigma_{z,\max}^2,
\end{equation}
whenever \(v(t)\) remains inside \(\mathcal{V}_{\mathrm{meter}}\).
\end{definition}

\begin{definition}[Coverage problem]
\label{def:p2}
Given a field polygon, a digital elevation model, a swath width and a turning radius, choose a heading \(\beta\) and a reference \((x_d(s),y_d(s),v_d(s))\) along arc length \(s\) minimising
\begin{equation}
\label{eq:p2}
J = W_{\mathrm{mech}} + w_T T + w_s\, N_{\mathrm{seed}}\, \E[\rho_{\mathrm{fail}}],
\end{equation}
where \(W_{\mathrm{mech}}\) is the non-conservative mechanical work, \(T\) is the time to cover the polygon, \(N_{\mathrm{seed}}\) is the number of seeds the plan will release, and \(\rho_{\mathrm{fail}}\) is the probability that a release is a miss, a damaging multiple, or a depth error outside the band.
\end{definition}

\begin{definition}[Joint problem]
\label{def:p3}
Solve Definition~\ref{def:p2} subject to the speed remaining in the set for which Definition~\ref{def:p1} is feasible, including on the accelerations of every headland, and subject to the flexure magnification remaining below a declared bound \(M_{\max}\).
\end{definition}

The joint problem is deliberately not posed as a single monolithic program from the first sample of the day to the last. The heading is a slow decision. It can be computed before the machine moves and held for the field, or recomputed when the elevation model is updated. The plate command and the depth command are fast decisions. They must run at the hole period of Equation~\eqref{eq:period}. The steering and speed commands sit between them. Splitting the joint problem along those time scales is a design choice of Chapter~\ref{ch:design}, and it is justified by the time-scale separation already used in agricultural predictive control \cite{xu2022,plessen2017}. The formal risk of the split is sub-optimality: a heading chosen before the day starts does not see a patch of soil that shifts \(f_n\). The build therefore allows the predictive tracker to reduce speed inside a heading even when it cannot rotate the whole field plan in one sample.

The failure probability in Equation~\eqref{eq:p2} has to be a function the planner can evaluate without resimulating every seed. The structure used throughout the thesis is
\begin{equation}
\label{eq:fail}
\rho_{\mathrm{fail}} = 1 - \bigl(1 - \mis - \kappa\mul\bigr)\, S(\sigma_z),
\end{equation}
with \(\kappa\in(0,1]\) because a multiple wastes seed and crowds a neighbour but is not always a total loss, and with \(S\) a decreasing function of the depth root-mean-square error. Chapter~\ref{ch:estimate} takes \(\kappa=1/2\) and \(S(\sigma_z)=\exp(-\sigma_z^2/(2\sigma_0^2))\) with \(\sigma_0=\SI{6.5}{mm}\). Both choices are estimation devices. They encode the statement, taken from the depth literature \cite{masarei2019,bond1999}, that depth error and count error both destroy a hill, and that they do not destroy it in identical ways.

\subsection{Requirements}

The requirements below are written so that a later design choice can be checked against them without reinterpretation. A requirement that says ``accurate'' would not do that work.

\begin{enumerate}[label=\textbf{R\arabic*.}]
\item The quality-of-feed index on the design spacing, inside \(\mathcal{V}_{\mathrm{meter}}\), shall be at least \num{0.95} in the estimation model when the detector operates at the published precision and recall.
\item The depth root-mean-square error shall be at most \SI{4}{mm} inside the same set, against the \SI{18}{mm} target.
\item The vision-and-actuation delay shall be treated as a hard constraint on speed through Equation~\eqref{eq:latency}, not as a penalty that a large weight might override.
\item The resonant neighbourhood of the timber flexure, where \(M(r)>3\), shall be excluded from steady running or traversed only while accelerating through it.
\item The coverage cost shall include mechanical work, time and failed placements, and the failed-placement term shall use the same miss model as the metering loop.
\item Each row shall be sensed and metered independently. A single vehicle-level camera does not satisfy the requirement.
\item The predictive tracker shall run at a sample period no longer than \SI{100}{ms} on the computer specified in the build, using a convex quadratic program.
\item The estimation shall publish the open-loop baseline and the resonant-speed baseline beside the proposed plan, so that a gain is not compared with an unspecified alternative.
\item No result shall be described as a field percentage for a named species. The language of Chapter~\ref{ch:estimate} is the language of model evaluation.
\item The two technologies shall share the speed interface of Equation~\eqref{eq:intersect} and shall not each carry a private speed set-point.
\end{enumerate}

Requirement~R3 is the lesson of the \SI{6}{km/h} reversal \cite{borja2021}. Requirement~R4 is the lesson of Equation~\eqref{eq:mag} applied to a flexure that a steel planter would have pushed out of the frequency band. Requirement~R10 is the architectural claim of the thesis. A vision meter that assumes a constant \SI{1.5}{m/s} and a coverage planner that assumes the meter can tolerate whatever speed saves fuel are two correct components and one incorrect system.

\subsection{Research questions and the route through the remaining chapters}

The problem set reduces to two research questions, matched to the two technologies, plus the joint question the estimation is built to answer.

\begin{enumerate}[label=\textbf{Q\arabic*.}]
\item How should plate speed and boot depth be computed from images and from a timber-flexure model so that misses caused by empty cells are repaired only when Equation~\eqref{eq:latency} still holds, and so that depth error caused by arm motion is estimated rather than averaged away?
\item How should heading and the speed profile along a swath be computed so that non-conservative grade work, headland accelerations and the resonant speed of the same flexure appear in one cost, with the miss model of the metering loop evaluated inside that cost?
\item On a sloped field that punishes contour driving with extra turns and punishes uphill driving with extra work, what heading and speed does the joint problem select, and how do its indices compare with an open-loop meter at a steady working speed and at a resonant speed?
\end{enumerate}

Question~1 is answered by the derivation in Section~\ref{sec:tech1}. Question~2 is answered by the derivation in Section~\ref{sec:tech2}. Question~3 is answered by the build and the numerical estimation in Chapter~\ref{ch:estimate}. The answers are design answers and model answers. They become experimental answers only after a timber row unit, a detector trained on the woody seed, and a sloped field are instrumented to the same indices. The thesis stops at the point where those experiments have a specification to be built against, which is the ordinary role of a system-design argument.

The quantities that cross the chapter boundary are fixed here so that Chapter~\ref{ch:design} does not retune them silently. Equivalent stiffness at the boot is \(k=\SI{8200}{N/m}\), equivalent mass \(m_{\mathrm{eq}}=\SI{6.8}{kg}\), damping ratio \(\zeta=0.08\). Equation~\eqref{eq:fn} then gives \(f_n=\SI{5.53}{Hz}\). A static soil force of \SI{80}{N} deflects the boot by \(80/k=\SI{9.8}{mm}\), which is already outside the \(\pm\SI{8}{mm}\) band if it is left uncompensated, and which a depth servo is therefore required to remove. Roughness is represented, for the estimation, by a force component of amplitude \SI{30}{N} at wavelength \(\lambda=\SI{0.45}{m}\). At \(v=\SI{2.4}{m/s}\) the disturbance frequency is \(v/\lambda=\SI{5.33}{Hz}\), the frequency ratio is \(r=0.96\), and Equation~\eqref{eq:mag} gives a magnification near \(6\). That is the resonant neighbourhood Requirement~R4 forbids as a steady speed. Vision latency is taken as \(\tau_v=\SI{22}{ms}\) for a small detector on an embedded graphics processor, and actuator latency as \(\tau_a=\SI{35}{ms}\). Their sum, \SI{57}{ms}, is the delay Chapter~\ref{ch:design} inserts into Equation~\eqref{eq:latency}. These numbers are a specification, chosen to be physically plausible for laminated softwood and for published detector speeds. They are not a measurement of one prototype.

\begin{table}[ht]
\centering
\caption{Specification quantities fixed by the problem set and used without retuning in later chapters.}
\label{tab:spec}
\begin{tabular}{@{}lll@{}}
\toprule
Quantity & Value & Role \\
\midrule
Rows and swath & 4 rows, \SI{2.4}{m} & Coverage geometry \\
Cells \(N\), spacing \(x_{\mathrm{ref}}\) & 24, \SI{0.20}{m} & Plate feedforward \\
Drop \(H\), flight \(t_f\) & \SI{0.42}{m}, \SI{0.29}{s} & Spacing transport \\
Target depth & \SI{18}{mm}, band \(\pm\SI{8}{mm}\) & Depth loop \\
\(k\), \(m_{\mathrm{eq}}\), \(\zeta\), \(f_n\) & \SI{8200}{N/m}, \SI{6.8}{kg}, \num{0.08}, \SI{5.53}{Hz} & Flexure \\
Static force and deflection & \SI{80}{N}, \SI{9.8}{mm} & Bias the servo must remove \\
Roughness & \SI{30}{N} at \SI{0.45}{m} & Disturbance spectrum \\
\(\tau_v\), \(\tau_a\) & \SI{22}{ms}, \SI{35}{ms} & Latency inequality \\
Field & \SI{80}{m} by \SI{180}{m}, grade \num{8}\% & Estimation site \\
Vehicle mass & \SI{460}{kg} & Energy model \\
\bottomrule
\end{tabular}
\end{table}

The problem set is now closed enough to design against. The next chapter does not add a third objective. It constructs the metering technology that makes Definition~\ref{def:p1} feasible inside a declared speed set, and the coverage technology that minimises Definition~\ref{def:p2} while staying inside that set.

\chapter{Proposed Technologies}
\label{ch:design}

The problem set asks for two technologies that share a speed limit. This chapter defines them so that each can be implemented and so that each can be evaluated from the same equations the estimator in Chapter~\ref{ch:estimate} will call. Section~\ref{sec:tech1} is compliance-aware vision metering. It takes images, force and inertial measurements and returns a plate-speed command, a depth command and a set of speeds on which those commands are claimed to meet the index bounds. Section~\ref{sec:tech2} is singulation-constrained terrain model-predictive coverage. It takes the field, the elevation model and that speed set, and returns a heading and a speed profile. Neither section assumes the other has already been tuned on a prototype. Where a coefficient is numerical, it is one of the specification values in Table~\ref{tab:spec} or it is introduced here and carried unchanged into the estimation.

\section{Compliance-aware vision metering}
\label{sec:tech1}

\subsection{Role of the loop and the signals it may use}

Compliance-aware vision metering is a per-row controller. It does not steer the vehicle and it does not choose the heading. Its outputs are the plate angular speed \(\omega\), the depth-actuator set-point \(z_{\mathrm{cmd}}\), and a Boolean feasibility flag that is false when the current speed has left \(\mathcal{V}_{\mathrm{meter}}\). Its inputs are a camera frame of the plate and the upper tube, the boot force \(F_b\), the ultrasonic range \(r_u\) from a sensor on the boot to the soil, the arm pitch rate and vertical acceleration from an inertial sensor, and the ground speed \(v\) supplied by the vehicle. Ground speed is an input, not a result. That is the interface Requirement~R10 demanded.

The loop is organised in four blocks that run at different rates, because a detector and a depth servo do not want the same clock. The detector runs at the camera rate, taken as \SI{30}{Hz}, which at the design spacing and at \SI{1.5}{m/s} gives several frames per hole. A tracking filter over hole events runs at \SI{100}{Hz} so that it can timestamp a release more finely than the frame period. The depth servo runs at \SI{50}{Hz}. The feasibility flag is recomputed whenever \(v\) or the estimated vibration changes, and is sent to the coverage controller at \SI{10}{Hz}. Figure~\ref{fig:cavm} shows the blocks and the only path by which this technology is allowed to influence the vehicle: the feasible-speed message.

\begin{figure}[ht]
\centering
\begin{tikzpicture}[font=\small, node distance=0.45cm and 0.6cm, >=Stealth, box/.style={draw, rounded corners, align=center, minimum height=0.9cm, fill=blue!6}]
  \node[box] (cam) {camera\\plate and tube};
  \node[box, right=of cam] (det) {detector\\empty / single / double};
  \node[box, right=of det] (ukf) {flexure and\\release filter};
  \node[box, below=of det] (depth) {depth servo\\ \(z_{\mathrm{cmd}}\)};
  \node[box, below=of ukf] (plate) {plate law\\ \(\omega\)};
  \node[box, fill=orange!12, below=of cam] (flag) {feasible set\\ \(\mathcal{V}_{\mathrm{meter}}\)};
  \node[draw, dashed, rounded corners, align=center, right=1.1cm of plate, minimum height=1.5cm] (mpc) {SCT-MPC\\uses \(v\) only};
  \draw[->] (cam) -- (det);
  \draw[->] (det) -- (ukf);
  \draw[->] (ukf) -- (plate);
  \draw[->] (ukf) -- (depth);
  \draw[->] (ukf) -- (flag);
  \draw[->] (flag) -- (mpc);
  \draw[->] (mpc.north) |- node[pos=0.25, above] {\(v\)} (ukf.north);
\end{tikzpicture}
\caption{Compliance-aware vision metering. The vehicle is influenced only through the feasible-speed set. Plate speed and boot depth stay on the row unit.}
\label{fig:cavm}
\end{figure}

\subsection{Timber flexure as a process model}

The arm is reduced to a single degree of freedom, the vertical displacement \(\delta\) of the boot relative to the chassis reference, positive downward. With the equivalent stiffness and damping of the problem set,
\begin{equation}
\label{eq:sddot}
m_{\mathrm{eq}}\ddot{\delta} + c_d \dot{\delta} + k\delta = F_{\mathrm{soil}}(t) + F_{\mathrm{act}}(t),
\end{equation}
where \(c_d = 2\zeta\sqrt{k m_{\mathrm{eq}}}\) and \(F_{\mathrm{act}}\) is the force of the depth actuator. The actuator is modelled as a first-order lag on the command, with time constant \(\tau_z=\SI{80}{ms}\), which is deliberately slower than the resonant period. A depth loop that tries to be faster than \(f_n\) on a lightly damped flexure will need a lead term and a measurement of \(\dot{\delta}\); the inertial sensor provides that measurement, and the lag model keeps the design from claiming a bandwidth the wood cannot be asked for. Substituting the specification values gives \(c_d = 2\cdot 0.08\cdot\sqrt{8200\cdot 6.8}\approx \SI{37.8}{N\,s/m}\).

Soil force is split into a slow part and a roughness part,
\begin{equation}
\label{eq:fsoil}
F_{\mathrm{soil}} = F_0(z) + F_r \sin\bigl(2\pi s/\lambda + \phi\bigr),
\end{equation}
where \(s\) is the distance travelled and \(F_0\) depends on the nominal penetration. Godwin and O'Dogherty's family of models justifies a monotone increase of vertical force with depth \cite{godwin2007}. For the controller a local linearisation is enough,
\begin{equation}
\label{eq:flinear}
F_0(z) \approx F_0(z^{\star}) + k_{\mathrm{soil}}\bigl(z-z^{\star}\bigr),
\end{equation}
with \(k_{\mathrm{soil}}\) identified from a short calibration pass, not assumed from a national soil map. The ultrasonic range measures the distance from the boot datum to the surface. If the boot has sunk a depth \(z_{\mathrm{sink}}\) below the surface, the range and the flexure coordinate are related by a constant offset \(c_r\) that is calibrated when the machine is lowered onto a hard reference,
\begin{equation}
\label{eq:range}
r_u = c_r - \delta + n_r,
\end{equation}
where \(n_r\) is measurement noise. Equation~\eqref{eq:range} is what allows the filter to see \(\delta\) even when the inertial sensor has been integrated long enough to drift.

The transfer from a harmonic force of amplitude \(F_r\) to a displacement amplitude is \(F_r M(r)/k\), with \(M\) from Equation~\eqref{eq:mag}. This is the calculation that places the forbidden steady speed near \SI{2.4}{m/s} for \(\lambda=\SI{0.45}{m}\). It is also the calculation the filter uses as its process-noise schedule: when the estimated frequency ratio approaches one, the process noise on \(\delta\) is inflated, the filter trusts the ultrasonic range more than the open prediction, and the feasibility flag is cleared if the predicted depth variance exceeds \(\sigma_{z,\max}^2\).

\subsection{Seed state and the observation the camera actually makes}

A cell is a discrete object. The filter does not pretend otherwise by burying the cell index inside a purely continuous state. The continuous state at row \(i\), suppressing the row index until the multi-row note below, is
\begin{equation}
\label{eq:state}
\mathbf{x} =
\begin{bmatrix}
\delta & \dot{\delta} & z_{\mathrm{bias}} & t_{\mathrm{rel}} & v_{\mathrm{ex}}
\end{bmatrix}^{\mathsf{T}}.
\end{equation}
Here \(z_{\mathrm{bias}}\) is a slowly varying offset between the actuator command and the true depth, accounting for wear of the wooden pivot and for a change in surface residue. \(t_{\mathrm{rel}}\) is the predicted time of the next release relative to the current sample. \(v_{\mathrm{ex}}\) is the horizontal velocity of the seed at release in the soil frame, which Equation~\eqref{eq:carry} turns into a spacing error. The cell label \(\ell\in\{\mathrm{empty},\mathrm{single},\mathrm{double}\}\) is carried beside \(\mathbf{x}\) as a discrete mode, updated by the detector, not as a fifth continuous coordinate.

Between camera frames the continuous state is predicted by Equation~\eqref{eq:sddot} for the flexure coordinates, by \(\dot{z}_{\mathrm{bias}}=0\) plus a small process noise, and by
\begin{equation}
\label{eq:trel}
\dot{t}_{\mathrm{rel}} = -1,\qquad
\dot{v}_{\mathrm{ex}} = 0,
\end{equation}
with \(t_{\mathrm{rel}}\) reset to the predicted time of the next cell whenever a release is declared. The predicted release time of cell \(n\) under the feedforward plate speed is
\begin{equation}
\label{eq:tn}
\hat{t}_n = t_{n-1} + \frac{2\pi}{N\omega}.
\end{equation}
A nonzero \(\dot{\delta}\) at the gate changes the actual release. The correction used here is geometric. The gate is a lip of angle \(\gamma\) relative to the arm. A vertical arm velocity \(\dot{\delta}\) moves the lip across the seed and shifts the release time by
\begin{equation}
\label{eq:tshift}
\Delta t_{\mathrm{rel}} = \frac{\dot{\delta}\,\sin\gamma}{r_p\,\omega},
\end{equation}
where \(r_p\) is the radius of the cell circle. Equation~\eqref{eq:tshift} is the compliance-aware part of the vision loop. A detector that sees a seed still in the cell at a time when a rigid-body schedule expected it to have left will, without Equation~\eqref{eq:tshift}, declare a miss. With the equation, the filter first asks whether the arm velocity explains the delay. Only the residual delay is offered to the miss logic.

The detector returns a class and a confidence \(p_{\ell}\) for each visible cell. Following the published operating points, the estimation will credit the detector with a recall of \num{0.99} and a precision of \num{0.985} on the empty-versus-occupied decision when the window is clean \cite{blaser2024,zhao2025yolo}. The filter does not use those two numbers as if they were probabilities of the current frame. It uses the frame confidence, and it uses the published numbers only when Chapter~\ref{ch:estimate} averages over many seeds. A double is declared when two connected components, or two detector boxes, occupy one cell above a confidence threshold. The action is not symmetric with the empty-cell action. An empty cell can be repaired by advancing the next filled cell if Equation~\eqref{eq:latency} holds. A double already in the gate cannot be unreleased. The plate law below therefore slows the plate and raises a singulator duty, if a mechanical singulator is fitted, but it does not pretend to retrieve the extra seed.

\begin{figure}[ht]
\centering
\begin{tikzpicture}[font=\small, >=Stealth, scale=1.0]
  \draw[thick] (0,0) circle (1.7);
  \foreach \a in {0,15,...,345} {
    \fill[blue!30] (\a:1.35) circle (1.5pt);
  }
  \fill[red!70] (20:1.35) circle (2.4pt);
  \node[red!70!black] at (20:1.9) {empty};
  \draw[thick] (1.55,-0.35) -- (2.6,-1.3);
  \draw[thick] (1.85,-0.15) -- (2.9,-1.1);
  \node[align=left] at (4.0,-0.7) {wooden tube};
  \fill (2.75,-1.2) circle (1.5pt) node[right] {\ seed};
  \draw[->] (0,0) -- (0.9,0.5) node[midway, above left] {\(\omega\)};
  \draw[dashed] (0,0) -- (20:1.7);
  \node at (-2.3,1.3) {camera};
  \draw[->] (-1.7,1.15) -- (-0.4,0.55);
\end{tikzpicture}
\caption{Plate observation. The detector labels cells before they reach the gate. One empty cell is still repairable. A seed already in the wooden tube is only a measurement of release time, not a seed that can be called back.}
\label{fig:plate}
\end{figure}

\subsection{Unscented filter for a mixed observation}

The process model is nonlinear because of Equation~\eqref{eq:tshift} and because the ultrasonic geometry saturates when the boot leaves the soil. An extended Kalman filter would linearise both. An unscented filter propagates a small set of sigma points through the nonlinear map and matches the mean and covariance up to second order, which is the property Julier and Uhlmann designed it for \cite{julier1997}. Welch and Bishop's notes and S{\"a}rkk{\"a}'s text are the implementation references for the weights and the Joseph-stabilised covariance update \cite{welch2006,sarkka2013}. Thrun, Burgard and Fox discuss the same update in the setting of a robot that must stamp sensor events onto a moving base \cite{thrun2005}. The pattern here is closer to their event filter than to a simultaneous-localisation map: the ``landmark'' is the next cell, and it is consumed when the seed passes the gate.

Let \(n_x=5\) and let \(\lambda_{\mathrm{ukf}}=\alpha_{\mathrm{ukf}}^2(n_x+\kappa_{\mathrm{ukf}})-n_x\) with the usual \(\alpha_{\mathrm{ukf}}=10^{-2}\), \(\kappa_{\mathrm{ukf}}=0\). Given mean \(\hat{\mathbf{x}}\) and covariance \(P\), the sigma points are the mean and the symmetric columns of the matrix square root of \((n_x+\lambda_{\mathrm{ukf}})P\). Each point is pushed through Equations~\eqref{eq:sddot}--\eqref{eq:tshift} over one filter step \(h_f=\SI{10}{ms}\). The predicted mean and covariance are the weighted sums given by Julier and Uhlmann, with process noise \(Q\) diagonal. The diagonal entries are not constant. The entry on \(\delta\) is scaled by \(1+4\max(0, M(r)-1)^2\), so that a resonant arm is not assigned a calm process. The entry on \(t_{\mathrm{rel}}\) is scaled by the detector's reported confidence: a low-confidence frame is not allowed to pull the release time hard.

The observation vector stacks the quantities a sample actually contains,
\begin{equation}
\label{eq:obs}
\mathbf{y} =
\begin{bmatrix}
r_u & F_b & a_{\mathrm{imu}} & \mathbf{1}_{\mathrm{cross}}
\end{bmatrix}^{\mathsf{T}}.
\end{equation}
The indicator \(\mathbf{1}_{\mathrm{cross}}\) is one when the tube window sees a seed cross a fixed image row. Its predicted value is a smooth step of \(t_{\mathrm{rel}}\), \(\psi(t_{\mathrm{rel}})=(1+\exp(t_{\mathrm{rel}}/\tau_{\psi}))^{-1}\) with \(\tau_{\psi}=\SI{8}{ms}\), which restores differentiability for diagnostics even though the unscented filter does not need the Jacobian. The force observation is \(F_b = k\delta + c_d\dot{\delta} + n_F\), the inertial observation is the second derivative from Equation~\eqref{eq:sddot} divided by \(m_{\mathrm{eq}}\), and the range observation is Equation~\eqref{eq:range}. Standard unscented update equations apply. A cell label is applied after the continuous update, as a mode switch: if the label is empty and the residual of Equation~\eqref{eq:tshift} does not explain the absence, the discrete miss hypothesis is raised.

The initial covariance is set from the calibration pass: \(\delta\) within \(\SI{5}{mm}\), \(\dot{\delta}\) within \(\SI{0.05}{m/s}\), bias within \(\SI{3}{mm}\), release-time uncertainty equal to a quarter of the hole period, and exit-velocity uncertainty \(\SI{0.05}{m/s}\). Those values are priors for the filter, not results.

\subsection{Plate law and the condition for attempting a repair}

The plate command is the feedforward of Equation~\eqref{eq:rate} plus a correction that exists only when a repair is both necessary and timely,
\begin{equation}
\label{eq:plate}
\omega = \omega^{\star}(v) + \Delta\omega_{\mathrm{repair}} + \Delta\omega_{\mathrm{double}} + K_{\omega}\bigl(\hat{t}_{\mathrm{rel}} - t_{\mathrm{sched}}\bigr).
\end{equation}
The last term is a small proportional trim, with gain \(K_{\omega}\), that advances or retards the plate when the predicted release drifts from the schedule but the cell is occupied. It is saturated at \(\pm\num{8}\%\) of \(\omega^{\star}\) so that a noisy filter cannot double the seed rate. The double term \(\Delta\omega_{\mathrm{double}}\) is non-positive. It reduces plate speed for one hole period after a double is seen, on the reasoning that a double often comes from overspeed or from too much vacuum, and that the next hole should not be asked to repeat the condition. The vacuum itself, if the plate is pneumatic rather than a mechanical cell, is not placed under high-rate feedback in this design. Equation~\eqref{eq:vacuum} is used to set the nominal \(\Delta p\), including a margin for \(a_{\mathrm{vib}}\) evaluated at the largest \(M(r)\) allowed inside \(\mathcal{V}_{\mathrm{meter}}\). A slow outer trim, once per ten metres, may lower \(\Delta p\) if doubles dominate misses. Fast vacuum modulation is left out because a wooden hopper does not contain a pressure plenum that settles in one hole period.

The repair term is the one that can cause harm if it is late. Let \(\tau=\tau_v+\tau_a\) and let \(\Delta t=x_{\mathrm{ref}}/v\). Define the repair fraction
\begin{equation}
\label{eq:cfrac}
c(v,a) =
\begin{cases}
1 & \text{if }\Delta t \ge 1.35\,\tau + 0.012\,|a| + 0.030,\\
0.15 & \text{if }\Delta t \le 1.35\,\tau + 0.012\,|a|,\\
\text{linear in between.}
\end{cases}
\end{equation}
Equation~\eqref{eq:cfrac} is a design function, not a fit to a new experiment. The coefficient \(0.012\) seconds per metre-per-second-squared charges the extra time the stepper needs when the vehicle is accelerating and the scheduled \(\omega^{\star}\) is itself a moving target. The floor of \(0.15\) expresses Borja's observation that a late loop does not become a zero loop; it becomes a harmful one, so the controller should mostly stop repairing rather than repair at random \cite{borja2021}. The command is then
\begin{equation}
\label{eq:repair}
\Delta\omega_{\mathrm{repair}} =
\begin{cases}
\dfrac{2\pi}{N}\, c(v,a)\, \mathbf{1}_{\mathrm{empty}} & \text{over the next hole period,}\\
0 & \text{otherwise.}
\end{cases}
\end{equation}
When \(c=1\), an empty cell is replaced by bringing the following cell forward one pitch. When \(c\) is near the floor, the indicator is ignored and the plate stays on the feedforward. The feasibility set reported to the coverage controller is
\begin{equation}
\label{eq:vset}
\mathcal{V}_{\mathrm{meter}} = \bigl\{ v : c(v,0) \ge 0.85 \bigr\}
\setminus \bigl\{ v : M\bigl(v/(\lambda f_n)\bigr) > 3 \bigr\}.
\end{equation}
With the specification numbers, \(1.35\tau+0.030=\SI{0.107}{s}\) and \(x_{\mathrm{ref}}/v \ge 0.107\) gives \(v\le \SI{1.87}{m/s}\) for full repair authority with a margin. The magnification bound removes a neighbourhood of \SI{2.4}{m/s} even if some other spacing were to make the latency look acceptable. Both cuts are reported upward. The coverage controller does not re-derive them.

\begin{algorithm}[ht]
\caption{Plate and feasibility update on one row}
\label{alg:plate}
\begin{algorithmic}[1]
\Require speed \(v\), acceleration \(a\), filter state, detector label \(\ell\), confidences
\State \(\omega^{\star} \gets 2\pi v / (N x_{\mathrm{ref}})\)
\State predict \(\mathbf{x}\) through the flexure and Equation~\eqref{eq:tshift}
\State update \(\mathbf{x}\) with \(\mathbf{y}\) by the unscented step
\State \(c \gets c(v,a)\) from Equation~\eqref{eq:cfrac}
\If{\(\ell=\mathrm{empty}\) and residual delay is unexplained and \(c\ge 0.85\)}
  \State apply one-pitch advance over the next hole period
\ElsIf{\(\ell=\mathrm{double}\)}
  \State \(\omega \gets \omega^{\star} - \Delta\omega_{\mathrm{double}}\) for one hole period
\Else
  \State \(\omega \gets \omega^{\star} + K_{\omega}(\hat{t}_{\mathrm{rel}}-t_{\mathrm{sched}})\), then saturate
\EndIf
\State publish \(\omega\), and publish whether \(v\in\mathcal{V}_{\mathrm{meter}}\)
\end{algorithmic}
\end{algorithm}

\subsection{Depth servo}

The depth loop is a proportional-integral law on the filtered boot displacement, with a damping term on \(\dot{\delta}\),
\begin{equation}
\label{eq:depth}
\begin{aligned}
e_z &= z^{\star} - \bigl(z_{\mathrm{ref,arm}} + \hat{\delta} + \hat{z}_{\mathrm{bias}}\bigr),\\
F_{\mathrm{act}} &= K_p e_z + K_i \int e_z\,dt - K_d \hat{\dot{\delta}}.
\end{aligned}
\end{equation}
The integral state is clamped so that a long lift of the machine, when the boot is in the air and the ultrasonic return is invalid, does not wind up a force that slams the boot into the soil on the next entry. The integral is also frozen whenever the ultrasonic innovation exceeds a gate, which is the filter's way of saying that residue or a stone has produced a false range. Gains are chosen by loop shaping on Equation~\eqref{eq:sddot}. The target closed-loop bandwidth is \SI{3}{Hz}, below \(f_n=\SI{5.53}{Hz}\), so the servo removes the \SI{9.8}{mm} static deflection and the slow part of the roughness, and it does not attempt to cancel the resonant peak. Cancelling that peak with actuator force would require a bandwidth beyond the wood and would pump energy into \(f_n\). The resonant peak is rejected by the speed constraint instead. This division of labour is the central design choice of the first technology.

A short Lyapunov argument records why the constrained servo does not fight the flexure. Take the flexure energy
\begin{equation}
\label{eq:energy}
V(\delta,\dot{\delta}) = \tfrac12 m_{\mathrm{eq}}\dot{\delta}^2 + \tfrac12 k\delta^2
\end{equation}
and the actuator policy of Equation~\eqref{eq:depth} with the integral held constant. Differentiating along Equation~\eqref{eq:sddot} yields
\begin{equation}
\label{eq:vdot}
\dot{V} = -c_d\dot{\delta}^2 + \dot{\delta}\, F_{\mathrm{soil}} + \dot{\delta}\, F_{\mathrm{act}}.
\end{equation}
The damping term is negative semidefinite. \(F_{\mathrm{act}}\) contributes \(\dot{\delta}(K_p e_z - K_d\dot{\delta})\) plus the integral contribution. The \(-K_d\dot{\delta}^2\) term augments wood damping, which is the reason for the rate feedback. The proportional term exchanges energy with the depth error and is bounded because \(F_{\mathrm{act}}\) is saturated at the force of a small linear actuator, taken as \SI{120}{N} in the build. Saturation prevents a noisy \(\hat{\delta}\) from commanding more force than the timber pivot is designed to carry. Asymptotic convergence to a point is not claimed, because \(F_{\mathrm{soil}}\) is persistently exciting. What is claimed, and what the estimation checks, is that the steady component of \(\delta\) is driven inside the band and that the resonant component is left to the speed constraint.

\begin{proposition}
\label{prop:split}
Under Assumption~\ref{as:flex}, a depth loop of bandwidth below \(f_n\) together with the exclusion of \(\{v:M(r)>3\}\) from \(\mathcal{V}_{\mathrm{meter}}\) keeps the resonant mode from being both strongly excited and strongly actuated. The loop is responsible for the mean depth. The vehicle speed is responsible for the dynamic magnification.
\end{proposition}

Proposition~\ref{prop:split} is a design proposition. Its content is the separation, not a new stability theorem. Khalil's treatment of singular perturbation is the justification for controlling the slow depth error while treating the fast resonant response as a disturbance to be avoided rather than inverted \cite{khalil2002}.

\subsection{Predicted indices at a constant speed}

The coverage layer needs a miss probability it can evaluate without running the filter. The map used for that purpose is an explicit function, constructed so that its shape agrees with the qualitative findings of Section~\ref{sec:prev-meter}: misses rise with speed, rise again with acceleration, and rise when the arm magnification grows. For the open-loop plate, which is the baseline,
\begin{equation}
\label{eq:openmiss}
\begin{aligned}
\mis_{\mathrm{ol}}(v,a) &= \mathrm{clip}_{[0,0.55]}\Bigl(
0.055 + 0.085\max(0,v-0.7) + 0.040|a| \\
&\qquad + 0.0045\max\bigl(0,\, d_{\mathrm{mm}}(v)-3\bigr)\Bigr),\\
\mul_{\mathrm{ol}}(v,a) &= \mathrm{clip}_{[0,0.30]}\Bigl(
0.028 + 0.022\max(0,v-0.8) + 0.015|a|\Bigr),
\end{aligned}
\end{equation}
where \(d_{\mathrm{mm}}(v) = 10^3 F_r M(r)/k\) is the dynamic displacement amplitude in millimetres and \(\mathrm{clip}\) denotes saturation to the interval. Equation~\eqref{eq:openmiss} is not identified from a woody-seed trial. The intercept and the slopes are chosen so that, at ordinary planting speeds, the open-loop quality-of-feed index sits in the \numrange{0.75}{0.90} range that irregular seed produces on a plate that was not tuned for it, and so that the acceleration term is visible, as Mangus et al.\ found it to be \cite{mangus2017}. Chapter~\ref{ch:estimate} states this status again next to every table that uses the equation.

The closed-loop indices apply the repair fraction and a double reduction, and they add a small floor that represents false repairs,
\begin{equation}
\label{eq:closedmiss}
\begin{aligned}
\mis_{\mathrm{cl}} &= \mis_{\mathrm{ol}}\bigl(1 - 0.99\cdot 0.97\cdot c(v,a)\bigr) + 0.003,\\
\mul_{\mathrm{cl}} &= \mul_{\mathrm{ol}}\bigl(1 - 0.75\, c(v,a)\bigr) + 0.004.
\end{aligned}
\end{equation}
The product \(0.99\cdot 0.97\) is the estimation's way of inserting published recall and a further loss for cells that are empty for a mechanical reason the detector sees but the stepper cannot fill, such as an empty hopper sector. The factor \(0.75\) on doubles expresses that seeing a double prevents some, not all, of the subsequent doubles. Depth error for the two loops is
\begin{equation}
\label{eq:deptherr}
\begin{aligned}
\sigma_{z,\mathrm{ol}} &= \sqrt{d_{\mathrm{mm}}(v)^2 + 4^2},\\
\sigma_{z,\mathrm{cl}} &= \sqrt{\bigl(d_{\mathrm{mm}}(v)\,\eta_z(v)\bigr)^2 + 1^2},
\end{aligned}
\end{equation}
where the \(4\)~mm and \(1\)~mm terms are a bias that the open loop leaves in place and a residual bias the integral term does not fully remove. The rejection factor \(\eta_z\) equals \(0.30+0.08(v/2)\) when the disturbance frequency is below \SI{3.2}{Hz}, and it worsens linearly once the roughness frequency enters the band the \SI{3}{Hz} servo cannot track, up to a cap of \num{0.85}. Equation~\eqref{eq:deptherr} is how Proposition~\ref{prop:split} appears in a number: below the servo bandwidth the closed loop is credited with a large reduction, and near resonance it is not.

\begin{figure}[ht]
\centering
\begin{tikzpicture}
\begin{axis}[
  width=0.92\textwidth, height=6.6cm,
  xlabel={Ground speed \(v\) (m/s)},
  ylabel={Quality of feed index},
  xmin=0.7, xmax=2.9, ymin=0.6, ymax=1.02,
  grid=both, legend pos=south west, legend style={font=\small}
]
\addplot[thick, mark=*] coordinates {
(0.8,0.904) (1.0,0.881) (1.2,0.858) (1.4,0.834) (1.6,0.809)
(1.8,0.781) (2.0,0.750) (2.2,0.709) (2.4,0.653) (2.6,0.643) (2.8,0.657)
};
\addplot[thick, mark=square*] coordinates {
(0.8,0.983) (1.0,0.981) (1.2,0.980) (1.4,0.978) (1.6,0.976)
(1.8,0.974) (2.0,0.927) (2.2,0.848) (2.4,0.752) (2.6,0.686) (2.8,0.697)
};
\legend{Open-loop model, CAVM model}
\end{axis}
\end{tikzpicture}
\caption{Modelled quality-of-feed index against speed at zero acceleration, from Equations~\eqref{eq:openmiss} and \eqref{eq:closedmiss}. The closed-loop curve holds near \num{0.98} while repairs are timely and falls once Equation~\eqref{eq:cfrac} withdraws repair authority, then falls again through the timber resonance near \SI{2.4}{m/s}.}
\label{fig:qfi}
\end{figure}

Figure~\ref{fig:qfi} is a plot of the model, not of a trial. Two features are the ones the second technology will exploit. First, there is a wide plateau on which spending effort to go slightly faster buys almost no singulation loss and, as the energy model will show, almost no energy change, until the plateau ends. Second, the end is sharp. A planner that treats singulation as a linear penalty in speed will happily ride up to the cliff. A planner that treats \(\mathcal{V}_{\mathrm{meter}}\) as a hard set will stop at the plateau. Requirement~R3 chooses the hard set.

\subsection{Four rows, clocking and what the filter does not estimate}

The four row units are four copies of the same filter. They share the vehicle speed and the sample clock, and nothing else in the nominal design. Sharing the clock matters because a repair advance on row~1 and a repair advance on row~3 must not be queued on one processor so that the second repair misses its gate. The build therefore gives each row a detector thread with a reserved budget of \SI{22}{ms}, and it drops a frame rather than extending \(\tau_v\). A dropped frame lowers confidence, inflates the release-time process noise, and withholds repair for that hole. Withholding repair is the safe side of Borja's reversal.

Chassis coupling, neglected by Assumption~\ref{as:flex}, would appear as a common \(\delta\) component on all four rows when the whole machine pitches. The inertial sensors will see that common mode. A later revision can subtract the mean of the four \(\dot{\delta}\) estimates before the repair law runs. The present design does not. The estimation treats the omission as an unmodelled disturbance inside the process noise already scheduled by \(M(r)\), and it does not claim the four rows are independent in a statistical test. They are independent in the software architecture.

Seed identity across a double is not estimated. The filter counts a double. It does not decide which of the two seeds is the one that should have been released. For a single-species woody seed the distinction does not change the command. For a mixed-species restoration hopper it would, and that case is outside Assumption~\ref{as:seed}.

\subsection{A numerical pass through one empty cell}

The plate law is easier to audit on one hole than in the general notation. Take the design point \(v=\SI{1.4}{m/s}\), \(a=0\), \(x_{\mathrm{ref}}=\SI{0.20}{m}\) and \(N=24\). The hole period is
\begin{equation}
\label{eq:dtnum}
\Delta t = \frac{0.20}{1.4} = \SI{0.143}{s}.
\end{equation}
The feedforward plate speed is
\begin{equation}
\label{eq:omeganum}
\omega^{\star} = \frac{2\pi \cdot 1.4}{24 \cdot 0.20} = \SI{1.833}{rad/s},
\end{equation}
which is \num{17.5} revolutions per minute. One cell pitch is \(2\pi/24=\SI{0.262}{rad}\). The latency sum is \(\tau=\SI{0.057}{s}\), so
\begin{equation}
\label{eq:cnum}
1.35\tau + 0.030 = \SI{0.107}{s} < 0.143~\text{s},
\end{equation}
and Equation~\eqref{eq:cfrac} returns \(c=1\). A repair is allowed.

Suppose the detector, at a frame whose exposure began \SI{90}{ms} before the empty cell reaches the gate, returns an empty label with high confidence, and the flexure residual does not explain the absence. The remaining time before the gate, after inference, is about \(90-22=\SI{68}{ms}\). The stepper must gain one pitch in that window if the following cell is to arrive at the original gate time. The extra angular speed, held for those \SI{68}{ms}, is
\begin{equation}
\label{eq:domega}
\Delta\omega = \frac{0.262}{0.068} = \SI{3.85}{rad/s},
\end{equation}
about twice the feedforward. The actuator budget \(\tau_a=\SI{35}{ms}\) is the time to begin the move, not the time to finish a two-fold speed change if the following cell is still far away. In this numerical hole the following cell is one pitch behind, so the move is exactly the one-pitch advance of Equation~\eqref{eq:repair}, and \SI{68}{ms} is longer than \SI{35}{ms}. The repair fits. The same arithmetic at \(v=\SI{2.2}{m/s}\) gives \(\Delta t=\SI{0.091}{s}\) and \(c=0.55\), and the plate law withholds the advance. The empty cell becomes a miss that is logged and is not ``fixed'' by a late seed dropped into the next interval.

The flight that follows a successful release is independent of that logic and still has to be accounted for in the seed map. With \(H=\SI{0.42}{m}\) and negligible vertical release speed,
\begin{equation}
\label{eq:tfnum}
t_f = \sqrt{2\cdot 0.42/9.81} = \SI{0.293}{s}.
\end{equation}
At \SI{1.4}{m/s} the machine travels \(1.4\cdot 0.293=\SI{0.410}{m}\) during the flight. The tube and the furrow closer are built to cancel that mean transport, so the nominal landing sits on the target spacing. What they cannot cancel is a gate velocity that changes from seed to seed. An arm velocity \(\dot{\delta}=\SI{0.04}{m/s}\), which is a modest fraction of the resonant motion in Table~\ref{tab:speed}'s dynamic amplitude, shifts the landing by
\begin{equation}
\label{eq:jitternum}
\dot{\delta}\, t_f = 0.04\cdot 0.293 = \SI{11.7}{mm}.
\end{equation}
Relative to a \SI{200}{mm} spacing that jitter is a coefficient of variation contribution of about \num{0.06} if it is the dominant random term, which is inside the single class of Equation~\eqref{eq:classes} and is exactly why \(\mathrm{CV}_s\) is reported separately from \(\qfi\). The filter's \(v_{\mathrm{ex}}\) state exists so that this \SI{11.7}{mm} is written into the seed map rather than absorbed into a story about the plate having been wrong.

Depth on the same hole uses the static and dynamic pieces separately, as Proposition~\ref{prop:split} requires. The \SI{80}{N} mean soil force on \(k=\SI{8200}{N/m}\) is a \SI{9.8}{mm} offset. The integral term in Equation~\eqref{eq:depth} is what removes it; the open-loop column of the later tables leaves a \SI{4}{mm} remnant of that offset as the bias inside Equation~\eqref{eq:deptherr}. The roughness amplitude at this speed is \(d_{\mathrm{mm}}(1.4)=\SI{5.31}{mm}\). The servo rejection factor at a disturbance frequency \(1.4/0.45=\SI{3.11}{Hz}\), just under the \SI{3.2}{Hz} knee, is
\begin{equation}
\label{eq:etaz}
\eta_z = 0.30 + 0.08\cdot(1.4/2) = 0.356,
\end{equation}
so the rejected dynamic part is \(5.31\cdot 0.356=\SI{1.89}{mm}\). Combined with the \SI{1}{mm} residual bias,
\begin{equation}
\label{eq:sznum}
\sigma_{z,\mathrm{cl}} = \sqrt{1.89^2 + 1^2} = \SI{2.14}{mm},
\end{equation}
which is the closed-loop entry at \SI{1.4}{m/s} in the speed table of Chapter~\ref{ch:estimate}. The point of writing it here is that the table is not a separate experiment. It is Equation~\eqref{eq:sznum} evaluated on a grid.

\subsection{One unscented step on a reduced flexure state}

The five-dimensional filter is the one that runs. A two-dimensional reduction is the one a reader can check by hand, and it uses the same weights. Keep only \(\delta\) and \(\dot{\delta}\), with a predicted mean \(\hat{\delta}=\SI{2}{mm}\), \(\hat{\dot{\delta}}=\SI{0.01}{m/s}\), and a diagonal covariance \(P=\diag(10^{-6},\, 4\cdot 10^{-4})\) in metres and metres per second. Then \(n_x=2\), \(\lambda_{\mathrm{ukf}}=\alpha_{\mathrm{ukf}}^2 n_x - n_x\) with \(\alpha_{\mathrm{ukf}}=10^{-2}\), so \(\lambda_{\mathrm{ukf}}\approx -1.9998\) and \(n_x+\lambda_{\mathrm{ukf}}\approx 2\cdot 10^{-4}\). The sigma-point spread is small, which is what a small \(\alpha_{\mathrm{ukf}}\) is for: the points stay near the mean and the approximation remains local \cite{julier1997}. Over a \SI{10}{ms} step with \(F_{\mathrm{act}}\) cancelling the mean of \(F_{\mathrm{soil}}\) and with the linear stiffness, each point moves by
\begin{equation}
\label{eq:euler}
\begin{aligned}
\delta^{+} &= \delta + h_f \dot{\delta},\\
\dot{\delta}^{+} &= \dot{\delta} + h_f m_{\mathrm{eq}}^{-1}\bigl(F_{\mathrm{soil}}-c_d\dot{\delta}-k\delta\bigr).
\end{aligned}
\end{equation}
With \(h_f=\SI{0.01}{s}\) the mean displacement advances by \(10^{-4}\)~m, a tenth of a millimetre, and the mean velocity is pulled back toward zero by the stiffness term \(k\hat{\delta}/m_{\mathrm{eq}} = 8200\cdot 0.002/6.8 \approx \SI{2.41}{m/s^2}\), times \(h_f\), a change of about \(\SI{0.024}{m/s}\) before damping. Damping \(c_d\dot{\delta}/m_{\mathrm{eq}}\) at the predicted velocity is only about \(\SI{0.056}{m/s^2}\), so on this particular step stiffness dominates damping. The numerical content of the step is ordinary. What the unscented construction adds, and what an extended filter would approximate with one Jacobian, is that a sigma point sitting \SI{1}{mm} higher sees a different stiffness force and the predicted covariance of velocity grows in the direction the stiffness pushes. Process noise is then added on the diagonal, inflated by the magnification schedule when \(M(r)\) is large, and the ultrasonic update pulls \(\delta\) toward Equation~\eqref{eq:range}. The implementation note that matters for the build is the period: this step is cheap beside the detector, and it is the reason the filter runs at \SI{100}{Hz} while the camera runs at \SI{30}{Hz}. Between frames the release-time coordinate simply counts down, as Equation~\eqref{eq:trel} says, and the flexure coordinates keep integrating force.

\subsection{Vacuum margin and the cases the plate law refuses}

If the plate is pneumatic, Equation~\eqref{eq:vacuum} is evaluated at the edge of \(\mathcal{V}_{\mathrm{meter}}\) rather than at the working speed, so that a brief excursion inside the feasible set does not drop seeds the nominal vacuum cannot hold. Take a design seed of mass \(m_s=\SI{18}{mg}=\SI{1.8e-5}{kg}\) and an orifice area \(A_o=\SI{4}{mm^2}=\SI{4e-6}{m^2}\). The ratio \(m_s/A_o=\SI{4.5}{kg/m^2}\). At plate radius \(r_p=\SI{0.12}{m}\) and \(\omega^{\star}=\SI{1.833}{rad/s}\), the centripetal acceleration is \(r_p\omega^2\approx\SI{0.40}{m/s^2}\), small beside gravity. The vibration term is not small beside gravity once magnification grows. An arm acceleration of order \((2\pi f)^2\) times a few millimetres, at \(f=\SI{3}{Hz}\) and amplitude \SI{5}{mm}, is
\begin{equation}
\label{eq:avib}
a_{\mathrm{vib}} \approx (2\pi\cdot 3)^2 \cdot 0.005 \approx \SI{1.8}{m/s^2},
\end{equation}
about \(0.18\,g\). The holding pressure at the plateau edge is then of order
\begin{equation}
\label{eq:dpnum}
\Delta p \ge 4.5 \cdot (9.81 + 0.40 + 1.8) \approx \SI{54}{Pa}
\end{equation}
before orifice efficiency, seal leakage and the extra factor a real vacuum model puts on seed shape \cite{karayel2004}. Fifty pascals is a small gauge pressure. The practical content of the calculation is the opposite of its magnitude. Because the required pressure is small, operators turn the blower up until doubles appear, which is the ridge Singh, Singh and Saraswat mapped for cotton \cite{singh2005}. The vision loop's double term exists because the pressure that guarantees a pickup of the lightest seed in the lot also picks up a second seed when two small seeds cover one orifice. The nominal \(\Delta p\) is set with the \(a_{\mathrm{vib}}\) margin included, and the slow outer trim is allowed to lower it only when the logged double rate exceeds the logged miss rate for ten metres. It is not allowed to chase a single double with a pressure step inside one hole period. A wooden hopper's plenum, if a plenum is fitted at all, does not settle that fast.

The plate law also refuses three situations that look, from a distance, like opportunities to be clever. It refuses a repair when the hopper-level sensor on that row is low, because advancing the plate cannot fill a cell from an empty sector and the advance would only desynchronise the schedule. It refuses a repair when the realised delay, measured from exposure stamp to stepper issue, has exceeded the budget for three consecutive holes, because the \(c\) computed from the specification would otherwise keep authorising repairs the hardware is no longer fast enough to perform. It refuses to increase vacuum and plate speed in the same hole. Each of those actions can reduce misses and each of them can manufacture doubles. Stacking them on one observation is how a correct detector produces the sign change Borja measured \cite{borja2021}. One action per hole, and no action when \(c\) is low, is the rule the numerical hole above was written to illustrate.

\section{Singulation-constrained terrain model-predictive coverage}
\label{sec:tech2}

\subsection{Decisions the coverage technology is allowed to make}

Singulation-constrained terrain model-predictive coverage makes two decisions. Before work, or whenever the machine stops long enough to replan, it chooses a heading \(\beta\) relative to the uphill direction and a steady working speed \(v_{\mathrm{work}}\in\mathcal{V}_{\mathrm{meter}}\). During work it tracks the resulting swaths and is allowed to vary speed inside the horizon, including a reduction below \(v_{\mathrm{work}}\) when a headland acceleration would otherwise push \(c(v,a)\) down or when a local grade would push the flexure. It is not allowed to exceed the upper end of \(\mathcal{V}_{\mathrm{meter}}\), and it is not allowed to retune the miss model. The miss model is Equations~\eqref{eq:openmiss}--\eqref{eq:deptherr}, evaluated with the closed-loop branch because the meter on the machine is the meter of Section~\ref{sec:tech1}.

The technology is split into an outer coverage program and an inner quadratic-program tracker, following the time-scale split already argued for the joint problem. The outer program is allowed seconds. The inner program is required to finish inside the \SI{100}{ms} of Requirement~R7.

\subsection{Work model on a grade}

The mechanical work in Definition~\ref{def:p2} counts only the non-conservative part. A vehicle that drives up a slope and later drives down it returns the gravitational potential. What it does not return is the work against rolling resistance, the work dissipated in slip, and the aerodynamic drag, which at these speeds is small but is kept so that the model does not claim drag is zero. With vehicle mass \(M\), rolling coefficient \(C_{rr}\), a slip coefficient \(C_s\), drivetrain efficiency \(\eta_d\), air density \(\rho\), and drag area \(C_dA\),
\begin{equation}
\label{eq:fpm}
\begin{aligned}
\Gamma(\beta) &= 1 + \gamma\cos^2\beta,\\
F(v,\beta) &= C_{rr} Mg\cos\alpha + C_s Mg\sin^2\alpha\;\Gamma(\beta) + \tfrac12\rho\, C_dA\, v^2,\\
P(v,\beta) &= F(v,\beta)\, v / \eta_d .
\end{aligned}
\end{equation}
The heading \(\beta=0\) points uphill and \(\beta=\ang{90}\) lies on the contour. The factor \(\Gamma\) is larger when the vehicle points up the slope, because that is when the tangential weight, and therefore the slip loss, is organised against the motion rather than across it. The coefficients used in the estimation are \(M=\SI{460}{kg}\), \(C_{rr}=0.085\), \(C_s=2.0\), \(\gamma=1.15\), \(\eta_d=0.70\), \(\rho=\SI{1.15}{kg/m^3}\) and \(C_dA=\SI{0.85}{m^2}\). They are chosen so that a change of heading on an \num{8}\% grade changes work by several percent, which is the same order as the \numrange{9.5}{15.6}\% direction effect reported on the fields of Vahdanjoo, Gislum and S{\o}rensen, without copying a percentage from those fields onto this rectangle \cite{vahdanjoo2023}. A softer tyre or a wetter soil would raise \(C_{rr}\) and \(C_s\) together. The estimation does not sweep that change; it records the coefficients so that a later bench test of rolling resistance can replace them without touching the optimiser.

Headland distance is not free. A turn of radius \(R_{\mathrm{turn}}=\SI{4.2}{m}\) at the end of each swath contributes length \(\pi R_{\mathrm{turn}}\) and is driven at \SI{0.8}{m/s}, with an extra \num{8}\% on the force to represent the side-slip of the turn. If a heading produces \(n\) swaths of length \(L_{\mathrm{row}}\),
\begin{equation}
\label{eq:work}
W_{\mathrm{mech}} = \frac{F(v_{\mathrm{work}},\beta)}{\eta_d}\, n L_{\mathrm{row}} + 1.08\,\frac{F(0.8,\beta)}{\eta_d}\, n\pi R_{\mathrm{turn}}.
\end{equation}
Time is \(T = n L_{\mathrm{row}}/v_{\mathrm{work}} + n\pi R_{\mathrm{turn}}/0.8\). The number of swaths is \(n=\lceil S_{\perp}/w\rceil\) with swath width \(w=\SI{2.4}{m}\) and \(S_{\perp}\) the field width measured across the heading. On the estimation field the uphill heading has long rows and \(n=34\) turns, while the contour heading has short rows and \(n=75\) turns. Equation~\eqref{eq:work} therefore punishes the contour with turn work and turn time, and punishes the uphill heading with \(\Gamma(0)=1+\gamma\). That is the trade the outer program exists to settle. A planner that minimises only \(W_{\mathrm{mech}}\) will lean toward the contour. A planner that also counts failed seeds during the accelerations at each of \(75\) turns will lean back.

\begin{figure}[ht]
\centering
\begin{tikzpicture}[scale=0.95, font=\small, >=Stealth]
  \draw[thick] (0,0) rectangle (3.2,7.2);
  \node[rotate=90] at (-0.45,3.6) {uphill, long rows};
  \foreach \x in {0.3,0.7,...,2.9} {
    \draw[blue!70!black] (\x,0.25) -- (\x,6.95);
  }
  \draw[->] (3.5,0.3) -- (3.5,1.5) node[midway, right] {grade};
  \begin{scope}[xshift=6.4cm]
    \draw[thick] (0,0) rectangle (7.2,3.2);
    \node at (3.6,-0.45) {contour, short rows, many turns};
    \foreach \y in {0.28,0.68,...,2.9} {
      \draw[orange!80!black] (0.25,\y) -- (6.95,\y);
    }
    \draw[->] (7.5,0.2) -- (7.5,1.2) node[midway, right] {grade};
  \end{scope}
\end{tikzpicture}
\caption{The same rectangle covered along the slope and along the contour. Long uphill rows pay a sustained slip loss. Short contour rows pay a turn, and an acceleration, at every swath.}
\label{fig:field}
\end{figure}

The distance spent accelerating at each end of a swath is estimated as \(v_{\mathrm{work}}^2/(2a_{\mathrm{turn}})\) on entry and the same on exit. With \(a_{\mathrm{turn}}=\SI{0.7}{m/s^2}\), the fraction of row length spent at elevated acceleration is
\begin{equation}
\label{eq:frac}
\phi_{\mathrm{acc}} = \min\Bigl(1,\; \frac{n\, v_{\mathrm{work}}^2 / a_{\mathrm{turn}}}{n L_{\mathrm{row}}}\Bigr)
= \min\Bigl(1,\; \frac{v_{\mathrm{work}}^2}{a_{\mathrm{turn}} L_{\mathrm{row}}}\Bigr).
\end{equation}
Short rows make \(\phi_{\mathrm{acc}}\) large. The indices that enter the cost are the mixture
\begin{equation}
\label{eq:mix}
\mis = (1-\phi_{\mathrm{acc}})\,\mis_{\mathrm{cl}}(v,0) + \phi_{\mathrm{acc}}\,\mis_{\mathrm{cl}}(v,a_{\mathrm{turn}}),
\end{equation}
and likewise for multiples and for the depth error, with an extra \SI{0.6}{mm} on the depth error during the accelerated fraction to represent pitch into the headland. Equation~\eqref{eq:fail} then supplies \(\rho_{\mathrm{fail}}\). The seed count is \(N_{\mathrm{seed}} = 4\, n L_{\mathrm{row}} / x_{\mathrm{ref}}\), four rows along the same swath length.

\subsection{Outer program}

The outer program searches a finite set. Headings are taken at five-degree steps from \ang{0} to \ang{90}, using the symmetry of a rectangular field with a single slope direction. Speeds are taken at \SI{0.1}{m/s} steps inside \(\mathcal{V}_{\mathrm{meter}}\), with the resonant interval \([\SI{2.15}{m/s},\SI{2.55}{m/s}]\) removed even if a future change of \(x_{\mathrm{ref}}\) were to widen the latency set. For each candidate the cost of Definition~\ref{def:p2} is evaluated in closed form from Equations~\eqref{eq:work}--\eqref{eq:mix}. Infeasible candidates, those with modelled \(\qfi<0.955\) or \(\sigma_z>\SI{4}{mm}\), are discarded before the cost is compared. This is a lexicographic filter in front of a scalar cost, and it implements Requirements~R1 and~R2 as constraints rather than as terms that a weight could bargain away.

The weights used for the scalar comparison are \(w_T=\SI{0.55}{MJ/h}\) and \(w_s\) such that the failure term enters as \(4.0\,\rho_{\mathrm{fail}}\) in megajoule units on this field. The weights are planning preferences. They say that an hour of machine time is worth a fraction of a megajoule of the work budget, and that a failed placement is expensive relative to the energy differences, which are themselves only a few tenths of a megajoule. A user who prices contractor time higher can raise \(w_T\). The constraint filter does not move when the weights move. That separation is intentional. It keeps a cheap hour from buying a resonant speed.

For a non-rectangular field the same cost is evaluated along a coverage produced by the trapezoidal decomposition of Oksanen and Visala, or by the three-dimensional decomposition of Jin and Tang when the elevation model is not a single plane \cite{oksanen2009,jin2011}. The rectangle in the estimation is a plane of constant \(\alpha\), so the decomposition adds no region boundary and the heading search is the whole outer program. The algorithm is written for the general case so that a polygonal field does not require a different controller, only a different enumeration of swaths.

\begin{algorithm}[ht]
\caption{Heading and working-speed selection}
\label{alg:outer}
\begin{algorithmic}[1]
\Require polygon, elevation model, \(\mathcal{V}_{\mathrm{meter}}\), weights, index limits
\State build a swath set for each heading \(\beta\) in the candidate grid
\For{each \((\beta, v)\) with \(v\in\mathcal{V}_{\mathrm{meter}}\) outside the resonant interval}
  \State compute \(W_{\mathrm{mech}}\), \(T\), \(\phi_{\mathrm{acc}}\), \(\mis\), \(\mul\), \(\sigma_z\), \(\rho_{\mathrm{fail}}\)
  \If{\(\qfi \ge 0.955\) and \(\sigma_z \le \SI{4}{mm}\)}
    \State \(J \gets W_{\mathrm{mech}} + w_T T + w_s N_{\mathrm{seed}}\rho_{\mathrm{fail}}\)
    \State keep \((\beta,v)\) if \(J\) improves the incumbent
  \EndIf
\EndFor
\State \Return incumbent heading, speed, and the smoothed swath references
\end{algorithmic}
\end{algorithm}

Smoothing of the swath-to-swath join follows Graf Plessen and Bemporad: straight swaths are joined by a curvature-continuous turn whose curvature stays inside \(1/R_{\mathrm{turn}}\) and whose curvature rate stays inside the steering-rate limit \cite{plessen2017}. The speed profile along the join is then reduced, if needed, until \(c(v,a)\ge 0.85\) at every sample of the join. If no speed profile inside \(\mathcal{V}_{\mathrm{meter}}\) meets that test, the outer program lengthens the headland rather than violating the test. Lengthening the headland costs unplanted area at the field edge, which is already the usual price of a turning region \cite{bochtis2008}. The unplanted area is added to the miss term of those edge seeds so that a planner cannot hide a singulation problem by calling it headland.

\subsection{Inner tracker}

The inner tracker follows the smoothed reference with a linear time-varying model-predictive controller. The state and input are
\begin{equation}
\label{eq:xu}
\mathbf{x}_v =
\begin{bmatrix} x & y & \theta & v \end{bmatrix}^{\mathsf{T}},\qquad
\mathbf{u} =
\begin{bmatrix} \delta & a \end{bmatrix}^{\mathsf{T}}.
\end{equation}
The continuous kinematics are Equation~\eqref{eq:bicycle} together with \(\dot{v}=a\). Along a reference \((\mathbf{x}_{v,d},\mathbf{u}_d)\) the discrete linearisation at period \(h=\SI{100}{ms}\) is
\begin{equation}
\label{eq:ltv}
\mathbf{x}_{k+1} - \mathbf{x}_{v,d,k+1} = A_k (\mathbf{x}_k-\mathbf{x}_{v,d,k}) + B_k (\mathbf{u}_k-\mathbf{u}_{d,k}),
\end{equation}
with \(A_k\) and \(B_k\) the standard Jacobian of the bicycle, evaluated on the reference, exactly as in the agricultural linear time-varying trackers \cite{plessen2017}. The horizon is \(N_h=15\) steps, hence \SI{1.5}{s}, which at \SI{1.4}{m/s} looks \SI{2.1}{m} ahead, a little more than one turn radius of preview if the reference curvature is already rising.

The stage cost is not a pure tracking cost. With \(\tilde{\mathbf{x}}=\mathbf{x}-\mathbf{x}_{v,d}\) and \(\tilde{\mathbf{u}}=\mathbf{u}-\mathbf{u}_d\),
\begin{equation}
\label{eq:stage}
\begin{aligned}
\ell_k &= \tilde{\mathbf{x}}_k^{\mathsf{T}} Q \tilde{\mathbf{x}}_k + \tilde{\mathbf{u}}_k^{\mathsf{T}} R \tilde{\mathbf{u}}_k + w_{\rho}\,\rho_{\mathrm{fail}}\bigl(v_k,a_k,\sigma_z(v_k)\bigr) \\
&\quad + w_{M}\,\psi_{\mathrm{res}}(v_k).
\end{aligned}
\end{equation}
The resonance penalty \(\psi_{\mathrm{res}}\) is zero outside \([\SI{2.15}{m/s},\SI{2.55}{m/s}]\) and rises as a quadratic bowl inside it, centred on the speed where \(r=1\). It is a soft version of a constraint the outer program already enforced, repeated in the inner loop so that a disturbance that pushes the speed up cannot cross the bowl without paying. The hard constraints are
\begin{equation}
\label{eq:hard}
\begin{aligned}
&v_k \in \mathcal{V}_{\mathrm{meter}},\qquad
|a_k| \le a_{\max},\qquad
|\delta_k| \le \delta_{\max},\qquad
|\delta_k-\delta_{k-1}| \le \dot{\delta}_{\max}\, h,\\
&c(v_k,a_k) \ge 0.85.
\end{aligned}
\end{equation}
The last of these is linear in \(1/v\) and in \(|a|\) inside each piece of Equation~\eqref{eq:cfrac}, and it is implemented by two linear inequalities on each piece, selecting the piece from the reference speed. The fail term \(\rho_{\mathrm{fail}}\) is convex and increasing in \(|a|\) and, inside the plateau of Figure~\ref{fig:qfi}, nearly flat in \(v\). On the plateau it is approximated by its value at \(v_{\mathrm{work}}\) plus a linear slope taken from the model, which preserves convexity of the quadratic program. Near the boundary of \(\mathcal{V}_{\mathrm{meter}}\) the hard constraint is doing the work and the linear slope is only a tie-break. The quadratic program is solved with qpOASES \cite{ferreau2014}. The linearisation and the constraint matrices are formed with CasADi in the estimation and would be formed by the same generated code on the machine \cite{andersson2019}.

Terminal ingredients follow Mayne, Rawlings, Rao and Scokaert in reduced form \cite{mayne2000}. The terminal cost is the infinite-horizon cost of the linear quadratic regulator computed on the straight-swath linearisation, and the terminal set is a sublevel set of that cost on which the hard constraints are inactive. On a straight swath this is the classical construction. On a turn the reference is already feasible by the smoothing step, and the terminal set is replaced by a terminal equality to the reference at step \(N_h\), which is conservative and cheap. Recursive feasibility is then the statement that a feasible solution shifted by one step, with the terminal law appended, remains feasible. It holds while the disturbance in \(v\), from slip on the grade, stays inside the margin between \(v_{\mathrm{work}}\) and the boundary of \(\mathcal{V}_{\mathrm{meter}}\). The outer program is required to leave that margin. At the design point \(v_{\mathrm{work}}=\SI{1.4}{m/s}\) and the latency boundary near \SI{1.87}{m/s}, the margin is about \SI{0.47}{m/s}, which is larger than the speed error a \num{8}\% grade produces under the rolling-resistance model. If a steeper patch consumes the margin, the inner law brakes. Braking is always feasible inside Equation~\eqref{eq:hard} because lower speed improves \(c(v,a)\) and moves away from the resonance.

\begin{proposition}
\label{prop:feas}
If the smoothed reference lies in the interior of Equation~\eqref{eq:hard} and the one-step slip disturbance is bounded by the speed margin, the shifted warm-start remains feasible and the vehicle speed never settles inside the resonant interval.
\end{proposition}

The proof is the standard shifted-sequence argument of \cite{mayne2000} applied to a constraint set that does not depend on the measured singulation. The resonant interval is excluded by construction of \(\mathcal{V}_{\mathrm{meter}}\) and by \(\psi_{\mathrm{res}}\), not by an on-line estimate of seed spacing. That choice is deliberate. An on-line estimate of spacing arrives one flight time late, about \SI{0.29}{s} after release, and is the wrong signal to use as a hard constraint that must be feasible at the current step. The on-line spacing estimate is logged for the seed map \cite{griepentrog2005,ehsani2004}. It is not placed inside Equation~\eqref{eq:hard}.

\begin{figure}[ht]
\centering
\begin{tikzpicture}[font=\small, >=Stealth]
  \draw[thick] (0,0) -- (10,0);
  \draw[blue!70!black, thick] plot[smooth] coordinates {(0,0.1) (2,0.15) (4,0.55) (6,1.05) (8,1.15) (10,1.2)};
  \draw[dashed] (0,0) -- (10,0.15);
  \foreach \x in {0,1.3,2.6,3.9,5.2} {
    \draw[orange!80!black, ->] (\x,0) -- (\x+1.15,0.22);
  }
  \node[below] at (5,-0.15) {samples of the \SI{1.5}{s} horizon};
  \node[above] at (7.2,1.25) {reference};
  \node[anchor=west] at (0.2,1.5) {predicted speed stays off the resonant bowl};
\end{tikzpicture}
\caption{Inner horizon. Tracking error is penalised, and so is any predicted speed that approaches the timber resonance or any predicted acceleration that would make a repair late.}
\label{fig:horizon}
\end{figure}

\subsection{What the inner cost will not do}

Xu, Lai, Guo and Lu showed that a tracker allowed to care about efficiency can beat a tracker that only cares about path error, by \num{8.56}\% in their simulation \cite{xu2022}. The same freedom can be misused. If \(w_{\rho}\) is set to zero, Equation~\eqref{eq:stage} becomes a conventional tracker plus a weak resonance bowl, and a high \(w_T\) in the outer program will push \(v_{\mathrm{work}}\) to the boundary of the constraint. The design therefore keeps the index limits in the outer filter even when the weights change, and it keeps \(c(v,a)\ge 0.85\) as a hard inner constraint even when \(w_{\rho}\) is small. Efficiency is allowed to choose among admissible speeds. It is not allowed to define admissibility.

The inner program also does not schedule plate speed. Plate speed remains Equation~\eqref{eq:plate}, driven by the speed the vehicle actually has, not by the speed the planner wanted a second ago. Using the realised \(v\) closes the interface in the safe direction. If slip on the grade drops the realised speed, the plate slows and the spacing is held. If the planner's \(v_{\mathrm{work}}\) were sent straight to the plate while the wheels were still slow, the meter would bunch seeds at the start of every swath. Mangus's acceleration result is the field version of that bunching \cite{mangus2017}. Realised speed as the feedforward input is how the first technology protects itself from a confident plan.

Variable rate along a swath, in the precision-agriculture sense of a prescription map, fits the same interface. A prescription \(x_{\mathrm{ref}}(s)\) changes Equation~\eqref{eq:rate} and changes \(\Delta t\), and therefore changes \(\mathcal{V}_{\mathrm{meter}}\) locally, because a shorter spacing at the same speed is a shorter hole period. The outer program, if a prescription is present, evaluates \(c\) with the local spacing and lowers \(v_{\mathrm{work}}\) on the dense parts of the map. The estimation field uses a constant spacing, so this branch is specified and not exercised. It is the reason \(x_{\mathrm{ref}}\) is an input of Definition~\ref{def:p1} rather than a constant buried in the plate law.

\subsection{How the two technologies avoid a circular design}

There is an apparent circle. The metering technology publishes \(\mathcal{V}_{\mathrm{meter}}\) using \(\tau_v\), \(\tau_a\), \(f_n\) and \(\lambda\). The coverage technology chooses \(v\) inside that set and then drives the flexure at that \(v\), which is what \(f_n\) and \(\lambda\) were worried about. The circle is broken by the order of computation. The flexure parameters and the latency are properties of the hardware and of the detector. They do not depend on the heading. Equation~\eqref{eq:vset} can be evaluated before any coverage search. The coverage search never writes back into \(k\), \(\zeta\), \(\tau_v\) or \(\tau_a\). The only write-back is the realised speed, which the plate law accepts as an exogenous input. A fixed point is not being solved. A constraint computed from hardware is being obeyed by a planner.

The write-back that does exist is diagnostic. After a pass, the logged release times and the logged \(\hat{\delta}\) can update the estimate of \(f_n\) and of \(\lambda\), if the roughness of that field is not the design wavelength. The updated parameters change \(\mathcal{V}_{\mathrm{meter}}\) for the next pass. They are not allowed to change it in the middle of a hole period. This is the same discipline as the dropped-frame rule: a late adaptation is a worse fault than a slightly conservative speed.

\subsection{Arithmetic of the work model at two headings}

The outer program is a search over a closed form, so the closed form should be visible on one heading before Chapter~\ref{ch:estimate} prints a table of many. Take \(v=\SI{1.5}{m/s}\), grade \num{8}\% so \(\alpha=\arctan 0.08\) and \(\sin\alpha=\num{0.0797}\), \(\sin^2\alpha=\num{0.00636}\), and \(Mg=460\cdot 9.81=\SI{4513}{N}\). Rolling resistance is
\begin{equation}
\label{eq:frr}
C_{rr} Mg\cos\alpha = 0.085\cdot 4513\cdot 0.9968 = \SI{382}{N}.
\end{equation}
The slip term at a heading uses \(\Gamma(\beta)=1+1.15\cos^2\beta\). Uphill, \(\beta=0\), \(\Gamma=2.15\), and
\begin{equation}
\label{eq:fslip}
C_s Mg\sin^2\alpha\;\Gamma = 2.0\cdot 4513\cdot 0.00636\cdot 2.15 = \SI{123}{N}.
\end{equation}
At the contour, \(\beta=\ang{90}\), \(\Gamma=1\) and the slip term falls to \SI{57}{N}. Aerodynamic drag at \SI{1.5}{m/s} with \(C_dA=\SI{0.85}{m^2}\) and \(\rho=\SI{1.15}{kg/m^3}\) is \(\tfrac12\rho C_dA v^2\approx\SI{1.1}{N}\), which is why work barely moves when speed moves inside the plateau. Dividing by drivetrain efficiency \(\eta_d=0.70\) gives a force-per-metre, in energy units, of about \SI{724}{J/m} uphill and about \SI{630}{J/m} on the contour. Those are the values the estimation uses. They are not a dynamometer record.

On the rectangle, an uphill heading runs rows of \SI{180}{m} and needs \(n=\lceil 80/2.4\rceil=34\) swaths. Planting distance is \(34\cdot 180=\SI{6120}{m}\). Turn distance is \(34\cdot\pi\cdot 4.2\approx\SI{448}{m}\). Equation~\eqref{eq:work} then charges planting at the uphill force and turns at the lower turn speed with the eight-percent turn factor. The resulting work is \SI{4.78}{MJ}, which is the uphill entry the estimation table will show at this speed. A heading \ang{55} from uphill shortens the grade factor, increases the number of swaths, and lands near \SI{4.42}{MJ}. The difference, about \SI{0.36}{MJ}, is the whole energy prize available on this field. It is smaller than the energy a reader might expect from the phrase ``terrain-aware'', and it is large enough, against a time weight of \SI{0.55}{MJ} per hour, that the optimiser will not ignore it. It is not large enough to justify leaving \(\mathcal{V}_{\mathrm{meter}}\). That comparison is the reason the index filter sits in front of the scalar cost rather than inside it.

The accelerated fraction on the long uphill rows at \SI{1.5}{m/s}, from Equation~\eqref{eq:frac} with \(a_{\mathrm{turn}}=\SI{0.7}{m/s^2}\), is
\begin{equation}
\label{eq:phi}
\phi_{\mathrm{acc}} = \frac{1.5^2}{0.7\cdot 180} = 0.018,
\end{equation}
under two percent of the row. On the short contour rows of \SI{80}{m} the same formula gives \(1.5^2/(0.7\cdot 80)=0.040\). Mixing the zero-acceleration and the turn-acceleration miss models with those weights, as Equation~\eqref{eq:mix} does, changes the closed-loop miss probability in the third digit at this speed. Heading is therefore not a singulation decision at \SI{1.5}{m/s}. It becomes one if someone proposes \(v=\SI{2.4}{m/s}\), where both the base miss model and the magnification are already bad on the straight, and the accelerated fraction makes them worse. The outer program never reaches that comparison as a cost trade, because the index filter has already discarded the candidate.

\subsection{Linearisation used by the inner tracker}

On a straight swath at \(v_d=\SI{1.4}{m/s}\), \(\theta_d=0\) and \(\delta_d=0\), the bicycle Jacobian with respect to \((x,y,\theta,v)\) and \((\delta, a)\), sampled at \(h=\SI{0.1}{s}\), is dominated by a simple structure. Position \(x\) integrates speed. Position \(y\) integrates \(v\theta\) for small heading error, so the discrete \(A\) carries a term \(h v_d\) from \(\theta\) into \(y\). Heading integrates \((v/\ell)\delta\) for small steering, so \(B\) carries \(h v_d/\ell\) from steering into heading, and from there, one step later, into lateral error. With a wheelbase \(\ell=\SI{1.6}{m}\),
\begin{equation}
\label{eq:jac}
\frac{\partial \dot{\theta}}{\partial \delta}\bigg|_{\delta=0} = \frac{v}{\ell} = \frac{1.4}{1.6} = \SI{0.875}{s^{-1}},
\end{equation}
and one tracker step of steering equal to \SI{0.05}{rad} changes heading by about \(0.875\cdot 0.05\cdot 0.1=\SI{0.0044}{rad}\), then changes lateral position on the following step by about \(1.4\cdot 0.0044\cdot 0.1\approx\SI{0.6}{mm}\). The tracker does not steer in one-step impulses of that size. The arithmetic shows why a \SI{1.5}{s} horizon is long enough to see a lateral error coming and short enough that the linearisation, recomputed each step on the reference, does not have to carry the vehicle around a full turn inside one linear model. The turn itself is in the reference, through the smoothing of Graf Plessen and Bemporad, and the linearisation is about that curved reference rather than about a straight line \cite{plessen2017}. Recomputing \(A_k\) and \(B_k\) every \SI{100}{ms} is the ``time-varying'' part of the linear time-varying controller. Freezing them for the whole headland would be a linear controller about the wrong point.

The quadratic weights used with this linearisation are diagonal,
\begin{equation}
\label{eq:QR}
\begin{aligned}
Q &= \diag(2,\, 8,\, 4,\, 1),\\
R &= \diag(0.4,\, 0.3),
\end{aligned}
\end{equation}
on states in metres, radians and metres per second, and on inputs in radians and metres per second squared. Lateral error is weighted more heavily than longitudinal error because a row that has drifted by ten centimetres plants the neighbouring row's strip, whereas a longitudinal lag of ten centimetres is removed by the plate law as soon as the realised speed is known. Speed error is weighted lightly inside \(Q\) because speed is already constrained by Equation~\eqref{eq:hard}, and a large speed weight would fight the resonance bowl and the repair constraint with an objective they have already settled. These weights are part of the specification of the inner tracker. They are not identified from a human driver.

A warm start for qpOASES is the previous solution with the first move discarded and the terminal linear-quadratic move appended \cite{ferreau2014,mayne2000}. On a straight swath the active set is usually empty: the inequalities are slack by the margin the outer program left, about \SI{0.47}{m/s} between \SI{1.4}{m/s} and the latency boundary near \SI{1.87}{m/s}. The solver then returns in a few factorisations. The \SI{40}{ms} budget in the build is set for the worse case, a headland in which the acceleration inequality and the repair inequality are both close to active. If a solve exceeds the budget, the tracker applies the second move of the previous feasible solution and marks the step as a hold. A hold is feasible by the shifted-sequence argument of Proposition~\ref{prop:feas}. An infeasible scramble for a new solution is not.

\subsection{A local spacing change, and why the speed set moves with it}

A prescription that asks for \(x_{\mathrm{ref}}=\SI{0.12}{m}\) on one patch of the field, and \SI{0.20}{m} elsewhere, changes the hole period at the same ground speed from \SI{0.143}{s} to \(0.12/1.4=\SI{0.086}{s}\). Against the repair threshold of \SI{0.107}{s} from Equation~\eqref{eq:cnum}, the tighter spacing is no longer inside the full-repair set at \SI{1.4}{m/s}. The speed that restores \(\Delta t\ge\SI{0.107}{s}\) is
\begin{equation}
\label{eq:vpresc}
v \le \frac{0.12}{0.107} = \SI{1.12}{m/s}.
\end{equation}
The outer program, on that patch, lowers the reference speed to \SI{1.1}{m/s} before the tracker runs. The plate law does not have to discover the violation by attempting repairs and then withdrawing them. This is the local form of Requirement~R10. One speed set, recomputed from the local spacing, is the only message the prescription is allowed to send into the tracker. A prescription that instead sent a plate speed directly would bypass the repair test and recreate, on the dense patch, the late-repair regime the whole design is built to avoid.

The same patch does not require a new heading. Heading was chosen for the field's grade and for the number of turns, neither of which the prescription changes. Replanning the heading because one patch wants a slower speed would rotate every other swath for a constraint that a speed change already satisfies. The inner tracker is allowed to deviate from \(v_{\mathrm{work}}\) inside the horizon. The outer heading stays put until the elevation model or the field boundary changes. That is a second time-scale split, inside the coverage technology, and it is why Algorithm~\ref{alg:outer} returns a single heading for the rectangle even though the tracker is free to slow down.

\subsection{What has been fixed for the estimation}

The estimation is allowed to evaluate, and is not allowed to revise, the following objects: the flexure parameters and the latency that define \(\mathcal{V}_{\mathrm{meter}}\); the miss and depth maps of Equations~\eqref{eq:openmiss}--\eqref{eq:deptherr}; the work coefficients inside Equation~\eqref{eq:fpm}; the weights \(w_T\) and the failure weight; the index thresholds \num{0.955} and \SI{4}{mm}; the resonant interval \([\SI{2.15}{m/s},\SI{2.55}{m/s}]\); and the tracker weights of Equation~\eqref{eq:QR}. Chapter~\ref{ch:estimate} may only choose which plans to place side by side and may only report the numbers those objects produce. The hardware chapter that precedes the tables exists so that a reader who doubts a latency or a sample period can see the task that is supposed to meet it, and can see which table entry moves if that task is slower than its budget.

\subsection{One headland, sample by sample}

The inner tracker is easiest to distrust in a turn, because that is where speed, acceleration and curvature change together and where Mangus and colleagues found the singulation errors organised themselves \cite{mangus2017}. Consider the exit from an uphill swath into a turn of radius \(R_{\mathrm{turn}}=\SI{4.2}{m}\) at the moment the reference begins to slow from \SI{1.4}{m/s} toward the turn speed \SI{0.8}{m/s}. The distance available for that slowing, if the reference is to arrive at the turn at \SI{0.8}{m/s} with \(a=-\SI{0.7}{m/s^2}\), is
\begin{equation}
\label{eq:brake}
\frac{v_1^2-v_2^2}{2|a|} = \frac{1.4^2-0.8^2}{1.4} = \SI{0.94}{m}.
\end{equation}
At a \SI{100}{ms} sample the vehicle covers about \SI{0.14}{m} at the start of the braking and about \SI{0.08}{m} at the end, so the manoeuvre occupies roughly eight to ten tracker steps, well inside the fifteen-step horizon. The horizon that begins a second and a half before the brake therefore already contains the whole speed change. The predicted accelerations are about \(-\SI{0.7}{m/s^2}\). Equation~\eqref{eq:cfrac} at \(v=\SI{1.4}{m/s}\) and \(|a|=\SI{0.7}{m/s^2}\) still returns \(c=1\), because the hole period is \SI{0.143}{s} and the acceleration surcharge \(0.012\cdot 0.7\approx\SI{0.008}{s}\) does not push the threshold above the hole period. The hard constraint is slack. The stage cost still penalises the acceleration through \(\rho_{\mathrm{fail}}\), but only by the mixture weight of a short distance, which is why \(\phi_{\mathrm{acc}}\) on the long rows stayed near two percent.

What the tracker must not do, and what a pure time-optimal tracker would do, is postpone the braking and then brake harder. A brake of \(\SI{1.2}{m/s^2}\) over a shorter distance would finish the swath earlier by a fraction of a second and would drop \(c\) toward the \num{0.71} computed earlier for that acceleration at a slightly higher speed. Equation~\eqref{eq:hard} forbids that brake even if the time weight would have liked it. The headland is therefore not an exception to the metering technology. It is the place the exception would have been, and the constraint is written there on purpose.

During the turn itself the plate may be inhibited by the boot-lift interlock if the headland is unplanted, which is the usual layout \cite{bochtis2008}. Inhibition is cleaner than asking the meter to singulate while the boot is out of the soil and the ultrasonic range is invalid. When the next swath begins, the integral of the depth loop is released from its freeze, the plate feedforward tracks the realised speed as the vehicle climbs back from \SI{0.8}{m/s} to \SI{1.4}{m/s}, and repairs are authorised only once \(c\) has returned. The first metre of the new swath is the other place a naive plan bunches seeds, by commanding the working plate speed before the wheels are at the working speed. Using realised \(v\) in Equation~\eqref{eq:plate} spreads that first metre at the correct spacing instead. The cost is a slightly later start of full-rate planting, already paid for in the headland length.

\subsection{Logging that makes the next pass stricter rather than braver}

Each release writes a row to the seed map: vehicle pose, filter release time, estimated exit velocity, detector label, repair flag, realised delay, and filtered depth. The row is an estimate, not a proof that a seed lies at the computed coordinate. Its first use is the one Griepentrog, N{\o}rremark, Nielsen and Blackmore described, a map a later operation can follow \cite{griepentrog2005}. Its second use is internal, and it is deliberately asymmetric. If the realised delay over a swath exceeds the budget, the next swath's \(\tau\) is raised and \(\mathcal{V}_{\mathrm{meter}}\) shrinks. If the delay is better than the budget, \(\tau\) is not lowered on the same day. A fast swath can be luck of the scheduler. A slow swath is evidence the inequality is tighter than the specification. The same asymmetry applies to \(f_n\). A resonance that appears in \(\hat{\delta}\) widens the forbidden interval for the next pass. A quiet swath does not delete the interval, because the roughness wavelength of one swath is not the wavelength of the field. Adaptation that only ever tightens the constraint cannot talk a machine into the regime Borja's loop entered when the processor fell behind \cite{borja2021}. Adaptation that updates in both directions can.

The log is also the object a later germination count has to be joined to, if the survival function \(S(\sigma_z)\) is ever replaced by a measured curve. Joining requires the landing coordinate from Equation~\eqref{eq:carry}, not only the vehicle coordinate at the moment the detector fired. Those two coordinates differ by the flight transport, about \SI{0.41}{m} at the design speed, which is two target spacings. A survival count joined to the wrong coordinate would teach the next version of the cost to fear the wrong places. The build therefore stores both stamps. None of this changes the quadratic program. It is the reason the release log is a row in the interface table and not an optional debug print.

\subsection{A calibration pass before the first planted swath}

The stiffness \(k\) is a specification of the timber link. The soil coefficient \(k_{\mathrm{soil}}\) in Equation~\eqref{eq:flinear} is not, because it belongs to the field. The build therefore begins a planting day with a calibration pass rather than with a planted swath. The pass is one swath length, plates disengaged, boot in the soil, depth screw held at three set-points around the \SI{18}{mm} target, for a few metres each. Boot force and ultrasonic depth are logged. The slope of force against depth is \(k_{\mathrm{soil}}\). A value of order a few newtons per millimetre is what the linearisation expects; the \SI{80}{N} static force at the target depth is the intercept \(F_0(z^{\star})\), and it is the number that produced the \SI{9.8}{mm} undeflected offset \(F_0/k\). If the calibration intercept is \SI{110}{N} rather than \SI{80}{N}, the undeflected offset is \(110/8200=\SI{13.4}{mm}\), still inside the screw travel, and the integral term is required to work harder. The filter does not need a new structure for that. It needs the new intercept written into \(F_{\mathrm{soil}}\) before the first repair is authorised.

The same pass estimates the roughness wavelength that the design took as \(\lambda=\SI{0.45}{m}\). The force residual after the linear depth term is removed is a time series along the swath. Its dominant spatial frequency, if one is visible, replaces \(\lambda\) in the magnification formula before \(\mathcal{V}_{\mathrm{meter}}\) is published. If the residual is broadband, the design wavelength is kept and the process noise on \(\delta\) stays on the magnification schedule already stated. Publishing \(\mathcal{V}_{\mathrm{meter}}\) after this pass, rather than from the workshop specification alone, is the one adaptation allowed before motion. It is still not an adaptation inside a hole period. The tracker does not move until the set has been published, which is the order of computation that broke the apparent circle between the two technologies.

A calibration that cannot be finished, because the ultrasonic return is invalid or the force signal is saturated, inhibits planting. The alternative is to run the design values and hope the field matches the specification. The interlock refuses that alternative. A wooden machine is light enough that an unexpected soil force moves the boot by millimetres, which is the whole reason the depth loop exists, and a depth loop without a force intercept is an integral winding up against an unknown load. The pass costs one swath of time, already small beside the \SI{1.4}{h} the estimation assigns to the field, and it is the price of keeping Equation~\eqref{eq:deptherr} attached to the soil under the machine rather than to the soil imagined in the problem set.

Table~\ref{tab:cal} is the decision the pass is allowed to make. The screw travel specified in the build has to cover the largest intercept in the table, \SI{24}{mm} of undeflected offset plus the \SI{18}{mm} target plus a lift clearance. An intercept of \SI{200}{N} is a stiffer or a deeper soil than the problem set assumed, and it is still a calibration result rather than a reason to abandon the flexure model. Only a saturated load cell, which cannot produce an intercept at all, inhibits the day. The resonant speed column shows the other output of the same pass. A shorter roughness wavelength moves the forbidden steady speed down, and the tracker inherits that number with the feasible set. No weight in the coverage cost is allowed to override the last column.

\begin{table}[ht]
\centering
\caption{Calibration outcomes the first pass may publish. Offsets use \(k=\SI{8200}{N/m}\). The resonant speed uses the measured wavelength in \(v \approx \lambda f_n\).}
\label{tab:cal}
\begin{tabular}{@{}rrll@{}}
\toprule
Intercept (N) & Offset (mm) & Measured \(\lambda\) (m) & Forbidden near (m/s) \\
\midrule
80 & 9.8 & 0.45 & 2.4 \\
110 & 13.4 & 0.45 & 2.4 \\
110 & 13.4 & 0.30 & 1.7 \\
200 & 24.4 & 0.60 & 3.3 \\
\bottomrule
\end{tabular}
\end{table}

The third row is the one that changes the plan rather than only the depth bias. A wavelength of \SI{0.30}{m} puts the resonance near \(0.30\cdot 5.53\approx\SI{1.66}{m/s}\), inside the plateau where the latency constraint would still have allowed \SI{1.4}{m/s} to \SI{1.8}{m/s}. The feasible set then has a hole where the magnification exceeds the bound used in Equation~\eqref{eq:vset}. The low side of that hole keeps the arm quiet and spends time. The high side is already where repair authority falls and where the \SI{3}{Hz} depth servo loses the roughness. On the estimation coefficients the index filter refuses the high side, and the machine plants more slowly than the \SI{1.4}{m/s} of Table~\ref{tab:plans} after a calibration that discovers short roughness. That is the correct refusal. A planner that kept \SI{1.4}{m/s} because the workshop wavelength had been \SI{0.45}{m} would be driving on the resonance the pass had just measured.

Chapter~\ref{ch:estimate} can now be a calculation rather than a narrative. The hardware that would host the calculation is specified first, so that the sample periods and the latency budget used above have a place to live. The calculation then evaluates Equations~\eqref{eq:openmiss}--\eqref{eq:work} on the field of the problem set, for the open-loop baselines and for the plan Algorithm~\ref{alg:outer} returns.

\chapter{System Build and Estimation}
\label{ch:estimate}

The two technologies are specified as equations and as constraints. This chapter turns the specification into a build that could host them, and then evaluates the equations on the field fixed in the problem set. Section~\ref{sec:build} is the system build: mechanical arrangement, sensors, computers, clocks and the software tasks that correspond to the blocks of Chapter~\ref{ch:design}. Section~\ref{sec:est} is the estimation: the protocol, the baselines, the numbers produced by the models, and the limits of what those numbers mean. The estimation does not contain a measurement that was not already implied by Table~\ref{tab:spec} and by Equations~\eqref{eq:openmiss}--\eqref{eq:work}. Where a sentence would be easier to read if it said ``the machine achieved'', the sentence says ``the model assigns''.

\section{System build}
\label{sec:build}

\subsection{Mechanical arrangement}

The chassis is a ladder of laminated softwood, two longitudinal members and three cross members, sized so that a four-row toolbar can be bolted on without a welded steel subframe as the primary structure. The Wood Handbook values for structural softwood are the design data for the members: modulus parallel to grain in the range used to justify a timber machine, and a check that moisture content is controlled before lamination so that the modulus assumed in \(k\) is the modulus the arm still has after a season \cite{woodhandbook}. Each row unit is a parallelogram. The two long links are timber. The pivots are steel pins in replaceable bushings, because a timber-on-timber pivot would frett and would change \(z_{\mathrm{bias}}\) faster than the filter's bias state is meant to follow. The boot is a short steel or cast opener on the timber link. Putting the wear surface in steel does not make the machine a steel planter. The compliance that moves the seed gate is still the timber link, and \(k=\SI{8200}{N/m}\) is the specification of that link, not of the boot.

The hopper on each row holds a single woody species, consistent with Assumption~\ref{as:seed}. It is a timber box with a polymer liner. The liner, not the wood, is the surface the seed slides on, so that pitch and dust do not change the friction the plate was tuned against. The plate has \(N=24\) cells, cut for the design seed, and is driven by a stepper through a reduction that can advance one cell pitch in the \(\tau_a=\SI{35}{ms}\) already charged to the repair law. If a vacuum plate is used instead of a mechanical cell, the same stepper drives the plate and a small blower holds the nominal \(\Delta p\) of Equation~\eqref{eq:vacuum}, including the vibration margin at the highest magnification allowed in \(\mathcal{V}_{\mathrm{meter}}\). The seed tube below the gate is wood-lined, of vertical drop \(H=\SI{0.42}{m}\), with a window for the lower camera view. A wooden tube is part of the object the detector must tolerate: the background is brown, and the augmentation practice of Blaser et al., who trained across light and viewpoint, is a requirement of the training set rather than an option \cite{blaser2024}.

The depth actuator is a linear screw on each row, rated for the \SI{120}{N} saturation used in the Lyapunov argument, with the \SI{80}{ms} lag of the model. It acts on the parallelogram and does not try to bend the timber link faster than the \SI{3}{Hz} bandwidth. Travel is set so that the \SI{18}{mm} target sits in the middle of the stroke, with enough extra stroke to remove the \SI{9.8}{mm} static deflection and to lift the boot clear of the soil in transport. Lift is a separate command. When the boot is up, the integral in Equation~\eqref{eq:depth} is frozen, as specified.

Row spacing is \SI{0.60}{m}, so four rows occupy the \SI{2.4}{m} swath used in the coverage arithmetic. The machine is drawn or self-propelled. The estimation uses a mass of \SI{460}{kg} and a wheelbase \(\ell\) inside the bicycle model, and it does not depend on which of those two traction arrangements supplies the force, because \(W_{\mathrm{mech}}\) is computed from resistance and slip rather than from an engine map. A self-propelled arrangement makes the speed command of the inner tracker a motor torque. A drawn arrangement makes it a request to the tractor. ISO~11783 is the request path in the drawn case \cite{iso11783}. In both cases the realised ground speed, not the request, is what Equation~\eqref{eq:plate} uses.

\begin{figure}[ht]
\centering
\begin{tikzpicture}[font=\small, >=Stealth, node distance=0.35cm]
  \node[draw, minimum width=10.2cm, minimum height=1.05cm, fill=brown!20] (frame) {laminated chassis, \SI{2.4}{m} toolbar};
  \node[draw, fill=brown!10, below left=0.6cm and -0.2cm of frame, minimum width=2.2cm, minimum height=0.9cm] (r1) {row 1};
  \node[draw, fill=brown!10, right=0.25cm of r1, minimum width=2.2cm, minimum height=0.9cm] (r2) {row 2};
  \node[draw, fill=brown!10, right=0.25cm of r2, minimum width=2.2cm, minimum height=0.9cm] (r3) {row 3};
  \node[draw, fill=brown!10, right=0.25cm of r3, minimum width=2.2cm, minimum height=0.9cm] (r4) {row 4};
  \node[draw, fill=blue!8, below=0.55cm of r2, minimum width=4.6cm, minimum height=0.85cm] (edge) {embedded computer, one detector thread per row};
  \node[draw, fill=orange!12, below=0.55cm of r3, minimum width=4.6cm, minimum height=0.85cm] (trk) {tracker computer, SCT-MPC at \SI{10}{Hz}};
  \draw[->] (r1.south) -- ++(0,-0.25) -| (edge.north);
  \draw[->] (r2.south) -- (edge.north);
  \draw[->] (r3.south) -- (trk.north);
  \draw[->] (r4.south) -- ++(0,-0.25) -| (trk.north);
  \draw[->] (edge) -- node[below, font=\scriptsize] {feasible \(v\)} (trk);
\end{tikzpicture}
\caption{Four-row build. Metering tasks stay on the row computers' detector threads. The tracker sees a feasible speed set and a realised speed, not a private copy of the plate state.}
\label{fig:build}
\end{figure}

\subsection{Sensors and the time base}

Each row carries one camera, one boot load cell, one ultrasonic range sensor and one inertial sensor. The camera frame is the observation the detector consumes. The load cell measures \(F_b\). The ultrasonic sensor measures \(r_u\). The inertial sensor measures the acceleration and the pitch rate that observe \(\ddot{\delta}\) and, by a known lever, \(\dot{\delta}\). A vehicle-level satellite and inertial navigation unit supplies \(x\), \(y\), \(\theta\) and \(v\) to the tracker. Centimetre-level satellite positioning is the class of fix Ehsani, Upadhyaya and Mattson used to map seed locations, and it is sufficient for a swath tracker at this speed \cite{ehsani2004}. It is not sufficient, by itself, to decide whether a cell was empty. That decision stays on the row camera.

All of these samples are useless to the filter if they do not share a clock. The build uses one clock domain. The tracker computer is the time master. Each row board disciplines its sampling to that clock over the same network that carries the feasible-speed message. A camera frame is stamped at exposure start, not at arrival of the image at the detector, because \(\tau_v\) is defined to include exposure. A plate command is stamped at the moment it is issued to the stepper. The difference between the exposure stamp of the frame that saw an empty cell and the issue stamp of the repair is the realised \(\tau_v+\tau_a\), and it is logged. If the realised delay exceeds the value assumed in Equation~\eqref{eq:vset} for more than a short burst, the row withdraws repair authority by lowering \(c\), which tightens \(\mathcal{V}_{\mathrm{meter}}\) from the machine rather than from the specification. This is the on-line honesty the Borja reversal requires: the inequality is checked with the delay the hardware is actually producing \cite{borja2018,borja2021}.

Table~\ref{tab:budget} is the timing budget the computers are required to meet. The numbers are allocations, and they sum to the latencies the control laws already use. A build that cannot meet an allocation has to change \(\mathcal{V}_{\mathrm{meter}}\) before it changes the law.

\begin{table}[ht]
\centering
\caption{Timing budget allocated to one row and to the vehicle tracker.}
\label{tab:budget}
\begin{tabular}{@{}lrr@{}}
\toprule
Task & Period & Budget \\
\midrule
Exposure and readout & \SI{33}{ms} frame & \SI{8}{ms} of \(\tau_v\) \\
Detector inference & per frame & \SI{14}{ms} of \(\tau_v\) \\
Filter update and plate law & \SI{10}{ms} & inside the remaining latency \\
Stepper execution of one pitch & on repair & \(\tau_a=\SI{35}{ms}\) \\
Depth servo & \SI{20}{ms} & bandwidth \SI{3}{Hz} \\
Feasible-set publish & \SI{100}{ms} & to the tracker \\
Tracker quadratic program & \SI{100}{ms} & finish inside \SI{40}{ms} \\
Heading replan & on request & seconds, off the hole period \\
\bottomrule
\end{tabular}
\end{table}

The detector budget of \SI{14}{ms} is plausible for a small model of the YOLOv10 class on an embedded graphics device, and it is not plausible for the class of microcontroller Borja used \cite{wang2024yolov10,borja2021}. The build therefore specifies an embedded graphics module per two rows, or one module with four threads and a hard preemption so that a slow frame on row~2 cannot borrow the \SI{14}{ms} of row~3. A frame that is still running at the deadline is cancelled. Cancellation is the dropped-frame rule of Section~\ref{sec:tech1}.

\subsection{Software tasks}

The software is a set of tasks with the periods in Table~\ref{tab:budget}, not a single loop. Bechar and Vigneault's subsystem split---sensing, decision, action, navigation---maps onto these tasks directly \cite{bechar2016}. Sensing is the camera and the analog acquisition. Decision on the row is Algorithm~\ref{alg:plate} and Equation~\eqref{eq:depth}. Action on the row is the stepper and the depth screw. Navigation is Algorithm~\ref{alg:outer} together with the quadratic program of Equation~\eqref{eq:stage}. The farm-scale record is a seed map: for each release, the filter's \(\hat{t}_{\mathrm{rel}}\), the vehicle pose at that time, and the estimated landing from Equation~\eqref{eq:carry}. Griepentrog and colleagues argued for exactly this record as the input to later weeding \cite{griepentrog2005}. A restoration planting uses it as the input to a later survival count, which is also the measurement that would one day replace the estimation function \(S(\sigma_z)\).

Task priorities follow the hole period. On each row the plate command outranks the depth servo, and the depth servo outranks logging. On the vehicle the quadratic program outranks a replan of heading. A replan that is still running when the machine is moving is discarded rather than applied half-computed, because a half-applied heading is not a heading the outer program evaluated. The tracker continues on the last feasible reference. This priority is the software form of Proposition~\ref{prop:feas}: feasibility of the current horizon is more important than optimality of the next heading.

The quadratic program is warm-started from the shifted previous solution, as the feasibility argument requires. qpOASES is the solver because it is an active-set method with a bounded, predictable cost on a problem of this size: fifteen steps, four states, two inputs, and linear inequalities \cite{ferreau2014}. A general nonlinear solver would also solve the problem and would not meet Requirement~R7 with the same predictability. CasADi is used off the machine to generate the Jacobian of the bicycle and the constraint matrices, and the generated evaluation is what runs at \SI{10}{Hz} \cite{andersson2019}. The estimation in the next section uses the same matrices in a direct evaluation of the outer cost. It does not need to simulate every quadratic-program step in order to rank headings, because Proposition~\ref{prop:feas} says that a reference inside the speed set is tracked without leaving the set. The estimation therefore ranks references by Equations~\eqref{eq:work}--\eqref{eq:mix}, and it treats tracking error as a separate, smaller contribution that does not reorder the headings.

\begin{figure}[ht]
\centering
\begin{tikzpicture}[x=0.78cm, font=\scriptsize, >=Stealth]
  \draw[->] (0,0) -- (14.2,0) node[below] {time};
  \foreach \x in {0,2,4,6,8,10,12} \draw (\x,0.05) -- (\x,-0.05) node[below] {\x0 ms};
  \draw[blue!70!black, line width=3pt] (0,1.0) -- (1.4,1.0) node[midway, above=2pt] {infer};
  \draw[blue!70!black, line width=3pt] (3.3,1.0) -- (4.7,1.0);
  \draw[blue!70!black, line width=3pt] (6.6,1.0) -- (8.0,1.0);
  \draw[orange!80!black, line width=3pt] (0.2,0.45) -- (1.0,0.45) node[midway, above=1pt] {QP};
  \draw[orange!80!black, line width=3pt] (10.2,0.45) -- (11.0,0.45);
  \node[anchor=west] at (0,1.55) {row detector, \SI{30}{Hz}};
  \node[anchor=west] at (0,-0.15) {};
  \node[anchor=east] at (0,0.45) {tracker};
\end{tikzpicture}
\caption{Two clocks. The row detector is fast enough to see every cell several times at the design speed. The tracker is an order of magnitude slower, which is the separation the joint problem was split along.}
\label{fig:clocks}
\end{figure}

\subsection{Safety interlocks}

A planter has a small number of states in which the correct command is to stop metering, and the build treats them as interlocks rather than as expensive regions of the cost. The plate is inhibited when the boot is lifted, when realised \(\tau_v+\tau_a\) exceeds the budget for longer than three hole periods, when the hopper-level sensor declares a row empty, and when the tracker reports that no speed inside \(\mathcal{V}_{\mathrm{meter}}\) is feasible on the current slope patch. The last of these is the situation Proposition~\ref{prop:feas} excludes on the design grade and includes on a grade or a roughness the specification did not see. Inhibiting the plate loses the rest of the swath and saves a resonant or a late-repair passage, which on the model of Figure~\ref{fig:qfi} is the right side of the cliff to stand on.

The depth screw is inhibited when the ultrasonic return is invalid for more than \SI{200}{ms}, and the last safe force is held. A held force is not a depth guarantee. It is a refusal to integrate noise. The vehicle is inhibited from entering the resonant speed interval even if a human operator overrides the working-speed set-point, unless a maintenance mode is selected with the plates disengaged. Override of a singulation constraint by a foot throttle is how a correct planner is turned back into the open-loop baseline. The interlock is unpopular and it is part of the technology.

\subsection{Interfaces between the two technologies}

Table~\ref{tab:iface} lists every quantity that crosses from one technology to the other. The list is short on purpose. A shared memory of filter covariances would couple the software releases of the row code and the tracker code, and it is not required by either algorithm.

\begin{table}[ht]
\centering
\caption{Quantities that cross the interface between the two technologies.}
\label{tab:iface}
\begin{tabularx}{\textwidth}{@{}l l X@{}}
\toprule
Quantity & Direction & Use \\
\midrule
\(\mathcal{V}_{\mathrm{meter}}\) & row to tracker & Hard constraint on \(v\) \\
Realised \(v\), \(a\) & tracker to row & Feedforward \(\omega^{\star}\) and the repair test \\
\(x_{\mathrm{ref}}(s)\) & prescription to both & Hole period and local feasible set \\
Release log & row to seed map & Later survival counts, not the live constraint \\
Realised delay & row to itself & Tightens \(c\) if the budget is missed \\
\bottomrule
\end{tabularx}
\end{table}

Nothing in the table is a learned weight. Replacing the detector with a later model changes \(\tau_v\) and the confidence calibration, and it does not change the tracker's quadratic program. Replacing the elevation model changes the outer search, and it does not change Equation~\eqref{eq:plate}. That replaceability is the practical form of the claim that the two technologies meet only at speed.

\subsection{Build sequence}

The build is staged so that a failure is attributable to one technology. The first stage is the timber row unit on a soil bin or a greased belt, with the vehicle speed simulated by the belt, as in the classical meter tests \cite{iso7256,panning2000}. At this stage \(\mathcal{V}_{\mathrm{meter}}\) is checked directly: repair success against hole period, and boot motion against a shaker set to \(\lambda\) and to speeds around \SI{2.4}{m/s}. The second stage connects two rows and looks for the common mode Assumption~\ref{as:flex} neglected. The third stage is the vehicle on a flat field, tracker only, plates disengaged, confirming that the quadratic program meets Table~\ref{tab:budget} and that the resonant interval is refused. The fourth stage is the joint system on the sloped field of the estimation, or on the nearest available slope, with the indices of Equation~\eqref{eq:qfi} measured by the same camera that the controller uses and by a second, recording camera that the controller does not see. The second camera is the check Zhao and colleagues used when they compared a detector with a high-speed recording \cite{zhao2025yolo}. The estimation in the next section is the prediction the fourth stage would be asked to confirm or to reject. It is not a substitute for that stage.

\section{Estimation}
\label{sec:est}

\subsection{Protocol}

The estimation evaluates the closed-form models on the field of Table~\ref{tab:spec}: \SI{80}{m} across the slope, \SI{180}{m} along the slope, grade \num{8}\%, area \num{1.44}~hectares. Swath width is \SI{2.4}{m}. Target spacing is \SI{0.20}{m}. The open-loop branch uses Equations~\eqref{eq:openmiss} and the open-loop half of Equation~\eqref{eq:deptherr}. The proposed branch uses Equations~\eqref{eq:cfrac}, \eqref{eq:closedmiss} and the closed-loop half of Equation~\eqref{eq:deptherr}, and it selects heading and speed by Algorithm~\ref{alg:outer}. Failure probability uses Equation~\eqref{eq:fail} with \(\kappa=1/2\) and \(S=\exp(-\sigma_z^2/(2\cdot 6.5^2))\) for \(\sigma_z\) in millimetres. Work uses Equations~\eqref{eq:fpm} and \eqref{eq:work} with the coefficients listed in Section~\ref{sec:tech2}. No coefficient is adjusted after the tables are seen. The script that evaluates the expressions is a direct transcription of those equations, so a change in a specification value changes the table in a way a reader can repeat.

Three plans are reported in every comparison.

\begin{enumerate}[label=\textbf{B\arabic*.}]
\item Open-loop plate, no depth servo, heading straight up the slope, steady speed \SI{1.6}{m/s}. This is the plan a careful operator might choose: long rows, a moderate speed, no electronics.
\item Open-loop plate, no depth servo, the same uphill heading, steady speed \SI{2.4}{m/s}. This is the plan that sits on the timber resonance in order to finish sooner.
\item The proposed pair of technologies. Heading and speed are whatever Algorithm~\ref{alg:outer} returns under the stated weights and under the index constraints.
\end{enumerate}

B1 and B2 are baselines in the sense of Requirement~R8. They use the same work model, the same field and the same flexure as B3. They differ in the metering branch and in the freedom to choose heading and speed. A comparison against a different field, or against a steel arm with a different \(f_n\), would mix those effects with the effect of the technologies.

\subsection{Speed response of the metering model}

Table~\ref{tab:speed} evaluates the metering model at zero acceleration, before any field geometry is applied. It is the numerical form of Figure~\ref{fig:qfi} together with the depth columns.

\begin{table}[ht]
\centering
\caption{Metering model at zero acceleration. Displacements are amplitudes or root-mean-square errors in millimetres. \(c\) is the repair fraction.}
\label{tab:speed}
\small
\begin{tabular}{@{}rrrrrrrrr@{}}
\toprule
\(v\) & \(f\) & \(M(r)\) & \(d\) & \(c\) & \(\qfi_{\mathrm{ol}}\) & \(\qfi_{\mathrm{cl}}\) & \(\sigma_{z,\mathrm{ol}}\) & \(\sigma_{z,\mathrm{cl}}\) \\
\midrule
0.8 & 1.78 & 1.11 & 4.07 & 1.00 & 0.904 & 0.983 & 5.71 & 1.68 \\
1.0 & 2.22 & 1.19 & 4.35 & 1.00 & 0.881 & 0.981 & 5.91 & 1.79 \\
1.2 & 2.67 & 1.30 & 4.74 & 1.00 & 0.858 & 0.980 & 6.21 & 1.93 \\
1.4 & 3.11 & 1.45 & 5.31 & 1.00 & 0.834 & 0.978 & 6.65 & 2.14 \\
1.6 & 3.56 & 1.68 & 6.15 & 1.00 & 0.809 & 0.976 & 7.33 & 2.65 \\
1.8 & 4.00 & 2.04 & 7.47 & 1.00 & 0.781 & 0.974 & 8.47 & 3.52 \\
2.0 & 4.44 & 2.66 & 9.73 & 0.80 & 0.750 & 0.927 & 10.52 & 5.00 \\
2.2 & 4.89 & 3.85 & 14.10 & 0.55 & 0.709 & 0.848 & 14.65 & 7.90 \\
2.4 & 5.33 & 5.92 & 21.64 & 0.33 & 0.653 & 0.752 & 22.01 & 13.20 \\
2.6 & 5.78 & 5.23 & 19.12 & 0.15 & 0.643 & 0.686 & 19.54 & 12.67 \\
2.8 & 6.22 & 3.10 & 11.34 & 0.15 & 0.657 & 0.697 & 12.03 & 8.15 \\
\bottomrule
\end{tabular}
\end{table}

Several facts in Table~\ref{tab:speed} are worth separating, because they have different consequences. From \SI{0.8}{m/s} to \SI{1.8}{m/s} the repair fraction stays at one and the closed-loop quality-of-feed index stays between \num{0.983} and \num{0.974}. The open-loop index over the same interval falls from \num{0.904} to \num{0.781}. The technology's benefit on count accuracy is therefore not a thin peak at one speed. It is a plateau. Depth tells a different story. Closed-loop depth error grows from \SI{1.68}{mm} to \SI{3.52}{mm} as the roughness frequency approaches the \SI{3}{Hz} servo bandwidth, and it is still inside the \SI{4}{mm} requirement at \SI{1.8}{m/s}. At \SI{2.0}{m/s} the latency margin starts to erode, \(c\) falls to \num{0.80}, closed-loop quality of feed drops to \num{0.927}, and depth error reaches \SI{5.00}{mm}. Both index constraints of the outer program fail together near that speed. The resonant row at \SI{2.4}{m/s} is worse than a smooth extrapolation of the plateau: dynamic displacement \SI{21.6}{mm}, open-loop depth error \SI{22.0}{mm}, and a closed-loop depth error of \SI{13.2}{mm} that shows, numerically, the content of Proposition~\ref{prop:split}. The servo was not given the bandwidth to remove a \SI{5.3}{Hz} motion, and the model does not pretend that it has.

The acceleration check at fixed \(v=\SI{1.8}{m/s}\) is the headland effect in isolation. At \(a=0\), \(0.4\), \(0.8\) and \(\SI{1.2}{m/s^2}\) the repair fraction is \num{1.00}, \num{0.98}, \num{0.85} and \num{0.71}, and the closed-loop miss probability is \num{0.010}, \num{0.014}, \num{0.041} and \num{0.072}. The constraint \(c\ge 0.85\) therefore bites between \num{0.8} and \(\SI{1.2}{m/s^2}\). The turn acceleration used in the field estimate, \(\SI{0.7}{m/s^2}\), sits just on the safe side of that bite at \SI{1.8}{m/s} and more comfortably on the safe side at the lower working speeds the optimiser actually selects.

\begin{figure}[ht]
\centering
\begin{tikzpicture}
\begin{axis}[
  width=0.92\textwidth, height=6.4cm,
  xlabel={Ground speed \(v\) (m/s)},
  ylabel={Depth error (mm)},
  xmin=0.7, xmax=2.9, ymin=0, ymax=24,
  grid=both, legend pos=north west, legend style={font=\small}
]
\addplot[thick, mark=*] coordinates {
(0.8,5.71) (1.0,5.91) (1.2,6.21) (1.4,6.65) (1.6,7.33)
(1.8,8.47) (2.0,10.52) (2.2,14.65) (2.4,22.01) (2.6,19.54) (2.8,12.03)
};
\addplot[thick, mark=square*] coordinates {
(0.8,1.68) (1.0,1.79) (1.2,1.93) (1.4,2.14) (1.6,2.65)
(1.8,3.52) (2.0,5.00) (2.2,7.90) (2.4,13.20) (2.6,12.67) (2.8,8.15)
};
\draw[dashed] (axis cs:0.7,4) -- (axis cs:2.9,4);
\legend{Open-loop model, CAVM model, \SI{4}{mm} limit}
\end{axis}
\end{tikzpicture}
\caption{Modelled depth root-mean-square error. The dashed line is Requirement~R2. The peak is the timber resonance, which the speed set is required to avoid.}
\label{fig:depth}
\end{figure}

\subsection{Heading and speed on the estimation field}

Table~\ref{tab:field} applies the field geometry at four speeds that the metering plateau still regards as interesting. Energy changes little with speed, because aerodynamic drag is a small term beside rolling resistance at \SI{1}{m/s} to \SI{2}{m/s}. Energy changes more with heading. At \SI{1.5}{m/s} the uphill heading costs \SI{4.782}{MJ} and the best of the oblique headings shown, \ang{55} from uphill, costs \SI{4.416}{MJ}. The reduction is \((4.782-4.416)/4.782=\num{7.7}\%\). The pure contour costs \SI{4.451}{MJ}, slightly more than the \ang{55} heading, because the saving in grade loss is spent again on \num{75} turns against \num{48} turns. Time moves the other way: the contour takes \SI{1.455}{h} at this speed and the uphill heading takes \SI{1.289}{h}. Count accuracy at a given speed is almost indifferent to heading. The closed-loop quality-of-feed index at \SI{1.5}{m/s} is \num{0.977} on every heading in the table, and the accelerated fraction is only a few percent of the row. Heading is an energy and time decision inside the plateau. It becomes a singulation decision only when the accelerated fraction grows, which is the short-row, high-speed corner of the table.

\begin{table}[ht]
\centering
\caption{Field model with the closed-loop meter. Work is non-conservative mechanical work. The accelerated fraction \(\phi_{\mathrm{acc}}\) is the share of row length run at turn acceleration.}
\label{tab:field}
\small
\begin{tabular}{@{}llrrrrrr@{}}
\toprule
\(v\) & Heading & \(n\) & Work & Time & \(\phi_{\mathrm{acc}}\) & \(\qfi\) & \(\sigma_z\) \\
(m/s) & & & (MJ) & (h) & & & (mm) \\
\midrule
1.2 & uphill & 34 & 4.778 & 1.572 & 0.011 & 0.980 & 1.94 \\
1.2 & \ang{55} & 48 & 4.413 & 1.609 & 0.016 & 0.980 & 1.94 \\
1.2 & contour & 75 & 4.447 & 1.733 & 0.026 & 0.980 & 1.95 \\
1.5 & uphill & 34 & 4.782 & 1.289 & 0.018 & 0.977 & 2.36 \\
1.5 & \ang{55} & 48 & 4.416 & 1.331 & 0.026 & 0.977 & 2.36 \\
1.5 & contour & 75 & 4.451 & 1.455 & 0.040 & 0.977 & 2.37 \\
1.8 & uphill & 34 & 4.786 & 1.100 & 0.026 & 0.973 & 3.53 \\
1.8 & \ang{55} & 48 & 4.420 & 1.146 & 0.037 & 0.973 & 3.54 \\
1.8 & contour & 75 & 4.455 & 1.270 & 0.058 & 0.972 & 3.55 \\
2.0 & uphill & 34 & 4.789 & 1.006 & 0.032 & 0.924 & 5.02 \\
2.0 & \ang{55} & 48 & 4.423 & 1.053 & 0.046 & 0.923 & 5.03 \\
2.0 & contour & 75 & 4.458 & 1.177 & 0.071 & 0.921 & 5.05 \\
\bottomrule
\end{tabular}
\end{table}

Algorithm~\ref{alg:outer}, with the index filter \(\qfi\ge 0.955\) and \(\sigma_z\le\SI{4}{mm}\), and with the resonant interval removed, returns a heading of \ang{55} from uphill and a working speed of \SI{1.4}{m/s}. Neighbouring speeds inside the plateau are close in cost. The optimiser's preference for \SI{1.4}{m/s} rather than \SI{1.2}{m/s} is the time weight. Its refusal of \SI{1.8}{m/s} and above is the depth constraint, not the time weight: at \SI{1.8}{m/s} the depth error is still under \SI{4}{mm} in Table~\ref{tab:field}, and at the next step the filter fails. The preference for \ang{55} over the pure contour is the turn cost. The preference for \ang{55} over straight uphill is the grade term in Equation~\eqref{eq:fpm}. No single literature objective---minimum length, minimum energy, or minimum time---would have returned this pair by itself. Minimum time wants a higher speed and the uphill heading. Minimum energy on this rectangle wants a heading near \ang{55} to \ang{90} and is almost indifferent to speed. The singulation constraint is what keeps the speed on the plateau, and the turn term is what keeps the heading from sliding all the way to the contour.

\begin{table}[ht]
\centering
\caption{Plans compared on the estimation field. B3 is the proposed pair at the heading and speed Algorithm~\ref{alg:outer} returns.}
\label{tab:plans}
\begin{tabular}{@{}lrrrrr@{}}
\toprule
Plan & Work & Time & \(\qfi\) & \(\sigma_z\) & \(\rho_{\mathrm{fail}}\) \\
 & (MJ) & (h) & & (mm) & \\
\midrule
B1 open, uphill, \SI{1.6}{m/s} & 4.783 & 1.218 & 0.808 & 7.35 & 0.561 \\
B2 open, uphill, \SI{2.4}{m/s} & 4.797 & 0.864 & 0.652 & 22.04 & 0.998 \\
B3 proposed, \ang{55}, \SI{1.4}{m/s} & 4.415 & 1.410 & 0.978 & 2.15 & 0.068 \\
\bottomrule
\end{tabular}
\end{table}

Table~\ref{tab:plans} is the joint answer to Question~3. Against B1, the proposed plan reduces non-conservative work from \SI{4.783}{MJ} to \SI{4.415}{MJ}, which is \num{7.7}\%, and it reduces the modelled failure probability from \num{0.561} to \num{0.068}. The time cost of that reduction is real: \SI{1.410}{h} against \SI{1.218}{h}, about \num{16}\% longer than B1. Against B2, the proposed plan uses slightly less work, takes longer, and avoids a regime in which the model assigns a failure probability of \num{0.998}. The resonant plan is the fastest and, on this model, agronomically empty. The survival function is already below one half at a depth error of \SI{8}{mm} and is exhausted at \SI{22}{mm}. That conclusion is only as strong as \(S(\sigma_z)\). It is, however, the same functional form at every row of the table, so the ranking does not depend on a special choice made for B2 alone. B2 fails the depth requirement inside the open-loop model even before \(S\) is applied: \(\sigma_z=\SI{22}{mm}\) against a band of \(\pm\SI{8}{mm}\).

The work reduction of \num{7.7}\% sits a little below the \numrange{9.5}{15.6}\% direction savings Vahdanjoo, Gislum and S{\o}rensen reported, on different fields and with their own energy model \cite{vahdanjoo2023}. The present rectangle has a moderate grade and a slip coefficient that was chosen not to dominate rolling resistance. A steeper restoration slope would increase the gap between uphill and oblique headings and would make the energy term louder relative to the time term. It would not remove the resonant exclusion, which does not depend on grade. The \num{8.56}\% efficiency gain Xu, Lai, Guo and Lu found for a tracker allowed to pursue efficiency is a different quantity, measured against a pure tracking controller rather than against an uphill heading \cite{xu2022}. It is cited here as evidence that putting a non-tracking term inside a predictive tracker moves the result by an amount of this order, not as a number the present field is supposed to reproduce.

Seed count on B3 follows from the swath length. At \ang{55} the search uses a row length of \SI{125}{m} and \num{48} swaths, so the distance travelled while planting is \SI{6000}{m} per row and \SI{24000}{m} over four rows. At \SI{0.20}{m} spacing that is \(1.2\times 10^{5}\) seeds if every swath is full, which is the count \(N_{\mathrm{seed}}\) used to weight \(\rho_{\mathrm{fail}}\). A failure probability of \num{0.068} against \num{0.561} is a difference of about \(5.9\times 10^{4}\) seeds on that count. The figure is a model difference in failed placements, not a count of emerged seedlings. It is large because the open-loop depth error of \SI{7.35}{mm} already sits where \(S(\sigma_z)\) is falling steeply, and because the open-loop miss probability at \SI{1.6}{m/s} is \num{0.146}. Both contributions are visible in the formulae, and neither is a field emergence count.

\subsection{Where the gain comes from}

The gain splits into a part the first technology would produce on a fixed heading and a part that appears only when the heading is free. On the uphill heading at \SI{1.4}{m/s}, the closed-loop model already reaches \(\qfi=\num{0.978}\) and \(\sigma_z=\SI{2.15}{mm}\), with failure probability about \num{0.067}. That is almost the whole of the index gain of B3 over B1, and it is available without any coverage technology at all, at the price of the uphill work \SI{4.780}{MJ}. The coverage technology then moves the same metering performance onto a heading that costs \SI{4.415}{MJ}. Readers who want the singulation result without a new planner can stop at the first technology and a disciplined speed. Readers who want the work result need the second technology, and they only keep the singulation result if the planner is forbidden from spending the work saving on a higher speed. The joint design is the statement that the work saving will not be spent that way.

A further split, inside the first technology, separates the repair of empty cells from the depth servo. At \SI{1.4}{m/s} the open-loop miss probability from Equation~\eqref{eq:openmiss} is \num{0.125} and the closed-loop miss probability is \num{0.008}. That gap is the repair term, and it exists only while \(c=1\). The open-loop depth error is \SI{6.65}{mm} and the closed-loop depth error is \SI{2.14}{mm}. That gap is the servo, and Proposition~\ref{prop:split} says it will shrink if the machine is driven toward resonance. Table~\ref{tab:speed} at \SI{2.4}{m/s} is the shrinkage: the closed loop still improves depth relative to the open loop, from \SI{22.0}{mm} to \SI{13.2}{mm}, and \SI{13.2}{mm} is still a failed band. Vision metering without the speed constraint is not the technology that was specified. The tables that show a \num{0.978} index are tables in which the constraint is active as a restriction on the search, even though at the chosen point it is not active as a saturated inequality.

\subsection{Sensitivity the estimation can still speak about}

Three specification numbers have enough leverage that a reader will want to know the direction of a change, even though the estimation does not rerun a full identification.

If the equivalent stiffness is doubled, \(f_n\) rises by \(\sqrt{2}\) and the resonant speed for \(\lambda=\SI{0.45}{m}\) moves from about \SI{2.4}{m/s} to about \SI{3.4}{m/s}. The latency limit near \SI{1.9}{m/s} then binds before the resonance does. The forbidden interval moves out of the speeds a planter is likely to request, which is the steel-arm situation the introduction described. The depth servo still has to remove the static deflection, which would itself fall because \(k\) is larger. A stiffer arm makes the first technology easier and makes the resonance penalty in the tracker less often active. It does not retire the latency inequality. Borja's sign change did not require a wooden arm \cite{borja2021}.

If the detector slips from \SI{14}{ms} to \SI{40}{ms} of inference, \(\tau_v\) becomes \SI{48}{ms} and \(\tau\) becomes \SI{83}{ms}. The speed at which \(c\) remains \num{1} falls in proportion to the hole period. The plateau shrinks from the right. Algorithm~\ref{alg:outer} then returns a slower \(v_{\mathrm{work}}\), the time term grows, and the index constraints can still be met if any plateau remains. If the plateau disappears, the outer program returns no feasible candidate, which is a more honest outcome than a plan with a decorative camera. This is the microcontroller outcome, predicted by the same inequality that organised Section~\ref{sec:prev-meter}.

If the grade doubles and the slip coefficient stays put, \(\sin^2\alpha\) nearly quadruples because \(\alpha\) is small, and the work gap between uphill and oblique headings widens. The optimiser's preference for an oblique heading strengthens. The metering table does not move. A field trial on a steeper site would be expected to show a larger energy difference and the same speed ceiling, which is a falsifiable split between the two technologies.

If woody-seed detector recall is worse than the peanut and maize figures used in Equation~\eqref{eq:closedmiss}, every closed-loop miss probability rises. The estimation has no new image set with which to put a number on that rise. The direction is unambiguous, and it is the reason Assumption~\ref{as:est} forbids reading \num{0.978} as an expected field index for a named species. The build sequence in Section~\ref{sec:build} makes the belt test, not this table, the place where that index is first allowed to be empirical.

\subsection{Comparison with the nearest published systems}

A fair comparison states what was measured there and what is only modelled here. Borja's vision-assisted maize meter reduced missed hills at \SI{2}{km/h} and \SI{4}{km/h} and increased them at \SI{6}{km/h} \cite{borja2021}. The present model reproduces the shape of that reversal, through \(c(v,a)\), and moves the speed at which the reversal begins by assuming a \SI{22}{ms} vision budget instead of a microcontroller budget. It does not reproduce their hill counts. Mangus's electric-drive meter held a mean singulation of \num{98.45}\% over a wide speed range and lost accuracy when acceleration or seed rate rose \cite{mangus2017}. The present closed-loop plateau at \num{0.97} to \num{0.98} is a model expression of that kind of singulation, restricted to the speeds at which a repair can finish, and the acceleration term in Equation~\eqref{eq:openmiss} is a model expression of their acceleration finding. Blaser's and Zhao's detectors supply the error rates the repair formula is allowed to believe \cite{blaser2024,zhao2025yolo}. They do not supply a wooden tube in the training set.

On the planning side, Vahdanjoo, Gislum and S{\o}rensen showed that a three-dimensional energy criterion moves the preferred seeding direction by enough to matter \cite{vahdanjoo2023}. The oblique heading in Table~\ref{tab:plans} is the same kind of decision, on a simpler elevation model, with extra constraints. Zhang and colleagues' seeding-robot planner reports a missed-sowing rate of \num{1.8}\% together with a large energy reduction in their own trials \cite{zhang2025path}. The present model at B3 assigns a miss probability near \num{0.008} on the plateau, which is a lower miss rate than \num{1.8}\% and is not a competing experimental claim. It is a statement that, if the detector and the stepper perform as the repair formula assumes, empty-cell misses are almost removed and the remaining failure is depth and the occasional multiple. Their planner and this one can be combined: their search can carry Equation~\eqref{eq:fail} as an edge cost, and the inner tracker of this thesis can track the path they return, provided the path's speed profile is projected onto \(\mathcal{V}_{\mathrm{meter}}\) before tracking begins.

No published system reviewed in Chapter~\ref{ch:previous} closes the loop from a timber natural frequency to a coverage constraint. That closure, rather than a new record on maize singulation, is the comparison this estimation is built to illustrate.

\subsection{Limits}

The limits are the assumptions, restated against the tables so that the largest numbers in the chapter stay attached to them.

The failure probabilities inherit \(S(\sigma_z)\) and \(\kappa=1/2\). A species whose emergence is tolerant of a \SI{7}{mm} depth error would shrink the gap between B1 and B3. A species as intolerant as the spinifex in Masarei, Guzzomi, Merritt and Erickson's trial would enlarge it and would also force a smaller \(\sigma_{z,\max}\) than \SI{4}{mm} \cite{masarei2019}. The design case at \SI{18}{mm} with an \SI{8}{mm} band is neither of those species.

The open-loop intercepts in Equation~\eqref{eq:openmiss} were chosen to sit in a plausible band for irregular seed. A plate that has already been tuned for the woody species might start from a better open-loop index, and the vision benefit on misses would then be smaller. The resonant depth peak would remain, because it is produced by \(M(r)\) and by the \SI{4}{mm} unresolved bias, not by the miss intercept.

Chassis coupling, mixed species, wheel slip beyond the force model, and detector performance on a dusty wooden tube are outside the numerical evaluation. The build sequence is where they stop being outside the project. Until that sequence is run, Table~\ref{tab:plans} is a comparison of plans under a declared model. It is evidence that the constraints are coherent and that they bind in the right places: latency at one speed, resonance at another, grade and turns in the heading. It is not evidence that a particular laminated arm has \(k=\SI{8200}{N/m}\), and it is not evidence that a stand of woody seedlings will emerge in the ratio \(1-0.068\).

\subsection{Checking the selected plan against the requirements}

Requirement~R8 asked for the baselines to sit in the same table as the proposal, and Table~\ref{tab:plans} is that table. The other requirements can be checked on B3 without a further search. The closed-loop quality-of-feed index at the selected point is \num{0.978}, which clears R1's floor of \num{0.95}. The depth error is \SI{2.15}{mm}, which clears R2's ceiling of \SI{4}{mm}. The speed \SI{1.4}{m/s} lies inside the full-repair side of Equation~\eqref{eq:latency}, so R3 is met as a constraint that discarded faster candidates rather than as a penalty the weights happened to respect. The resonant interval is not the selected speed, and it is not inside \(\mathcal{V}_{\mathrm{meter}}\), so R4 holds for the plan even though a human override could still violate it if the interlock in Section~\ref{sec:build} were omitted. The cost that ranked \ang{55} above uphill and above the pure contour contains work, time and \(\rho_{\mathrm{fail}}\), and \(\rho_{\mathrm{fail}}\) uses the closed-loop miss model, so R5 holds. Four independent row tasks are in the build, so R6 holds as an architectural requirement the estimation's single-row formulas do not by themselves prove. The tracker period and the quadratic program are specified, so R7 holds as a budget, not as a measured solve time. The language of the tables is model language, so R9 holds if the surrounding prose is read with the tables. One speed set is published by the meter and consumed by the tracker, so R10 holds by the interface in Table~\ref{tab:iface} and by the prescription example in which a shorter spacing moves that same set.

The check is worth doing explicitly because a plan can win a table and still break a requirement that was never a column. B2, for example, wins the time column and fails R1, R2 and R4 at once. B1 meets none of the index requirements and was never required to: it is the baseline those requirements are measured against. B3 is the only one of the three that is feasible for the joint problem as Definition~\ref{def:p3} stated it. That is a stronger statement than ``B3 has the smallest failure probability'', and it is the statement the requirement list was written to support.

\subsection{A measurement sheet for the fourth build stage}

The fourth stage of the build is the first moment at which Table~\ref{tab:plans} could be wrong in the empirical sense. The sheet below is the comparison that stage should record, so that a disagreement has an address. Each row is one operating point already present in the model tables. The recording camera, not the control camera, supplies \(\qfi\), \(\mis\) and \(\mul\) by the classification of Equation~\eqref{eq:classes} on at least five hundred spacings, which is the order of count at which a miss rate of one percent is distinguishable from a miss rate of three percent. Depth is the mean and the root-mean-square of boot penetration against the \SI{18}{mm} target, measured by the ultrasonic sensor and, on a calibration strip, by a manual check of uncovered seed. Work is not measured as a germination count. It is measured as drawbar pull or as motor current integrated over a swath of known length, on the uphill heading and on the \ang{55} heading, at the same speed, so that the \num{7.7}\% gap has a chance to appear or to fail without being mixed with a speed change.

\begin{table}[ht]
\centering
\caption{Points the field stage should record beside the model. Model columns repeat Chapter~2 and this chapter. The empty columns are the measurement.}
\label{tab:sheet}
\small
\begin{tabular}{@{}llrrrr@{}}
\toprule
Point & Quantity & Model & Recorded & Unit & Pass if \\
\midrule
\(v=\SI{1.4}{m/s}\), straight & \(\qfi_{\mathrm{cl}}\) & 0.978 & & -- & \(\ge 0.95\) \\
\(v=\SI{1.4}{m/s}\), straight & \(\sigma_z\) & 2.15 & & mm & \(\le 4\) \\
\(v=\SI{1.4}{m/s}\), \(a=\SI{0.7}{m/s^2}\) & \(c\) stays high & yes & & -- & repairs still timely \\
\(v=\SI{2.4}{m/s}\), straight & \(\sigma_{z,\mathrm{ol}}\) & 22 & & mm & resonance visible \\
\(v=\SI{2.4}{m/s}\), commanded & speed refused & yes & & -- & interlock holds \\
uphill versus \ang{55} & work gap & 7.7 & & \% & same sign \\
realised delay & \(\tau_v+\tau_a\) & 57 & & ms & budget holds \\
\bottomrule
\end{tabular}
\end{table}

Two outcomes of that sheet would change the design rather than only the numbers. If the resonant peak is absent, \(k\) or \(\lambda\) differs from the specification and the forbidden interval should be recomputed from the measured frequency, not deleted from the tracker on the strength of one smooth run. A single run can miss a wavelength. If the resonant peak is present and the closed-loop depth error at \SI{1.4}{m/s} is nevertheless above \SI{4}{mm}, the servo bandwidth or the bias estimate is wrong and Proposition~\ref{prop:split} still stands: the remedy is not to drive faster. If realised delay is near \SI{100}{ms}, the plateau shrinks and Algorithm~\ref{alg:outer} must be rerun with the new \(\tau\) before any claim about \SI{1.4}{m/s} is repeated. If the work gap changes sign, the slip coefficient or the grade of the site is not the one in Equation~\eqref{eq:fpm}, and the heading search should be rerun with the measured resistance. In every one of those cases the structure of the two technologies stays. What moves is a specification value that the estimation, by construction, held fixed.

The sheet also records what a negative singulation result would mean. A quality-of-feed index of \num{0.90} at \SI{1.4}{m/s} with a healthy delay budget would say that Equation~\eqref{eq:closedmiss} was optimistic about detector recall on this seed, which Assumption~\ref{as:est} already flagged. The speed constraint would remain, because it did not depend on the recall figure. The depth loop would remain, because it does not use the detector. The coverage ranking of headings would remain, because inside the plateau the miss rate barely changes with heading. The number \num{0.978} would be retired. The architecture would not. That separation is the useful product of having written the model in pieces that can fail separately.

\subsection{A short spacing tape, scored by hand}

The indices are easier to trust after they have been scored once by hand, on a tape short enough to add up. Suppose twenty successive spacings, in millimetres, are recorded on the belt at the design target of \SI{200}{mm}:
\begin{equation*}
\begin{aligned}
&210,\; 190,\; 205,\; 90,\; 310,\; 200,\; 198,\; 202,\; 215,\; 185,\\
&400,\; 188,\; 207,\; 195,\; 220,\; 70,\; 130,\; 203,\; 197,\; 208.
\end{aligned}
\end{equation*}
The class boundaries are \SI{100}{mm} and \SI{300}{mm}. The value \SI{90}{mm} is a multiple. The value \SI{310}{mm} that follows it is what a double often leaves behind, and it is a miss. The value \SI{400}{mm} is a miss. The pair \SI{70}{mm} and \SI{130}{mm} contains one multiple; \SI{130}{mm} is a single because it clears half the target. Sixteen spacings are singles. Then \(N_{\mathrm{tot}}=20\), \(N_u=2\), \(N_m=2\), \(N_s=16\), and
\begin{equation}
\label{eq:hand}
\qfi = 0.80,\qquad \mul = 0.10,\qquad \mis = 0.10.
\end{equation}
The single-class mean and the coefficient of variation use only those sixteen singles, not the \SI{90}{mm} or the \SI{400}{mm}. Including those two would inflate the coefficient of variation and would hide, inside one number, the fact that most seeds were close to the target and two events were categorical failures. That is the distinction Kachman and Smith insisted on, and it is why the estimation reports \(\qfi\), \(\mis\), \(\mul\) and a depth error rather than one coefficient of variation for the whole tape \cite{kachman1995}. A field stage that published only a coefficient of variation would not be able to confirm or reject Table~\ref{tab:plans}.

Twenty spacings are an illustration, not a sample size. Five hundred spacings, as the measurement sheet asks, make a miss rate of \num{0.01} a handful of events and a miss rate of \num{0.15}, the open-loop model at \SI{1.6}{m/s}, a hundred and fifty events. Those two samples cannot be confused. The estimation's claim that B3 and B1 are different is a claim about that separation under the model. The tape shows the arithmetic the field stage must not replace with a smoother and less informative summary.

\subsection{Reading the energy prize at the scale of one hectare}

The \SI{0.37}{MJ} difference between B1 and B3 is \SI{0.10}{kWh}. On a \num{1.44}~hectare field that is a small energy quantity beside the fuel a truck spends to deliver the machine, and it is the right order of magnitude for a direction choice on a moderate grade. Reported direction savings of \numrange{9.5}{15.6}\% on other seeding fields are the same kind of quantity: a fraction of an already-modest in-field energy, not a new power source \cite{vahdanjoo2023}. The estimation would mislead a reader who remembered only the percentage and forgot the base. \num{7.7}\% of \SI{4.78}{MJ} is the prize. The singulation and depth columns are where the technologies change the operation by a factor rather than by a fraction. Failure probability falls from \num{0.561} to \num{0.068}, a factor of about eight, under a survival function that treats a \SI{7}{mm} depth error as already serious. Even if a later germination trial halves that factor, the index gap remains the reason to build the vision loop, and the energy gap remains a constraint on how the planner is allowed to spend the time the loop requires.

Per hectare, using \(1.44\)~ha as the divisor, B3 uses about \SI{3.07}{MJ/ha} and about \SI{0.98}{h/ha}, and it assigns a failed-placement fraction of \num{0.068}. B1 uses about \SI{3.32}{MJ/ha} and about \SI{0.85}{h/ha}, with a failed-placement fraction of \num{0.561}. A contractor who bills only time will prefer B1 and will plant a stand the model regards as mostly outside the depth band. A contractor who bills for established plants, which is the usual contract in restoration, will prefer B3 even before the energy line is read. Pedersen, Fountas, Have and Blackmore made the same structural point for agricultural robots in general: the robot pays when the operation it enables is one the cheaper machine misses, not when it is a few percent cheaper at the same operation \cite{pedersen2006}. The wooden planter with these two technologies is in that category if, and only if, the depth and singulation gaps survive the measurement sheet. The energy line is not the category.

\subsection{The same rectangle turned the other way}

The estimation field was chosen so that a contour heading has short rows. That choice is doing real work, and it should be shown by computing the opposite geometry with the same equations and the same coefficients. Turn the rectangle so that the long side, \SI{180}{m}, lies on the contour and the short side, \SI{80}{m}, runs up the \num{8}\% grade. At \SI{1.4}{m/s} the work model then assigns \SI{4.16}{MJ} and \SI{1.37}{h} to the contour heading, with \num{34} turns, and \SI{5.12}{MJ} and \SI{1.53}{h} to the uphill heading, with \num{75} turns. Contour travel is cheaper per metre because \(\Gamma(\ang{90})=1\), and it is also shorter in time because the long rows cut the number of turns. Both terms agree. Algorithm~\ref{alg:outer} and a planner that minimises only work return the same heading. The singulation constraint still chooses the speed, because nothing in the rotation moves \(f_n\) or \(\tau\), and the selected speed remains \SI{1.4}{m/s} on the plateau of Figure~\ref{fig:qfi}.

On the original field the agreement breaks. Uphill travel at \SI{1.4}{m/s} costs \SI{4.78}{MJ} and \SI{1.37}{h}. Contour travel costs \SI{4.45}{MJ} and \SI{1.53}{h}. The \ang{55} heading returned by the joint search costs \SI{4.41}{MJ} and \SI{1.41}{h}. Pure minimum work and the joint search are close, and they are not identical: the turn term and the time weight pull the optimum off the pure contour, by a few tenths of a megajoule, toward a longer row. A reader who only remembers that ``the planner chose an oblique heading'' has missed the smaller and more useful fact. The oblique heading is the compromise on a field that punishes the contour with turns. On a field that rewards the contour with both energy and time, the compromise disappears and the constraint that remains is the speed ceiling from the wooden arm and the vision budget. The two technologies are still both active. One of them has nothing left to trade.

This rotation is the last numerical check the estimation makes. It uses no new coefficient. It shows that the heading result in Table~\ref{tab:plans} is a property of the field orientation fixed in the problem set, not a universal preference for \ang{55}. The speed result is the one that travels. Wherever the long side points, steady running stays off the resonance and inside the repair inequality, or the outer program returns no plan.

Within those limits the estimation answers the three research questions in the only sense a design thesis can answer them before a prototype. The plate law and the flexure filter of Section~\ref{sec:tech1} repair empty cells only while the repair fraction is high, and they assign mean depth to a servo whose bandwidth sits below \(f_n\). The outer search and the inner quadratic program of Section~\ref{sec:tech2} select a heading and a speed inside that repair set, with grade work, turn time and failed placements in one cost. On the estimation field the selected plan is \ang{55} from uphill at \SI{1.4}{m/s}, with modelled quality of feed \num{0.978}, depth error \SI{2.15}{mm} and non-conservative work \num{7.7}\% below the uphill open-loop plan at \SI{1.6}{m/s}. The resonant open-loop plan is faster and, in the same model, does not establish the crop. The prototype specified in Section~\ref{sec:build} is the instrument that would turn these statements into measurements.

\clearpage
\phantomsection
\addcontentsline{toc}{chapter}{References}
\begin{thebibliography}{99}

\bibitem{anantachar2010}
Anantachar, M., Kumar, G.~V.~P., and Guruswamy, T.
Neural network prediction of performance parameters of an inclined plate seed metering device and its reverse mapping for the determination of optimum design and operational parameters.
\textit{Computers and Electronics in Agriculture}, 72(2):87--98, 2010.

\bibitem{andersson2019}
Andersson, J.~A.~E., Gillis, J., Horn, G., Rawlings, J.~B., and Diehl, M.
CasADi: a software framework for nonlinear optimization and optimal control.
\textit{Mathematical Programming Computation}, 11:1--36, 2019.

\bibitem{asabe313}
ASABE.
Soil cone penetrometer.
ASAE S313.3, American Society of Agricultural and Biological Engineers, St.~Joseph, MI.

\bibitem{auat2013}
Auat Cheein, F.~A. and Carelli, R.
Agricultural robotics: unmanned robotic service units in agricultural tasks.
\textit{IEEE Industrial Electronics Magazine}, 7(3):48--58, 2013.

\bibitem{backman2012}
Backman, J., Oksanen, T., and Visala, A.
Navigation system for agricultural machines: nonlinear model predictive path tracking.
\textit{Computers and Electronics in Agriculture}, 82:32--43, 2012.

\bibitem{bechar2016}
Bechar, A. and Vigneault, C.
Agricultural robots for field operations: concepts and components.
\textit{Biosystems Engineering}, 149:94--111, 2016.

\bibitem{blackmore2005}
Blackmore, S., Stout, B., Wang, M., and Runov, B.
Robotic agriculture --- the future of agricultural mechanisation?
In \textit{5th European Conference on Precision Agriculture}, pages 621--628, Uppsala, 2005.

\bibitem{blaser2024}
Blaser, M., Porter, W., Fuhrer, L., and Bourlai, T.
YOLOv10 computer vision performance measurement for agricultural vacuum seed meters.
In \textit{IEEE International Conference on Big Data}, pages 1929--1936, 2024.

\bibitem{bochtis2008}
Bochtis, D.~D. and Vougioukas, S.~G.
Minimising the non-working distance travelled by machines operating in a headland field pattern.
\textit{Biosystems Engineering}, 101(1):1--12, 2008.

\bibitem{bochtis2014}
Bochtis, D.~D., S{\o}rensen, C.~G., and Busato, P.
Advances in agricultural machinery management: a review.
\textit{Biosystems Engineering}, 126:69--81, 2014.

\bibitem{bond1999}
Bond, W.~J., Honig, M., and Maze, K.~E.
Seed size and seedling emergence: an allometric relationship and some ecological implications.
\textit{Oecologia}, 120:132--136, 1999.

\bibitem{borja2018}
Borja, A.~A., Amongo, R.~M.~C., Suministrado, D.~C., and Pabico, J.~P.
A machine vision assisted mechatronic seed meter for precision planting of corn.
In \textit{3rd International Conference on Control and Robotics Engineering}, pages 183--187, 2018.

\bibitem{borja2021}
Borja, A.~A., Amongo, R.~M., Suministrado, D.~C., and Pabico, J.~P.
Design and evaluation of a machine vision system and mechatronic drive for a pneumatic seed meter for corn.
\textit{Philippine Journal of Agricultural and Biosystems Engineering}, 17(1), 2021.

\bibitem{boyd2004}
Boyd, S. and Vandenberghe, L.
\textit{Convex Optimization}.
Cambridge University Press, 2004.

\bibitem{bustos2025}
Bustos-Gayt{\'a}n, A., Salda{\~n}a-Robles, N., Cepeda-Negrete, J., and co-authors.
Performance evaluation of vacuum seed-metering devices using machine vision: influence of components and operating parameters on seeding uniformity.
\textit{Journal of Biosystems Engineering}, 2025.
doi:10.1007/s42853-025-00295-8.

\bibitem{choset2001}
Choset, H.
Coverage for robotics --- a survey of recent results.
\textit{Annals of Mathematics and Artificial Intelligence}, 31:113--126, 2001.

\bibitem{cundall1979}
Cundall, P.~A. and Strack, O.~D.~L.
A discrete numerical model for granular assemblies.
\textit{G{\'e}otechnique}, 29(1):47--65, 1979.

\bibitem{dijkstra1959}
Dijkstra, E.~W.
A note on two problems in connexion with graphs.
\textit{Numerische Mathematik}, 1:269--271, 1959.

\bibitem{dorigo2006}
Dorigo, M., Birattari, M., and St{\"u}tzle, T.
Ant colony optimization.
\textit{IEEE Computational Intelligence Magazine}, 1(4):28--39, 2006.

\bibitem{dibbler}
Duke, M.~D., McGuinness, B.~J., and Kunnemeyer, R.
Development of a mechatronic dibbling machine for improving the quality of forestry seedlings.
University of Waikato Research Commons, 2016.

\bibitem{edwards2017}
Edwards, G.~T.~C., Hinge, J., Skou-Nielsen, N., Villa-Henriksen, A., S{\o}rensen, C.~A.~G., and Green, O.
Route planning evaluation of a prototype optimised infield route planner for neutral material flow agricultural operations.
\textit{Biosystems Engineering}, 153:149--157, 2017.

\bibitem{ehsani2004}
Ehsani, M.~R., Upadhyaya, S.~K., and Mattson, M.~L.
Seed location mapping using RTK GPS.
\textit{Transactions of the ASAE}, 47(3):909--914, 2004.

\bibitem{ferreau2014}
Ferreau, H.~J., Kirches, C., Potschka, A., Bock, H.~G., and Diehl, M.
qpOASES: a parametric active-set algorithm for quadratic programming.
\textit{Mathematical Programming Computation}, 6(4):327--363, 2014.

\bibitem{fountas2015}
Fountas, S., Carli, G., S{\o}rensen, C.~G., Tsiropoulos, Z., Cavalaris, C., Vatsanidou, A., Liakos, B., Canavari, M., Wiebensohn, J., and Tisserye, B.
Farm management information systems: current situation and future perspectives.
\textit{Computers and Electronics in Agriculture}, 115:40--50, 2015.

\bibitem{galceran2013}
Galceran, E. and Carreras, M.
A survey on coverage path planning for robotics.
\textit{Robotics and Autonomous Systems}, 61(12):1258--1276, 2013.

\bibitem{godwin2007}
Godwin, R.~J. and O'Dogherty, M.~J.
Integrated soil tillage force prediction models.
\textit{Journal of Terramechanics}, 44(1):3--14, 2007.

\bibitem{griepentrog2005}
Griepentrog, H.~W., N{\o}rremark, M., Nielsen, H., and Blackmore, B.~S.
Seed mapping of sugar beet.
\textit{Precision Agriculture}, 6:157--165, 2005.

\bibitem{hameed2010}
Hameed, I.~A., Bochtis, D.~D., S{\o}rensen, C.~G., and N{\o}rremark, M.
Automated generation of guidance lines for operational field planning.
\textit{Biosystems Engineering}, 107(4):294--306, 2010.

\bibitem{hart1968}
Hart, P.~E., Nilsson, N.~J., and Raphael, B.
A formal basis for the heuristic determination of minimum cost paths.
\textit{IEEE Transactions on Systems Science and Cybernetics}, 4(2):100--107, 1968.

\bibitem{heege2013}
Heege, H.~J.
\textit{Precision in Crop Farming: Site Specific Concepts and Sensing Methods}.
Springer, Dordrecht, 2013.

\bibitem{iso7256}
International Organization for Standardization.
Sowing equipment --- test methods --- part 1: single seed drills (precision drills).
ISO 7256-1:1984.

\bibitem{iso11783}
International Organization for Standardization.
Tractors and machinery for agriculture and forestry --- serial control and communications data network.
ISO 11783.

\bibitem{jab2010}
Bioresources Engineering, University of KwaZulu-Natal.
Design and evaluation of a jab planter.
CSBE/SCGAB Paper CSBE101037, 2010.

\bibitem{jin2010}
Jin, J. and Tang, L.
Optimal coverage path planning for arable farming on 2D surfaces.
\textit{Transactions of the ASABE}, 53(1):283--295, 2010.

\bibitem{jin2011}
Jin, J. and Tang, L.
Coverage path planning on three-dimensional terrain for arable farming.
\textit{Journal of Field Robotics}, 28(3):424--440, 2011.

\bibitem{julier1997}
Julier, S.~J. and Uhlmann, J.~K.
A new extension of the Kalman filter to nonlinear systems.
In \textit{Proceedings of AeroSense: the 11th International Symposium on Aerospace/Defense Sensing, Simulation and Controls}, pages 182--193, 1997.

\bibitem{kachman1995}
Kachman, S.~D. and Smith, J.~A.
Alternative measures of accuracy in plant spacing for planters using single seed metering.
\textit{Transactions of the ASAE}, 38(2):379--387, 1995.

\bibitem{karayel2004}
Karayel, D., Barut, Z.~B., and {\"O}zmerzi, A.
Mathematical modelling of vacuum pressure on a precision seeder.
\textit{Biosystems Engineering}, 87(4):437--444, 2004.

\bibitem{kayacan2015}
Kayacan, E., Kayacan, E., Ramon, H., and Saeys, W.
Towards agrobots: identification of the yaw dynamics and trajectory tracking of an autonomous tractor.
\textit{Computers and Electronics in Agriculture}, 115:78--87, 2015.

\bibitem{khalil2002}
Khalil, H.~K.
\textit{Nonlinear Systems}.
Prentice Hall, 3rd edition, 2002.

\bibitem{kocher1998}
Kocher, M.~F., Lan, Y., Chen, C., and Smith, J.~A.
Opto-electronic sensor system for rapid evaluation of planter seed spacing uniformity.
\textit{Transactions of the ASAE}, 41(1):237--245, 1998.

\bibitem{kraus2013}
Kraus, T., Ferreau, H.~J., Kayacan, E., Ramon, H., De~Baerdemaeker, J., Diehl, M., and Saeys, W.
Moving horizon estimation and nonlinear model predictive control for autonomous agricultural vehicles.
\textit{Computers and Electronics in Agriculture}, 98:25--33, 2013.

\bibitem{lavalle2006}
LaValle, S.~M.
\textit{Planning Algorithms}.
Cambridge University Press, 2006.

\bibitem{lei2021}
Lei, X., Hu, H., Wu, W., Liu, H., Liu, L., Yang, W., and co-authors.
Seed motion characteristics and seeding performance of a centralised seed metering system for rapeseed investigated by DEM simulation and bench testing.
\textit{Biosystems Engineering}, 203:22--33, 2021.

\bibitem{lowenberg2019}
Lowenberg-DeBoer, J. and Erickson, B.
Setting the record straight on precision agriculture adoption.
\textit{Agronomy Journal}, 111(4):1552--1569, 2019.

\bibitem{maleki2006}
Maleki, M.~R., Jafari, J.~F., Raufat, M.~H., Mouazen, A.~M., and De~Baerdemaeker, J.
Evaluation of seed distribution uniformity of a multi-flight auger as a grain drill metering device.
\textit{Biosystems Engineering}, 94(4):535--543, 2006.

\bibitem{mangus2017}
Mangus, D.~L., Sharda, A., Engelhardt, A., Flippo, D., Strasser, R., Luck, J.~D., and Griffin, T.
Development of high-speed camera hardware and software package to evaluate real-time electric seed meter accuracy of a variable rate planter.
\textit{Computers and Electronics in Agriculture}, 142:314--325, 2017.

\bibitem{masarei2019}
Masarei, M., Guzzomi, A.~L., Merritt, D.~J., and Erickson, T.~E.
Improvements to mechanical direct seeder design guided by the optimal sowing depth of soft spinifex (\textit{Triodia pungens}).
In \textit{Mine Closure 2019}, pages 207--222. Australian Centre for Geomechanics, 2019.

\bibitem{mayne2000}
Mayne, D.~Q., Rawlings, J.~B., Rao, C.~V., and Scokaert, P.~O.~M.
Constrained model predictive control: stability and optimality.
\textit{Automatica}, 36(6):789--814, 2000.

\bibitem{navid2011}
Navid, H., Ebrahimian, S., Gassemzadeh, H.~R., and Mousavinia, M.~J.
Laboratory evaluation of seed metering device using image processing method.
\textit{Australian Journal of Agricultural Engineering}, 2(1):1--4, 2011.

\bibitem{norremark2008}
N{\o}rremark, M., Griepentrog, H.~W., Nielsen, J., and S{\o}gaard, H.~T.
The development and assessment of the accuracy of an autonomous GPS-based system for intra-row mechanical weed control in row crops.
\textit{Biosystems Engineering}, 101(4):396--410, 2008.

\bibitem{oksanen2009}
Oksanen, T. and Visala, A.
Coverage path planning algorithms for agricultural field machines.
\textit{Journal of Field Robotics}, 26(8):651--668, 2009.

\bibitem{panning2000}
Panning, J.~W., Kocher, M.~F., Smith, J.~A., and Kachman, S.~D.
Laboratory and field testing of seed spacing uniformity for sugarbeet planters.
\textit{Applied Engineering in Agriculture}, 16(1):7--13, 2000.

\bibitem{pedersen2006}
Pedersen, S.~M., Fountas, S., Have, H., and Blackmore, B.~S.
Agricultural robots --- system analysis and economic feasibility.
\textit{Precision Agriculture}, 7:295--308, 2006.

\bibitem{plessen2017}
Graf~Plessen, M.~M. and Bemporad, A.
Reference trajectory planning under constraints and path tracking using linear time-varying model predictive control for agricultural machines.
\textit{Biosystems Engineering}, 153:28--41, 2017.

\bibitem{rajamani2012}
Rajamani, R.
\textit{Vehicle Dynamics and Control}.
Springer, 2nd edition, 2012.

\bibitem{rawlings2017}
Rawlings, J.~B., Mayne, D.~Q., and Diehl, M.
\textit{Model Predictive Control: Theory, Computation, and Design}.
Nob Hill Publishing, 2nd edition, 2017.

\bibitem{sarkka2013}
S{\"a}rkk{\"a}, S.
\textit{Bayesian Filtering and Smoothing}.
Cambridge University Press, 2013.

\bibitem{shamshiri2018}
Shamshiri, R.~R., Weltzien, C., Hameed, I.~A., Yule, I.~J., Grift, T.~E., Balasundram, S.~K., Pitonakova, L., Ahmad, D., and Chowdhary, G.
Research and development in agricultural robotics: a perspective of digital farming.
\textit{International Journal of Agricultural and Biological Engineering}, 11(4):1--14, 2018.

\bibitem{singh2005}
Singh, R.~C., Singh, G., and Saraswat, D.~C.
Optimisation of design and operational parameters of a pneumatic seed metering device for planting cottonseeds.
\textit{Biosystems Engineering}, 92(4):429--438, 2005.

\bibitem{sorensen2010}
S{\o}rensen, C.~G. and Bochtis, D.~D.
Conceptual model of fleet management in agriculture.
\textit{Biosystems Engineering}, 105(1):41--50, 2010.

\bibitem{spekken2013}
Spekken, M. and de~Bruin, S.
Optimized routing on agricultural fields by minimizing maneuvering and servicing time.
\textit{Precision Agriculture}, 14:224--244, 2013.

\bibitem{stjack2013}
St~Jack, D., Hesterman, D.~C., and Guzzomi, A.~L.
Precision metering of \textit{Santalum spicatum} (Australian sandalwood) seeds.
\textit{Biosystems Engineering}, 115(2):171--183, 2013.

\bibitem{thrun2005}
Thrun, S., Burgard, W., and Fox, D.
\textit{Probabilistic Robotics}.
MIT Press, 2005.

\bibitem{usdaplanter}
USDA Forest Service, Missoula Technology and Development Center.
Evaluation of single-seed planters.
Tech Tip 1024-2306P-MTDC.

\bibitem{vahdanjoo2023}
Vahdanjoo, M., Gislum, R., and S{\o}rensen, C.~A.~G.
Three-dimensional area coverage planning tool for robotic application. Case: seeding operation.
SSRN abstract 4594418, 2023.

\bibitem{wang2024yolov10}
Wang, A., Chen, H., Liu, L., Chen, K., Lin, Z., Han, J., and Ding, G.
YOLOv10: real-time end-to-end object detection.
arXiv:2405.14458, 2024.

\bibitem{welch2006}
Welch, G. and Bishop, G.
An introduction to the Kalman filter.
Technical Report TR 95-041, University of North Carolina at Chapel Hill, 2006.

\bibitem{woodhandbook}
Forest Products Laboratory.
\textit{Wood Handbook: Wood as an Engineering Material}.
General Technical Report FPL-GTR-190, USDA Forest Service, Madison, WI, 2010.

\bibitem{xu2022}
Xu, J., Lai, J., Guo, R., and Lu, X.
Efficiency-oriented MPC algorithm for path tracking in autonomous agricultural machinery.
\textit{Agronomy}, 12(7):1662, 2022.

\bibitem{zhang2025path}
Zhang, J., Qian, Y., Li, P., Bai, Y., and Wang, S.
A path planning and optimization framework for autonomous seeding robots in intelligent agriculture.
\textit{Journal of Physics: Conference Series}, 3277:012196, 2025.

\bibitem{zhaoJ2023}
Zhao, J., Xi, X., Shi, Y., Zhang, B., Qu, J., and Zhang, Y.
An online method for detecting seeding performance based on an improved YOLOv5s model.
\textit{Agronomy}, 13(9):2391, 2023.

\bibitem{zhao2024hs}
Zhao, P., Gao, X., Su, Y., Xu, Y., and Huang, Y.
Investigation of seeding performance of a novel high-speed precision seed metering device based on numerical simulation and high-speed camera.
\textit{Computers and Electronics in Agriculture}, 217:108563, 2024.

\bibitem{zhao2025yolo}
Zhao, P., Wu, X., Cheng, H., Gao, X., Liu, Z., Lai, Q., and Li, H.
Performance evaluation of a high-speed maize seed-metering device using an improved YOLOv5s object detection and tracking algorithm.
\textit{Smart Agricultural Technology}, 11:100997, 2025.

\bibitem{bakker2011}
Bakker, T., van Asselt, K., Bontsema, J., M{\"u}ller, J., and van Straten, G.
Autonomous navigation using a robot platform in a sugar beet field.
\textit{Biosystems Engineering}, 109(4):357--368, 2011.

\end{thebibliography}


\end{document}
